% Généralités sur l'oxydoréduction en solution aqueuse — Chimie 1ère E — Cours signé OnBuch+
% Programme officiel MINESEC, Première C, D, E, TI. Compiler avec : tectonic N05-generalites-oxydoreduction-solution-aqueuse.tex
\newcommand{\DOCMATIERE}{Chimie}
\newcommand{\DOCNIVEAU}{1ère E}
\newcommand{\DOCMODULE}{Module 2 : Oxydoréduction}
\newcommand{\DOCLECON}{05}
\newcommand{\DOCTITRE}{Généralités sur l'oxydoréduction en solution aqueuse}
% OnBuch+ — préambule « cours signés » ENRICHI (compilé avec tectonic / XeLaTeX)
% Système visuel « Pop » d'OnBuch+ : fond crème (jamais blanc), bordures encre
% épaisses, ombres « dures » décalées, titres Archivo Black en capitales.
% Version enrichie pour les cours de chimie : chimie (mhchem, chemfig),
% graphes (pgfplots), boîtes pédagogiques supplémentaires (définition,
% propriété, méthode, attention, le savais-tu, exercice résolu, exercice,
% TP, bilan), page de garde refaite et signature OnBuch+ ancrée en bas.
%
% Macros attendues dans main.tex AVANT \input{preamble} :
%   \DOCMATIERE  \DOCNIVEAU  \DOCMODULE  \DOCLECON (numéro)  \DOCTITRE
\documentclass[11pt]{article}

\usepackage{fontspec}
\usepackage[french]{babel}
\usepackage[a4paper,margin=2cm,top=3.4cm,bottom=2.8cm,headheight=1.4cm,footskip=1.3cm]{geometry}
\usepackage{amsmath,amssymb,amsfonts}
\usepackage[table]{xcolor} % [table] : fournit aussi \rowcolors (alternance de lignes)
\usepackage{pagecolor}
\usepackage{tikz}
\usetikzlibrary{arrows.meta,positioning,calc,shapes.geometric,shapes.misc,shapes.arrows,decorations.pathmorphing,decorations.markings,patterns,shadows,mindmap,trees,fit,backgrounds,matrix,intersections}
\usepackage{pgfplots}
\pgfplotsset{compat=1.18}
\usepackage[version=4]{mhchem}
\usepackage{chemfig}
\usepackage{enumitem}
\usepackage{fancyhdr}
% [most] charge toutes les bibliothèques tcolorbox (listings, minted, raster,
% documentation...) et gonfle inutilement la mémoire TeX sur un document long
% (~300 boîtes) : on ne charge que ce qui est réellement utilisé.
\usepackage[skins,breakable]{tcolorbox}
\usepackage{graphicx}
\usepackage{colortbl}
\usepackage{booktabs}
\usepackage{tabularx}
\usepackage{array}
\usepackage{multicol}
\usepackage{siunitx}
% parse-numbers=false : accepte aussi \SI{2,5 \times 10^{-3}}{...}
\sisetup{locale=FR,output-decimal-marker={,},group-minimum-digits=4,parse-numbers=false}
\DeclareSIUnit\ohm{\ensuremath{\symup{\Omega}}} % U+2126 absent de la police
\usepackage{qrcode}
% \degree : commande fréquemment utilisée par les modèles pour un angle ou une
% température en dehors de \SI{}{} (non fournie par les paquets chargés).
\providecommand{\degree}{^{\circ}}
\usepackage{lastpage}
\usepackage{hyperref}
\hypersetup{hidelinks}

% ============================================================
% Palette « Pop » OnBuch+ — mêmes hex que la classe P (plus_ui.dart)
% ============================================================
\definecolor{poporange}{HTML}{F59321}
\definecolor{popdark}{HTML}{E07A0C}
\definecolor{popcream}{HTML}{FFF6E8}
\definecolor{popink}{HTML}{1C1714}
\definecolor{poppurple}{HTML}{7A5AE0}
\definecolor{popgreen}{HTML}{2E9E6A}
\definecolor{poppink}{HTML}{E0457B}
\definecolor{popblue}{HTML}{2F7DE1}
\definecolor{popgold}{HTML}{C98A12}
\definecolor{popmuted}{HTML}{8A7F76}
\definecolor{popline}{HTML}{EADFD2}
\definecolor{popgray}{HTML}{8A7F76}
\colorlet{popgrey}{popgray}
% Alias tolérants (le modèle nomme parfois les couleurs différemment)
\colorlet{popwhite}{white}
\colorlet{popblack}{popink}
\colorlet{popred}{poppink}
\colorlet{popyellow}{popgold}
% Teintes claires pour fonds de tableaux / zones
\colorlet{popblueL}{popblue!10!white}
\colorlet{poporangeL}{poporange!14!white}
\colorlet{popgreenL}{popgreen!12!white}
\colorlet{poppurpleL}{poppurple!12!white}
\colorlet{poppinkL}{poppink!12!white}
\colorlet{popgoldL}{popgold!14!white}

\pagecolor{popcream}

% ============================================================
% Typographie
% ============================================================
\defaultfontfeatures{Ligatures=TeX}
\setmainfont{PlusJakartaSans-Medium.ttf}[Path=../../../fonts/,BoldFont=PlusJakartaSans-Bold.ttf]
% Maths en sans-serif (Fira Math) avec chiffres et lettres droites de Plus
% Jakarta, pour que les formules se fondent dans le texte.
\usepackage[math-style=ISO,bold-style=ISO]{unicode-math}
\setmathfont{FiraMath-Regular.otf}[Path=../../../fonts/]
\setmathfont{PlusJakartaSans-Medium.ttf}[Path=../../../fonts/,range={up/num,up/{Latin,latin},\mathup}]
\usepackage{icomma} % virgule décimale sans espace en mode math
% Caractères absents de la police de texte (sinon carré vide) : rendus par
% leur équivalent mathématique, en texte comme en formule.
\usepackage{newunicodechar}
\newunicodechar{α}{\ensuremath{\alpha}}
\newunicodechar{Α}{\ensuremath{\alpha}}
\newunicodechar{β}{\ensuremath{\beta}}
\newunicodechar{δ}{\ensuremath{\delta}}
\newunicodechar{✓}{\popcheck{popgreen}}
\newfontfamily\popdisplay{ArchivoBlack-Regular.ttf}[Path=../../../fonts/]
\newfontfamily\poplogo{SpaceGrotesk-Bold.ttf}[Path=../../../fonts/]
\newfontfamily\popsemi{PlusJakartaSans-SemiBold.ttf}[Path=../../../fonts/]
\newfontfamily\popxbold{PlusJakartaSans-ExtraBold.ttf}[Path=../../../fonts/]
\newfontfamily\popxboldls{PlusJakartaSans-ExtraBold.ttf}[Path=../../../fonts/,LetterSpace=3.0]

\setlist[enumerate]{leftmargin=1.7em,itemsep=5pt,topsep=4pt}
\setlist[itemize]{leftmargin=1.7em,itemsep=4pt,topsep=4pt,label={\color{poporange}\textbullet}}
\setlength{\parindent}{0pt}
\setlength{\parskip}{5pt}
\renewcommand{\arraystretch}{1.25}

% chemfig : liaisons un peu plus épaisses, encre Pop
\setchemfig{atom sep=2em,bond style={line width=0.9pt,popink},cram width=3pt}

% pgfplots : style Pop par défaut
\pgfplotsset{
  every axis/.append style={axis line style={popink,line width=1pt},tick style={popink},
    grid style={popline,line width=0.5pt},label style={font=\small},tick label style={font=\footnotesize}},
  popcurve/.style={poporange,line width=1.8pt,no markers,smooth},
}

% ============================================================
% Pill / squiggle
% ============================================================
\newcommand{\poppill}[3]{%
  \tikz[baseline=-0.55ex]\node[draw=popink,line width=1.2pt,rounded corners=2.6pt,%
    fill=#2,inner xsep=6.5pt,inner ysep=3pt]{%
    \popxboldls\fontsize{7}{7}\selectfont\color{#3}\MakeUppercase{#1}};%
}
\newcommand{\popsquiggle}[1]{%
  \tikz[baseline=-0.4ex]{
    \draw[#1,line width=2.6pt,line cap=round]
      plot [smooth] coordinates {(0,0.09) (0.34,0.01) (0.68,0.21) (1.02,0.01) (1.36,0.10)};
  }%
}
% Mot-clé surligné (marqueur orange sous le texte)
\newcommand{\cle}[1]{{\popxbold\color{popdark}#1}}

% ============================================================
% En-tête / pied
% ============================================================
\pagestyle{fancy}
\fancyhf{}
\renewcommand{\headrulewidth}{0pt}
\renewcommand{\footrulewidth}{0pt}
\lhead{\raisebox{-0.15ex}{\poplogo\fontsize{13}{13}\selectfont\color{popink}OnBuch\color{poporange}+}}
\rhead{\poppill{\DOCMATIERE\ \textbullet\ \DOCNIVEAU\ \textbullet\ Leçon \DOCLECON}{white}{popink}}
\lfoot{\popxbold\fontsize{7.6}{7.6}\selectfont\color{popmuted}\MakeUppercase{Cours signé OnBuch+}\ \textbullet\ {\popsemi\DOCTITRE}}
\rfoot{\tikz[baseline=-0.6ex]\node[rounded corners=8.5pt,fill=popink,minimum height=17pt,minimum width=17pt,inner xsep=5pt,inner sep=0]{\popxbold\fontsize{8}{8}\selectfont\color{popcream}\thepage\,/\,\pageref*{LastPage}};}

% ============================================================
% Titres
% ============================================================
\newcommand{\chaptitre}[2]{%
  \begin{tcolorbox}[enhanced,colback=poporange,colframe=popink,boxrule=2.6pt,arc=11pt,
      drop shadow=popink,left=15pt,right=15pt,top=12pt,bottom=11pt,before skip=3pt,after skip=18pt]
    \poppill{#2}{popink}{popcream}\par\vspace{9pt}
    {\raggedright\hyphenpenalty=10000\exhyphenpenalty=10000\popdisplay\fontsize{22}{24}\selectfont\color{popcream}\MakeUppercase{#1}\par}
  \end{tcolorbox}
}
% Section numérotée : pastille orange avec le numéro + titre Archivo
\newcounter{popsec}
\newcommand{\coursec}[1]{%
  \stepcounter{popsec}\setcounter{popsub}{0}%
  \par\vspace{16pt}\noindent
  \begin{minipage}{\linewidth}
  \tikz[baseline=-0.7ex]\node[draw=popink,line width=1.6pt,fill=poporange,rounded corners=4pt,minimum size=22pt,inner sep=0]{\popdisplay\fontsize{12}{12}\selectfont\color{popcream}\thepopsec};%
  \hspace{8pt}{\popdisplay\fontsize{15}{17}\selectfont\color{popink}\MakeUppercase{#1}}\par
  \vspace{4pt}\noindent{\color{poporange}\rule{36pt}{3.6pt}}
  \end{minipage}\par\nopagebreak\vspace{8pt}
}
\newcounter{popsub}
\newcommand{\courssub}[1]{%
  \stepcounter{popsub}%
  \par\vspace{9pt}\noindent{\popxbold\fontsize{11.5}{13}\selectfont\color{popink}{\color{poporange}\thepopsec.\thepopsub}\hspace{6pt}#1}\par\nopagebreak\vspace{4pt}
}

% ============================================================
% Boîtes pédagogiques — même recette : fond blanc, bordure épaisse colorée,
% ombre dure encre, badge plein. Titre optionnel [#1].
% ============================================================
\tcbset{popbase/.style={breakable,enhanced,colback=white,boxrule=2.3pt,arc=9pt,
  shadow={3pt}{-3pt}{0pt}{popink},left=14pt,right=14pt,top=11pt,bottom=11pt,before skip=11pt,after skip=15pt,
  fontupper=\popsemi\fontsize{10.6}{14.5}\selectfont\color{popink},
  fontlower=\popsemi\fontsize{10.6}{14.5}\selectfont\color{popink}}}
% Coche dessinée (le glyphe ✓ n'existe pas dans les polices du document)
\newcommand{\popcheck}[1]{\tikz[baseline=-0.5ex]\draw[#1,line width=1.6pt,line cap=round,line join=round](-0.13,0.0)--(-0.04,-0.09)--(0.13,0.1);}
\newcommand{\popboxhead}[3]{\poppill{#1}{#2}{white}\ifx\relax#3\relax\else\hspace{6pt}{\popxbold\fontsize{10.6}{12}\selectfont\color{popink}#3}\fi\par\vspace{7pt}}

\newtcolorbox{aretenir}[1][]{popbase,colframe=popgreen,before upper={\popboxhead{À retenir}{popgreen}{#1}}}
\newtcolorbox{exemplebox}[1][]{popbase,colframe=poporange,before upper={\popboxhead{Exemple}{poporange}{#1}}}
\newtcolorbox{definition}[1][]{popbase,colframe=popblue,before upper={\popboxhead{Définition}{popblue}{#1}}}
\newtcolorbox{propriete}[1][]{popbase,colframe=poppurple,before upper={\popboxhead{Propriété}{poppurple}{#1}}}
\newtcolorbox{methode}[1][]{popbase,colframe=popgold,before upper={\popboxhead{Méthode}{popgold}{#1}}}
\newtcolorbox{attention}[1][]{popbase,colframe=poppink,before upper={\popboxhead{Attention}{poppink}{#1}}}
\newtcolorbox{experience}[1][]{popbase,colframe=popblue,colback=popblueL,before upper={\popboxhead{Expérience}{popblue}{#1}}}
% « Le savais-tu ? » : carte violette pleine (contexte, culture, Cameroun)
\newtcolorbox{savaistu}[1][]{popbase,colback=poppurple,colframe=popink,
  fontupper=\popsemi\fontsize{10.4}{14.2}\selectfont\color{white},
  before upper={\poppill{Le savais-tu\,?}{white}{poppurple}\ifx\relax#1\relax\else\hspace{6pt}{\popxbold\color{white}#1}\fi\par\vspace{7pt}}}
% Exercice résolu : énoncé en haut, \tcblower puis la solution sur fond crème
\newcounter{popexo}\newcounter{popexr}
\newtcolorbox{exoresolu}[1][]{popbase,colframe=poporange,colbacklower=popcream,
  segmentation style={popink,line width=1.2pt,dashed},
  before upper={\stepcounter{popexr}\popboxhead{Exercice résolu \thepopexr}{poporange}{#1}},
  before lower={\poppill{Solution}{popink}{popcream}\par\vspace{6pt}}}
% Exercice (non résolu en place) — la correction est dans la partie Corrigés
\newtcolorbox{exercice}[1][]{popbase,colframe=popink,
  before upper={\stepcounter{popexo}\popboxhead{Exercice \thepopexo}{popink}{#1}}}
% Corrigé
\newtcolorbox{corrige}[1][]{popbase,colframe=popgreen,colback=popgreenL,
  before upper={\popboxhead{Corrigé}{popgreen}{#1}}}
% Objectifs / prérequis / bilan
\newtcolorbox{objectifs}{popbase,colframe=popink,colback=poporangeL,
  before upper={\popboxhead{Objectifs de la leçon}{popink}{}}}
\newtcolorbox{prerequis}{popbase,colframe=popink,colback=popblueL,
  before upper={\popboxhead{Prérequis}{popblue}{}}}
\newtcolorbox{bilan}{popbase,colframe=popink,colback=white,boxrule=2.8pt,
  before upper={\popboxhead{Fiche bilan}{poporange}{L'essentiel à connaître pour l'examen}}}
% Figure encadrée avec légende
\newtcolorbox{popfigure}[1][]{enhanced,breakable=false,colback=white,colframe=popline,boxrule=1.4pt,arc=8pt,
  left=10pt,right=10pt,top=10pt,bottom=8pt,before skip=10pt,after skip=12pt,halign=center,
  lower separated=false,halign lower=center,
  fontlower=\popsemi\fontsize{9}{11.5}\selectfont\color{popmuted},
  #1}
\newcommand{\legende}[1]{\par\vspace{6pt}{\popsemi\fontsize{9}{11.5}\selectfont\color{popmuted}#1}}

% ============================================================
% Signature OnBuch+
% ============================================================
\newcommand{\poplogoinline}[1]{{\poplogo\fontsize{#1}{#1}\selectfont\color{popink}OnBuch\color{poporange}+}}
\newcommand{\signatureonbuch}{%
  \begin{tcolorbox}[enhanced,colback=popink,colframe=popink,boxrule=2.2pt,arc=10pt,
      drop shadow=poporange,left=14pt,right=14pt,top=10pt,bottom=10pt,before skip=0pt,after skip=0pt]
    \tikz[baseline=-0.7ex]\node[circle,fill=popgreen,draw=popcream,line width=1.3pt,minimum size=22pt,inner sep=0]{\popcheck{white}};%
    \hspace{9pt}\begin{minipage}[c]{0.62\linewidth}
      {\popxbold\fontsize{12}{13}\selectfont\color{popcream}Signé\ \textbullet\ {\poplogo OnBuch\color{poporange}+}}\par\vspace{2pt}
      {\raggedright\popsemi\fontsize{8}{10.5}\selectfont\color{popline}\MakeUppercase{Équipe pédagogique OnBuch+}\\\MakeUppercase{Conforme au programme officiel MINESEC}\par}
    \end{minipage}\hfill
    {\poplogo\fontsize{22}{22}\selectfont\color{popcream}OnBuch\color{poporange}+}
  \end{tcolorbox}%
}

% ============================================================
% Page de garde
% #1 titre · #2 matière · #3 niveau · #4 module · #5 n° leçon · #6 durée ·
% #7 description / objectifs (texte ou itemize)
% ============================================================
\newcommand{\pagegarde}[7]{%
  \thispagestyle{empty}
  \newgeometry{a4paper,margin=1.7cm,top=1.6cm,bottom=1.4cm}%
  \begin{tikzpicture}[remember picture,overlay]
    \fill[poporange!18] ([xshift=-3.2cm,yshift=-3.6cm]current page.north east) circle (4.6cm);
    \fill[poppurple!14] ([xshift=2.4cm,yshift=7.6cm]current page.south west) circle (3.8cm);
  \end{tikzpicture}%
  {\poplogo\fontsize{30}{30}\selectfont\color{popink}OnBuch\color{poporange}+}\hfill
  \raisebox{0.6ex}{\poppill{Cours signé \textbullet\ Premium}{popink}{popcream}}\par
  \vspace{1pt}\popsquiggle{poppurple}\par
  \vspace{8pt}
  \begin{tcolorbox}[enhanced,colback=poporange,colframe=popink,boxrule=2.8pt,arc=13pt,
      drop shadow=popink,left=18pt,right=18pt,top=12pt,bottom=13pt,after skip=10pt]
    \poppill{#2 \textbullet\ #3}{popink}{popcream}\hspace{5pt}\poppill{#4}{white}{popink}\par\vspace{8pt}
    {\popdisplay\fontsize{14}{14}\selectfont\color{popink}LEÇON #5}\par\vspace{4pt}
    {\raggedright\hyphenpenalty=10000\exhyphenpenalty=10000\popdisplay\fontsize{25}{27}\selectfont\color{popcream}\MakeUppercase{#1}\par}
    \vspace{8pt}
    {\popxbold\fontsize{9.5}{9.5}\selectfont\color{popink}\MakeUppercase{Durée conseillée : #6}}
  \end{tcolorbox}
  \begin{tcolorbox}[enhanced,colback=white,colframe=popink,boxrule=2.2pt,arc=10pt,
      drop shadow=popink,left=16pt,right=16pt,top=10pt,bottom=10pt,after skip=8pt]
    \poppill{Dans cette leçon}{poppurple}{white}\par\vspace{5pt}
    \popsemi\fontsize{9.3}{12.6}\selectfont\color{popink}#7
  \end{tcolorbox}
  \noindent\begin{minipage}[t]{0.487\linewidth}
    \begin{tcolorbox}[enhanced,colback=popgreen,colframe=popink,boxrule=2.2pt,arc=10pt,
        drop shadow=popink,left=11pt,right=11pt,top=10pt,bottom=10pt,height=3.1cm,valign=center]
      \tcbox[colback=white,colframe=popink,boxrule=1.3pt,arc=5pt,boxsep=0pt,nobeforeafter,
        left=3pt,right=3pt,top=3pt,bottom=3pt,valign=center]{\qrcode[height=1.9cm]{https://play.google.com/store/apps/details?id=com.luvvix.onbuch&hl=fr}}%
      \hspace{8pt}\begin{minipage}[c]{0.5\linewidth}
        {\popdisplay\fontsize{11}{12.5}\selectfont\color{white}\MakeUppercase{Télécharge OnBuch}}\par\vspace{4pt}
        {\raggedright\popsemi\fontsize{8.4}{11}\selectfont\color{white}Gratuit \textbullet\ Google Play\\Tuteur IA, annales, concours, résultats\par}
      \end{minipage}
    \end{tcolorbox}
  \end{minipage}\hfill
  \begin{minipage}[t]{0.487\linewidth}
    \begin{tcolorbox}[enhanced,colback=poppurple,colframe=popink,boxrule=2.2pt,arc=10pt,
        drop shadow=popink,left=11pt,right=11pt,top=10pt,bottom=10pt,height=3.1cm,valign=center]
      \tikz[baseline=-0.7ex]\node[circle,fill=white,draw=popink,line width=1.3pt,minimum size=1.6cm,inner sep=0]{\popdisplay\fontsize{24}{24}\selectfont\color{poppurple}+};%
      \hspace{8pt}\begin{minipage}[c]{0.6\linewidth}
        {\popdisplay\fontsize{11}{12.5}\selectfont\color{white}\MakeUppercase{Passe à OnBuch+}}\par\vspace{4pt}
        {\raggedright\popsemi\fontsize{8.4}{11}\selectfont\color{white}Tous les cours signés, Tuteur IA illimité, fiches PDF — onglet OnBuch+ de l'app\par}
      \end{minipage}
    \end{tcolorbox}
  \end{minipage}
  \vspace{8pt}
  \popcontenu
  \vfill
  \signatureonbuch
  \restoregeometry
  \newpage
}

% Rangée « Ce cours contient » (page de garde)
\newcommand{\poptile}[3]{% #1 couleur #2 gros texte #3 légende
  \begin{tcolorbox}[enhanced,colback=#1,colframe=popink,boxrule=2pt,arc=9pt,shadow={2.5pt}{-2.5pt}{0pt}{popink},
      left=6pt,right=6pt,top=9pt,bottom=9pt,halign=center,valign=center,height=1.75cm,width=\linewidth,nobeforeafter]
    {\popdisplay\fontsize{12.5}{13}\selectfont\color{white}\MakeUppercase{#2}}\par\vspace{3pt}
    {\centering\popsemi\fontsize{7.8}{9.5}\selectfont\color{white}#3\par}
  \end{tcolorbox}}
\newcommand{\popcontenu}{%
  \par\noindent{\popxboldls\fontsize{7.5}{7.5}\selectfont\color{popmuted}CE COURS SIGNÉ CONTIENT}\par\vspace{7pt}
  \noindent\begin{minipage}[t]{0.235\linewidth}\poptile{popblue}{Cours}{complet, illustré, pas à pas}\end{minipage}\hfill
  \begin{minipage}[t]{0.235\linewidth}\poptile{poporange}{Méthodes}{et exercices résolus}\end{minipage}\hfill
  \begin{minipage}[t]{0.235\linewidth}\poptile{poppink}{Exercices}{type examen + corrigés}\end{minipage}\hfill
  \begin{minipage}[t]{0.235\linewidth}\poptile{popgreen}{Fiche bilan}{l'essentiel à retenir}\end{minipage}\par}

% Sommaire
\newtcolorbox{sommaire}{enhanced,colback=white,colframe=popink,boxrule=2.2pt,arc=10pt,
  drop shadow=popink,left=18pt,right=18pt,top=13pt,bottom=13pt,before skip=6pt,after skip=20pt,
  before upper={\poppill{Sommaire}{poppurple}{white}\par\vspace{9pt}\popsemi\fontsize{10.6}{14.5}\selectfont\color{popink}}}

% Signature de fin de cours — poussée tout en bas de la dernière page
\newcommand{\findecours}{%
  \par\vfill\nopagebreak
  \begin{center}\popsquiggle{poporange}\end{center}\vspace{4pt}
  \signatureonbuch
}
% Glyphes absents de Fira Math : redéfinis à partir de glyphes disponibles
\AtBeginDocument{%
  \def\setminus{\mathbin{\backslash}}\let\smallsetminus\setminus
  \def\vdots{\mathord{\vbox{\baselineskip3.2pt\lineskiplimit0pt\kern4pt\hbox{.}\hbox{.}\hbox{.}}}}%
  \def\ddots{\mathinner{\mkern1mu\raise6.5pt\hbox{.}\mkern2mu\raise3.6pt\hbox{.}\mkern2mu\raise0.7pt\hbox{.}\mkern1mu}}%
  \def\triangle{\mathord{\scalebox{0.9}{$\symup{\Delta}$}}}%
  \def\nabla{\mathord{\raisebox{\depth}{\scalebox{1}[-1]{$\symup{\Delta}$}}}}%
  \def\checkmark{\popcheck{popdark}}%
  \def\star{\mathbin{\ast}}\def\bigstar{\mathord{\ast}}%
  \def\diamond{\mathbin{\scalebox{0.6}{$\Diamond$}}}%
  \def\bigcup{\mathop{\mathchoice{\raisebox{-0.3em}{\scalebox{1.7}{$\cup$}}}{\scalebox{1.25}{$\cup$}}{\cup}{\cup}}\displaylimits}%
  \def\bigcap{\mathop{\mathchoice{\raisebox{-0.3em}{\scalebox{1.7}{$\cap$}}}{\scalebox{1.25}{$\cap$}}{\cap}{\cap}}\displaylimits}%
  \def\lesseqqgtr{\mathrel{\lessgtr}}\def\bigoplus{\mathop{\oplus}\displaylimits}%
}
\providecommand{\ssi}{si et seulement si} % abréviation fréquente

%% FIGFIX-BEGIN v3 — correctifs globaux des figures (docs/onbuch-generation/tools/figfix)
\usepackage{adjustbox}\usepackage{etoolbox}
% toute figure TikZ/pgfplots plus large que la ligne est réduite à la largeur disponible
\BeforeBeginEnvironment{tikzpicture}{\begin{adjustbox}{max width=\linewidth}}
\AfterEndEnvironment{tikzpicture}{\end{adjustbox}}
% légendes pgfplots lisibles par-dessus les courbes
\pgfplotsset{every axis/.append style={legend style={fill=white,fill opacity=0.92,text opacity=1,draw=black!15,inner sep=3pt}}}
% étiquettes d'arêtes (syntaxe edge["..."]) avec fond blanc
\tikzset{every edge quotes/.append style={fill=white,inner sep=1.2pt}}
% bulles de cartes mentales : texte à la ligne dans la bulle, bulle un peu plus grande
\tikzset{every concept/.append style={text width=2.0cm,align=center,minimum size=2.3cm,inner sep=2pt}}
%% FIGFIX-END
\begin{document}

\pagegarde{Généralités sur l'oxydoréduction en solution aqueuse}{Chimie}{1ère E}{Module 2 : Oxydoréduction}{05}{8 h}{Cette leçon introduit l'oxydoréduction en solution aqueuse à travers l'action des acides dilués sur les métaux et la réaction entre un ion métallique et un métal. Elle définit les notions d'oxydation, de réduction, d'oxydant et de réducteur, et apprend à écrire demi-équations électroniques et équations-bilans.
\vspace{4pt}
\begin{multicols}{2}\small
\begin{itemize}
\item À la fin de cette leçon, je sais décrire l'action des acides chlorhydrique et sulfurique dilués sur les métaux usuels (zinc, fer, aluminium, cuivre, argent, or)
\item À la fin de cette leçon, je sais identifier le gaz dégagé lors de l'action d'un acide sur un métal et le caractériser
\item À la fin de cette leçon, je sais définir l'oxydation, la réduction, l'oxydant, le réducteur et une réaction d'oxydoréduction
\item À la fin de cette leçon, je sais écrire les demi-équations électroniques d'un couple oxydant/réducteur
\item À la fin de cette leçon, je sais établir l'équation-bilan d'une réaction d'oxydoréduction à partir des demi-équations
\item À la fin de cette leçon, je sais décrire et interpréter l'action du zinc ou du fer sur des ions métalliques (Cu2+, Zn2+, Fe2+)
\end{itemize}\end{multicols}}

\begin{sommaire}
\begin{multicols}{2}
\begin{enumerate}
\item Action d'un acide dilué sur un métal
\item Réaction entre un ion métallique et un métal
\item Oxydation, réduction, oxydant, réducteur
\item Demi-équations électroniques et équation-bilan
\item Prévision des réactions et applications
\item Travaux pratiques
\item Méthodes et exercices résolus
\item Exercices
\item Corrigés des exercices
\item Fiche bilan
\end{enumerate}
\end{multicols}
\end{sommaire}

\begin{objectifs}
\begin{itemize}
\item À la fin de cette leçon, je sais décrire l'action des acides chlorhydrique et sulfurique dilués sur les métaux usuels (zinc, fer, aluminium, cuivre, argent, or)
\item À la fin de cette leçon, je sais identifier le gaz dégagé lors de l'action d'un acide sur un métal et le caractériser
\item À la fin de cette leçon, je sais définir l'oxydation, la réduction, l'oxydant, le réducteur et une réaction d'oxydoréduction
\item À la fin de cette leçon, je sais écrire les demi-équations électroniques d'un couple oxydant/réducteur
\item À la fin de cette leçon, je sais établir l'équation-bilan d'une réaction d'oxydoréduction à partir des demi-équations
\item À la fin de cette leçon, je sais décrire et interpréter l'action du zinc ou du fer sur des ions métalliques (Cu2+, Zn2+, Fe2+)
\item À la fin de cette leçon, je sais prévoir si une réaction a lieu entre un métal et un ion métallique ou un acide
\item À la fin de cette leçon, je sais exploiter ces réactions pour identifier un métal inconnu
\end{itemize}
\end{objectifs}

\begin{prerequis}
\begin{itemize}
\item Structure de l'atome : électrons, ions, cations (Seconde)
\item Écriture et équilibrage d'une équation-bilan (Seconde)
\item Ions en solution aqueuse et formules des solutions ioniques (Première, leçons précédentes)
\item Tests d'identification des ions et des gaz (Seconde)
\item Notion de mole et concentration molaire (Première)
\end{itemize}
\end{prerequis}

\newpage

\coursec{Action d'un acide dilué sur un métal}

Au marché Mokolo à Yaoundé, un bijoutier propose à Maman Ngo Bell un collier « en or » à bas prix. Méfiante, elle se souvient d'un truc que lui avait confié sa grand-mère : l'or véritable ne réagit pas avec les acides, contrairement au cuivre ou au zinc. Une goutte d'acide, et le faux collier se trahit en quelques secondes ! Par ailleurs, dans les ateliers de soudure de Douala, on observe qu'un clou en fer plongé dans une solution bleue de sulfate de cuivre se recouvre progressivement d'un dépôt rougeâtre. Comment expliquer ces phénomènes ? Que se passe-t-il réellement entre un métal et une solution aqueuse ? Pourquoi certains métaux sont-ils attaqués par les acides alors que d'autres restent intacts ?

Cette leçon nous donne les clés pour \cle{prévoir} et \cle{interpréter} ces réactions. Commençons par la situation la plus simple : l'action d'un acide dilué sur un métal.

\courssub{Expérience 1 : action de l'acide chlorhydrique dilué sur le zinc}

\begin{experience}[Le zinc dans l'acide chlorhydrique]
\textbf{Dispositif et protocole.} Dans un tube à essai, introduisons quelques grenailles de zinc (petits grains de zinc métallique), puis versons environ \SI{5}{mL} d'acide chlorhydrique dilué (concentration de l'ordre de \SI{1}{mol.L^{-1}}). Bouchons légèrement le tube et attendons quelques instants.

\textbf{Observations.} Trois phénomènes se manifestent simultanément :
\begin{itemize}
\item il se produit un \cle{dégagement gazeux} abondant : des bulles se forment à la surface du zinc et remontent dans la solution ;
\item le zinc \emph{disparaît progressivement} : la quantité de métal diminue au fil du temps ;
\item le tube à essai \emph{chauffe} : la réaction est \cle{exothermique}, elle libère de l'énergie thermique.
\end{itemize}

\textbf{Interprétation provisoire.} Le métal zinc est « attaqué » par l'acide : il se transforme en une espèce soluble (invisible à l'œil nu) tandis qu'un gaz se dégage. Reste à identifier ce gaz et cette espèce dissoute.
\end{experience}

\begin{popfigure}
\begin{center}
\begin{tikzpicture}[scale=0.95, every node/.style={font=\small}]
% Tube à essai principal
\draw[thick, popdark] (-0.8,0) -- (-0.8,-3.2) arc[start angle=180, end angle=360, radius=0.8] -- (0.8,0);
% Liquide
\fill[popblueL] (-0.78,-1.1) -- (-0.78,-3.15) arc[start angle=180, end angle=360, radius=0.78] -- (0.78,-1.1) -- cycle;
\draw[popblue] (-0.78,-1.1) -- (0.78,-1.1);
% Grenailles de zinc
\foreach \p in {(-0.35,-3.35),(0.1,-3.45),(0.4,-3.3),(-0.05,-3.2)}{
  \fill[popmuted!70] \p circle (0.14);
  \draw[popdark] \p circle (0.14);}
% Bulles
\foreach \p in {(-0.3,-2.4),(0.25,-2.0),(-0.1,-1.6),(0.35,-2.7),(-0.45,-1.9),(0.05,-2.9)}{
  \draw[popblue] \p circle (0.06);}
% Étiquettes tube
\node[align=left, anchor=west] at (1.3,-3.3) {grenailles\\ de zinc \ce{Zn}};
\draw[->, popdark] (1.25,-3.3) -- (0.45,-3.35);
\node[align=left, anchor=west] at (1.3,-1.6) {acide chlorhydrique\\ dilué \ce{(H3O+ + Cl-)}};
\draw[->, popdark] (1.25,-1.7) -- (0.5,-1.9);
\node[align=left, anchor=west] at (1.3,-0.4) {dégagement\\ de gaz \ce{H2}};
\draw[->, popdark] (1.25,-0.5) -- (0.2,-0.9);
% Tube test flamme
\begin{scope}[xshift=6.2cm, yshift=-1.2cm]
\draw[thick, popdark] (-0.6,0.6) -- (-0.6,-1.6) arc[start angle=180, end angle=360, radius=0.6] -- (0.6,0.6);
\fill[popblueL!60] (-0.58,-0.2) -- (-0.58,-1.55) arc[start angle=180, end angle=360, radius=0.58] -- (0.58,-0.2) -- cycle;
% flamme
\fill[poporange] (0,1.55) .. controls (0.28,1.15) and (0.18,0.95) .. (0,0.8) .. controls (-0.18,0.95) and (-0.28,1.15) .. (0,1.55);
\fill[popgold] (0,1.25) .. controls (0.12,1.05) .. (0,0.92) .. controls (-0.12,1.05) .. (0,1.25);
\draw[->, poporange, thick] (-1.6,1.1) -- (-0.35,1.1);
\node[anchor=east, align=right] at (-1.7,1.1) {on approche\\ une flamme};
\node[popink, font=\bfseries] at (0,2.3) {détonation !};
\end{scope}
\end{tikzpicture}
\end{center}
\legende{Action de l'acide chlorhydrique dilué sur le zinc : dégagement de dihydrogène, caractérisé par la détonation (« aboiement ») lorsqu'on approche une flamme.}
\end{popfigure}

\courssub{Caractérisation des produits formés}

\textbf{Identification du gaz : le dihydrogène.} Recueillons le gaz dégagé dans un tube retourné, puis approchons une flamme (une allumette enflammée) de l'ouverture du tube.

\begin{methode}[Identifier le dihydrogène]
\begin{enumerate}
\item Recueillir le gaz dans un tube à essai renversé (le dihydrogène est beaucoup moins dense que l'air, il « monte »).
\item Approcher rapidement une flamme de l'ouverture du tube.
\item Observer : une \cle{détonation} brève, une sorte d'« aboiement » sec, caractérise le dihydrogène \ce{H2}.
\end{enumerate}
Ce test repose sur la combustion explosive du mélange dihydrogène-dioxygène : $\ce{2H2 + O2 -> 2H2O}$.
\end{methode}

\textbf{Mise en évidence des ions zinc \ce{Zn^2+}.} Le zinc a disparu, mais il ne s'est pas volatilisé : il est passé en solution sous forme d'ions. Pour le prouver, filtrons le contenu du tube et ajoutons quelques gouttes de soude (hydroxyde de sodium, \ce{Na+ + HO-}) au filtrat : il apparaît un \cle{précipité blanc} d'hydroxyde de zinc :
\[ \ce{Zn^2+ + 2HO- -> Zn(OH)2 (s)} \]
Ce précipité blanc est le test caractéristique des ions \ce{Zn^2+} (il se redissout dans un excès de soude, ce qui le distingue d'autres précipités blancs).

\begin{aretenir}[Bilan de l'expérience]
L'acide chlorhydrique dilué attaque le zinc. Il se forme :
\begin{itemize}
\item du \cle{dihydrogène} \ce{H2} (gaz qui détone à la flamme) ;
\item des \cle{ions zinc} \ce{Zn^2+} en solution (précipité blanc avec la soude, soluble en excès de soude) ;
\item de la \cle{chaleur} (réaction exothermique).
\end{itemize}
\end{aretenir}

\courssub{Expérience 2 : action de l'acide sulfurique dilué sur le fer}

Recommençons l'expérience avec de la limaille de fer et de l'acide sulfurique dilué \ce{H2SO4}. Les observations sont analogues : dégagement gazeux qui détone à la flamme (donc \ce{H2}), disparition progressive du fer, échauffement du tube. Une différence notable apparaît cependant : la solution prend une \cle{teinte verdâtre} pâle, caractéristique des ions fer(II) \ce{Fe^2+}. On confirme leur présence en ajoutant quelques gouttes de soude : un précipité \emph{vert} d'hydroxyde de fer(II) se forme :
\[ \ce{Fe^2+ + 2HO- -> Fe(OH)2 (s)} \]
L'acide chlorhydrique dilué donne exactement le même résultat avec le fer : l'essentiel, c'est la présence des ions hydrogène \ce{H+} (solvatés en \ce{H3O+}), commune à tous les acides.

\begin{attention}[Précipité vert \ce{Fe(OH)2} : instabilité à l'air]
Le précipité vert d'hydroxyde de fer(II) \ce{Fe(OH)2} s'oxyde très rapidement au contact de l'air ambiant (dioxygène) et devient \emph{brun-orangé} (formation d'hydroxyde de fer(III) \ce{Fe(OH)3}). L'observation doit être immédiate.
\end{attention}

\courssub{Cas de l'aluminium, du cuivre, de l'argent et de l'or}

\textbf{L'aluminium : un métal protégé.} Plongeons une plaque d'aluminium dans l'acide chlorhydrique dilué \emph{à froid} : la réaction est très lente, presque imperceptible au début. Pourquoi ? L'aluminium est recouvert d'une fine couche d'\cle{oxyde d'aluminium} \ce{Al2O3} (alumine), invisible et très adhérente, qui le protège de l'attaque acide. C'est d'ailleurs pour cela que les marmites en aluminium résistent à l'usage quotidien. À \emph{chaud}, ou après avoir gratté la surface (ou grâce à l'action des ions chlorure qui percent la couche d'alumine), la réaction devient vive : dégagement de dihydrogène et formation d'ions \ce{Al^3+}.

\textbf{Le cuivre, l'argent et l'or : les métaux nobles.} Plongeons une lame de cuivre dans l'acide chlorhydrique dilué : \emph{rien ne se passe}, même à chaud. Même constat avec l'acide sulfurique dilué, avec l'argent et avec l'or. Ces métaux, dits \cle{métaux nobles}, ne sont pas attaqués par les acides dilués non oxydants : aucun dégagement gazeux, aucune dissolution.

\begin{attention}[Erreur fréquente : « l'acide brûle le métal »]
Il est incorrect de dire que l'acide « brûle » ou « consume » le métal. Une \emph{combustion} est une réaction avec le dioxygène de l'air. Ici, il s'agit d'une \cle{réaction chimique} entre le métal et les ions \ce{H+} de la solution, au cours de laquelle le métal \emph{cède des électrons} aux ions hydrogène. C'est un transfert d'électrons, pas une combustion. On dit que le métal est \emph{attaqué} ou \emph{dissous} par l'acide.
\end{attention}

\begin{popfigure}
\begin{center}
\begin{tabularx}{\linewidth}{|l|>{\centering}X|>{\centering}X|X|}
\hline
\rowcolor{popblueL}
\textbf{Métal} & \textbf{\ce{HCl} dilué} & \textbf{\ce{H2SO4} dilué} & \textbf{Observations} \\
\hline
Zinc \ce{Zn} & \cellcolor{popgreen!20}\textbf{Réaction vive} & \cellcolor{popgreen!20}\textbf{Réaction vive} & \ce{H2} dégagé, \ce{Zn^2+} en solution \\
\hline
Fer \ce{Fe} & \cellcolor{popgreen!20}\textbf{Réaction} & \cellcolor{popgreen!20}\textbf{Réaction} & \ce{H2} dégagé, solution verdâtre \ce{Fe^2+} \\
\hline
Aluminium \ce{Al} & \cellcolor{poporange!20}\textbf{Lente à froid} & \cellcolor{poporange!20}\textbf{Lente à froid} & Couche \ce{Al2O3} protectrice ; vive à chaud \\
\hline
Cuivre \ce{Cu} & \cellcolor{poppink!20}\textbf{Aucune} & \cellcolor{poppink!20}\textbf{Aucune} & Métal non attaqué (noble) \\
\hline
Argent \ce{Ag} & \cellcolor{poppink!20}\textbf{Aucune} & \cellcolor{poppink!20}\textbf{Aucune} & Métal noble \\
\hline
Or \ce{Au} & \cellcolor{poppink!20}\textbf{Aucune} & \cellcolor{poppink!20}\textbf{Aucune} & Métal noble \\
\hline
\end{tabularx}
\end{center}
\legende{Action des acides chlorhydrique et sulfurique dilués sur quelques métaux usuels.}
\end{popfigure}

\courssub{Conclusion générale et équations-bilans}

\begin{definition}[Attaque d'un métal par un acide dilué]
Un métal est \cle{attaqué} par un acide dilué (non oxydant) s'il est capable de \emph{céder des électrons} aux ions hydrogène \ce{H+} de la solution. Le métal passe alors en solution sous forme de \cle{cations métalliques} pendant que les ions \ce{H+} sont transformés en \cle{dihydrogène} \ce{H2} qui se dégage. Les métaux qui ne peuvent pas effectuer ce transfert (cuivre, argent, or) ne réagissent pas.
\end{definition}

Avant d'écrire l'équation-bilan globale, mettons en évidence les deux demi-réactions électroniques (demi-équations) qui se superposent : l'\cle{oxydation} du métal et la \cle{réduction} des ions hydrogène.

\begin{propriete}[Demi-équations et équation-bilan pour l'attaque du zinc]
\begin{itemize}
\item \textbf{Demi-équation d'oxydation (perte d'électrons) :} $\ce{Zn (s) -> Zn^2+ (aq) + 2e-}$
\item \textbf{Demi-équation de réduction (gain d'électrons) :} $\ce{2H+ (aq) + 2e- -> H2 (g)}$
\item \textbf{Équation-bilan ionique (somme, électrons simplifiés) :}
\[ \ce{Zn (s) + 2H+ (aq) -> Zn^2+ (aq) + H2 (g)} \]
\end{itemize}
Vérification : atomes (1 Zn, 2 H) et charges ($+2$ des deux côtés) équilibrés. Réaction rapide, exothermique.
\end{propriete}

\begin{propriete}[Demi-équations et équation-bilan pour l'attaque du fer]
\begin{itemize}
\item \textbf{Oxydation :} $\ce{Fe (s) -> Fe^2+ (aq) + 2e-}$
\item \textbf{Réduction :} $\ce{2H+ (aq) + 2e- -> H2 (g)}$
\item \textbf{Bilan :}
\[ \ce{Fe (s) + 2H+ (aq) -> Fe^2+ (aq) + H2 (g)} \]
\end{itemize}
\end{propriete}

Pour l'aluminium, chaque atome cède 3 électrons et devient \ce{Al^3+} ; il faut ajuster les coefficients pour égaler les électrons (PPCM de 2 et 3 = 6 électrons) :
\begin{itemize}
\item \textbf{Oxydation ($\times 2$) :} $\ce{2Al (s) -> 2Al^3+ (aq) + 6e-}$
\item \textbf{Réduction ($\times 3$) :} $\ce{6H+ (aq) + 6e- -> 3H2 (g)}$
\item \textbf{Bilan :}
\[ \ce{2Al (s) + 6H+ (aq) -> 2Al^3+ (aq) + 3H2 (g)} \]
\end{itemize}
Vérification : 2 Al, 6 H de chaque côté ; charges : $+6$ à gauche, $2\times(+3)=+6$ à droite.

\textbf{Et les ions chlorure et sulfate ?} Dans la solution d'acide chlorhydrique, les ions \ce{Cl-} sont présents mais n'apparaissent pas dans l'équation-bilan ionique : ils ne participent pas à la transformation. Ce sont des \cle{ions spectateurs}. Il en va de même pour les ions sulfate \ce{SO4^2-} de l'acide sulfurique. Si l'on veut écrire les \emph{équations moléculaires complètes} (avec toutes les espèces), on obtient :
\[ \ce{Zn (s) + 2H3O+ (aq) + 2Cl- (aq) -> Zn^2+ (aq) + 2Cl- (aq) + H2 (g) + 2H2O (l)} \]
\[ \ce{Fe (s) + 2H3O+ (aq) + SO4^2- (aq) -> Fe^2+ (aq) + SO4^2- (aq) + H2 (g) + 2H2O (l)} \]
En simplifiant par les ions spectateurs (et les molécules d'eau liées à \ce{H3O+}), on retrouve les équations ioniques précédentes.

\begin{exoresolu}[Volume de dihydrogène dégagé par l'attaque du zinc]
On attaque \SI{6,5}{g} de zinc par un excès d'acide chlorhydrique dilué. On donne $M(\ce{Zn}) = \SI{65}{g.mol^{-1}}$ et le volume molaire dans les CNTP $V_m = \SI{24}{L.mol^{-1}}$. Déterminer la quantité puis le volume de dihydrogène dégagé.
\tcblower
\textbf{Étape 1 : équation-bilan ionique.}
\[ \ce{Zn + 2H+ -> Zn^2+ + H2} \]

\textbf{Étape 2 : quantité de matière de zinc.}
\[ n(\ce{Zn}) = \frac{m}{M} = \frac{6,5}{65} = \SI{0,10}{mol} \]

\textbf{Étape 3 : tableau d'avancement} (l'acide est en excès, le zinc est le réactif limitant).
\begin{center}
\begin{tabular}{|c|cccc|}
\hline
& \ce{Zn} & \ce{2H+} & \ce{Zn^2+} & \ce{H2} \\
\hline
État initial (mol) & $0{,}10$ & excès & $0$ & $0$ \\
Avancement (mol) & $-x$ & $-2x$ & $+x$ & $+x$ \\
État final (mol) & $0{,}10-x$ & excès & $x$ & $x$ \\
\hline
\end{tabular}
\end{center}
Le zinc disparaît totalement : $x_{max} = \SI{0,10}{mol}$, d'où $n(\ce{H2}) = \SI{0,10}{mol}$.

\textbf{Étape 4 : volume de gaz.}
\[ V(\ce{H2}) = n \times V_m = 0,10 \times 24 = \SI{2,4}{L} \]
L'attaque de \SI{6,5}{g} de zinc dégage donc \SI{2,4}{L} de dihydrogène dans les CNTP, soit environ le volume d'une grande bouteille d'eau !
\end{exoresolu}

\begin{savaistu}[La pierre de touche des bijoutiers]
C'est précisément parce que l'or et l'argent résistent aux acides dilués que les bijoutiers utilisent des acides pour tester l'authenticité des bijoux ! On frotte le bijou sur une « pierre de touche » (une pierre noire très fine), puis on dépose une goutte d'acide sur la trace laissée : si la trace se dissout, le métal n'est pas de l'or pur ; si elle résiste, c'est bon signe. Maman Ngo Bell avait donc raison de se méfier : un collier « en or » attaqué par l'acide est en réalité un alliage de cuivre ou de zinc doré en surface. À l'inverse, l'or ne réagit qu'avec un mélange très particulier, l'\emph{eau régale} (acide chlorhydrique concentré + acide nitrique concentré), capable d'oxyder même les métaux les plus nobles.
\end{savaistu}

\begin{exercice}[À toi de jouer]
On plonge \SI{2,8}{g} de limaille de fer dans un excès d'acide sulfurique dilué. On donne $M(\ce{Fe}) = \SI{56}{g.mol^{-1}}$ et $V_m = \SI{24}{L.mol^{-1}}$.
\begin{enumerate}
\item Écrire les deux demi-équations électroniques et l'équation-bilan ionique de la réaction.
\item Calculer la quantité de matière de fer utilisée.
\item En déduire la quantité puis le volume de dihydrogène dégagé dans les CNTP.
\item Quel test permet de mettre en évidence les ions formés en solution ? Qu'observe-t-on ? Préciser une particularité de ce précipité.
\end{enumerate}
\end{exercice}

\begin{corrige}[Exercice n° 1]
\textbf{1.} 
\begin{itemize}
\item Oxydation : $\ce{Fe (s) -> Fe^2+ (aq) + 2e-}$
\item Réduction : $\ce{2H+ (aq) + 2e- -> H2 (g)}$
\item Bilan : $\ce{Fe (s) + 2H+ (aq) -> Fe^2+ (aq) + H2 (g)}$
\end{itemize}

\textbf{2.} $n(\ce{Fe}) = \dfrac{m}{M} = \dfrac{2,8}{56} = \SI{0,050}{mol}$.

\textbf{3.} D'après l'équation, 1 mol Fe produit 1 mol \ce{H2}. L'acide est en excès, le fer est limitant.
$n(\ce{H2}) = n(\ce{Fe}) = \SI{0,050}{mol}$.
\[ V(\ce{H2}) = 0,050 \times 24 = \SI{1,2}{L} \]

\textbf{4.} On ajoute quelques gouttes de soude à la solution : un précipité \emph{vert} d'hydroxyde de fer(II) \ce{Fe(OH)2} apparaît, ce qui prouve la présence des ions \ce{Fe^2+} (la solution était déjà verdâtre). \textbf{Particularité :} ce précipité vert noircit rapidement à l'air (oxydation en \ce{Fe(OH)3} brun).
\end{corrige}

En résumé, les acides chlorhydrique et sulfurique dilués attaquent le zinc, le fer et (plus difficilement) l'aluminium, avec dégagement de dihydrogène et formation de cations métalliques ; en revanche, le cuivre, l'argent et l'or restent inertes. Mais pourquoi cette différence de comportement ? Que s'échange-t-il réellement entre le métal et les ions \ce{H+} ? La réponse tient en deux mots : des \emph{électrons}. C'est ce que nous formaliserons dans la suite de la leçon avec les notions d'oxydation et de réduction.

\coursec{Réaction entre un ion métallique et un métal}

Dans la section précédente, nous avons vu que certains métaux (zinc, fer, aluminium) sont attaqués par les acides dilués : le métal disparaît peu à peu tandis qu'un dégagement de dihydrogène se produit. Nous avons interprété ce phénomène comme un transfert d'électrons du métal vers les ions \ce{H+}. Une question surgit alors naturellement : les ions \emph{métalliques} présents dans une solution, comme les ions cuivre \ce{Cu^2+} qui donnent sa belle couleur bleue au sulfate de cuivre, peuvent-ils eux aussi « capter » des électrons d'un métal ? Autrement dit, un métal peut-il réagir avec les ions d'un \emph{autre} métal ? C'est toute la question de cette section, et la réponse va nous offrir un outil précieux : comparer les métaux entre eux et prévoir leurs réactions.

\courssub{Expérience 1 : le zinc plongé dans une solution de sulfate de cuivre}

\begin{experience}[Lame de zinc dans une solution de sulfate de cuivre]
\textbf{Dispositif et protocole.} On plonge une lame de zinc bien décapée (grattée à la toile émeri pour enlever la couche d'oxyde) dans un bécher contenant une solution bleue de sulfate de cuivre (\ce{Cu^2+} + \ce{SO4^2-}). On attend quelques minutes à température ambiante, sans chauffer ni agiter fortement.

\textbf{Observations.}
\begin{itemize}
\item La lame de zinc se recouvre progressivement d'un \cle{dépôt rouge-brique}, spongieux, qui adhère au métal.
\item La coloration \cle{bleue} de la solution \emph{pâlit} peu à peu ; si l'on attend assez longtemps, elle finit par se décolorer complètement.
\item La lame de zinc s'use : sa masse diminue.
\end{itemize}

\textbf{Tests d'identification.} Une goutte de solution prélevée en fin d'expérience et traitée par la soude donne un précipité blanc gélatineux d'hydroxyde de zinc \ce{Zn(OH)2} : des ions \ce{Zn^2+} sont apparus en solution. Le dépôt rouge, raclé et examiné, présente l'éclat et la couleur caractéristiques du \cle{cuivre métallique}.

\textbf{Interprétation.} Les ions \ce{Cu^2+} de la solution ont capté des électrons cédés par le zinc : ils se transforment en cuivre métallique \ce{Cu} qui se dépose, tandis que le zinc passe en solution sous forme d'ions \ce{Zn^2+}. Les ions sulfate \ce{SO4^2-} ne participent pas : ce sont des ions \emph{spectateurs}.
\end{experience}

\begin{popfigure}
\begin{tikzpicture}[scale=0.9]
% ===== AVANT =====
\begin{scope}
% bécher
\draw[thick,popdark] (-1.4,0) -- (-1.4,2.6) (-1.4,0) -- (1.4,0) (1.4,0) -- (1.4,2.6);
\draw[thick,popdark] (-1.55,2.6) -- (-1.4,2.6) (1.4,2.6) -- (1.55,2.6);
% solution bleue
\fill[popblue!35] (-1.36,0.04) rectangle (1.36,1.9);
\draw[popblue,thick] (-1.36,1.9) -- (1.36,1.9);
\node[popblue!70!black] at (0,0.45) {\small \ce{Cu^2+} + \ce{SO4^2-}};
% lame de zinc
\fill[popmuted!60] (-0.22,0.7) rectangle (0.22,3.1);
\draw[popdark] (-0.22,0.7) rectangle (0.22,3.1);
\node[popdark] at (0,3.35) {\small Zn};
\node at (0,-0.55) {\textbf{Avant}};
\end{scope}
% flèche
\draw[-{Stealth[length=3mm]},very thick,poporange] (2.1,1.5) -- (3.3,1.5);
% ===== APRÈS =====
\begin{scope}[xshift=6.9cm]
\draw[thick,popdark] (-1.4,0) -- (-1.4,2.6) (-1.4,0) -- (1.4,0) (1.4,0) -- (1.4,2.6);
\draw[thick,popdark] (-1.55,2.6) -- (-1.4,2.6) (1.4,2.6) -- (1.55,2.6);
% solution décolorée
\fill[popblue!8] (-1.36,0.04) rectangle (1.36,1.9);
\draw[popblue!50,thick] (-1.36,1.9) -- (1.36,1.9);
\node[popmuted] at (0,0.45) {\small \ce{Zn^2+} + \ce{SO4^2-}};
% lame amincie + dépôt rouge
\fill[popmuted!60] (-0.15,0.7) rectangle (0.15,3.1);
\draw[popdark] (-0.15,0.7) rectangle (0.15,3.1);
\fill[poporange!80!red,decorate,decoration={random steps,segment length=2pt,amplitude=1pt}] (-0.30,0.7) rectangle (-0.15,1.9);
\fill[poporange!80!red,decorate,decoration={random steps,segment length=2pt,amplitude=1pt}] (0.15,0.7) rectangle (0.30,1.9);
\node[poporange!80!red] at (1.05,1.35) {\small \ce{Cu}};
\node at (0,-0.55) {\textbf{Après}};
\end{scope}
\end{tikzpicture}
\legende{Lame de zinc plongée dans une solution de sulfate de cuivre : dépôt rouge de cuivre sur le zinc et décoloration progressive de la solution bleue.}
\end{popfigure}

L'équation-bilan de la réaction s'écrit :

\[ \ce{Zn(s) + Cu^2+(aq) -> Zn^2+(aq) + Cu(s)} \]

Cette équation est équilibrée \emph{en atomes} (un zinc, un cuivre de chaque côté) et \emph{en charges} ($+2$ de chaque côté). Elle est \cle{spontanée}, se déroule à température ambiante, sans catalyseur, et elle est pratiquement \emph{totale} : elle s'arrête seulement quand l'un des réactifs est épuisé. Remarque essentielle : on n'observe \textbf{aucun dégagement gazeux}, contrairement à l'action d'un acide. Le seul transfert qui a lieu est un transfert d'électrons du zinc vers les ions cuivre.

\begin{definition}[Réaction entre un métal et un ion métallique]
La réaction entre un métal et un ion métallique en solution aqueuse est un \cle{transfert d'électrons} du métal vers l'ion métallique : le métal perd des électrons et passe en solution sous forme d'ions, tandis que l'ion métallique gagne des électrons et se transforme en métal qui se dépose. Une telle réaction est une réaction d'\cle{oxydoréduction}.
\end{definition}

\courssub{Expérience 2 : le fer plongé dans une solution de sulfate de cuivre}

\begin{experience}[Clou en fer dans une solution de sulfate de cuivre]
\textbf{Protocole.} On plonge un clou en fer, préalablement décapé, dans une solution bleue de sulfate de cuivre. On laisse reposer une dizaine de minutes.

\textbf{Observations.}
\begin{itemize}
\item Le clou se recouvre d'un \cle{dépôt rouge} de cuivre métallique : on dit parfois que le clou s'est « cuivré ».
\item La solution bleue pâlit et prend progressivement une teinte \emph{verdâtre} très pâle, caractéristique des ions \ce{Fe^2+} en solution diluée.
\item Le test à la soude sur un prélèvement de solution donne un précipité vert d'hydroxyde de fer(II) \ce{Fe(OH)2}, confirmant la présence d'ions \ce{Fe^2+}.
\end{itemize}

\textbf{Interprétation.} Le fer a cédé des électrons aux ions \ce{Cu^2+} : le fer passe en solution sous forme d'ions \ce{Fe^2+} et le cuivre se dépose sur le clou. L'équation-bilan s'écrit :
\[ \ce{Fe(s) + Cu^2+(aq) -> Fe^2+(aq) + Cu(s)} \]
Réaction spontanée, totale, à température ambiante, sans catalyseur ni dégagement gazeux.
\end{experience}

\begin{popfigure}
\begin{tikzpicture}[scale=0.9]
% bécher
\draw[thick,popdark] (-1.5,0) -- (-1.5,2.8) (-1.5,0) -- (1.5,0) (1.5,0) -- (1.5,2.8);
\draw[thick,popdark] (-1.65,2.8) -- (-1.5,2.8) (1.5,2.8) -- (1.65,2.8);
% solution bleue qui verdit
\fill[popblue!25] (-1.46,0.04) rectangle (1.46,2.0);
\fill[popgreen!12] (-1.46,0.04) rectangle (1.46,0.9);
\draw[popblue,thick] (-1.46,2.0) -- (1.46,2.0);
\node[popblue!70!black] at (0,2.35) {\small solution de sulfate de cuivre};
% clou incliné
\begin{scope}[rotate=18]
\fill[popmuted!70] (-0.12,-0.4) rectangle (0.12,3.3);
\draw[popdark] (-0.12,-0.4) rectangle (0.12,3.3);
% tête du clou
\fill[popmuted!70] (-0.28,3.3) rectangle (0.28,3.45);
\draw[popdark] (-0.28,3.3) rectangle (0.28,3.45);
% dépôt de cuivre sur la partie immergée
\fill[poporange!80!red,decorate,decoration={random steps,segment length=2pt,amplitude=1pt}] (-0.24,-0.4) rectangle (-0.12,1.6);
\fill[poporange!80!red,decorate,decoration={random steps,segment length=2pt,amplitude=1pt}] (0.12,-0.4) rectangle (0.24,1.6);
\end{scope}
\node[poporange!80!red] at (1.9,0.7) {\small dépôt de \ce{Cu}};
\draw[-{Stealth},poporange!80!red] (1.55,0.75) -- (0.75,0.85);
\node[popgreen!60!black] at (-2.15,0.5) {\small ions \ce{Fe^2+}};
\draw[-{Stealth},popgreen!60!black] (-1.75,0.55) -- (-0.95,0.6);
\node[popdark] at (-0.9,3.6) {\small clou en fer};
\end{tikzpicture}
\legende{Clou en fer plongé dans une solution de sulfate de cuivre : le clou se « cuivre » et des ions \ce{Fe^2+} apparaissent en solution.}
\end{popfigure}

\begin{attention}[Le dépôt rouge n'est pas de la rouille !]
Erreur très fréquente chez les élèves : voir un dépôt rougeâtre sur le clou et conclure que « le clou a rouillé ». C'est faux. La \cle{rouille} est un oxyde de fer hydraté, de couleur brun-rouille, qui se forme par action lente de l'air humide sur le fer. Ici, le dépôt est \emph{rouge cuivré}, brillant par endroits, et il apparaît en quelques minutes : c'est du \cle{cuivre métallique} \ce{Cu}. D'ailleurs, la rouille ne peut pas se former ici : le fer \emph{disparaît} en donnant des ions \ce{Fe^2+}, il ne s'oxyde pas en rouille. Retiens : dépôt rouge sur un métal plongé dans une solution de \ce{Cu^2+} = cuivre métallique.
\end{attention}

\courssub{Expérience 3 : les témoins négatifs — le cuivre face aux ions \ce{Zn^2+} et \ce{Fe^2+}}

Faisons maintenant l'expérience « dans l'autre sens ». Puisque le zinc et le fer déposent le cuivre de ses ions, le cuivre peut-il, à son tour, déposer le zinc ou le fer à partir de leurs ions ?

\begin{experience}[Cuivre dans des solutions de \ce{Zn^2+} et de \ce{Fe^2+}]
\textbf{Protocole.} On plonge une lame de cuivre décapée dans une solution de sulfate de zinc (\ce{Zn^2+} + \ce{SO4^2-}), puis une autre lame de cuivre dans une solution de sulfate de fer(II) (\ce{Fe^2+} + \ce{SO4^2-}). On attend longuement, éventuellement en chauffant modérément.

\textbf{Observation.} Dans les deux cas, \cle{aucune réaction observable} : pas de dépôt sur le cuivre, pas de changement de couleur des solutions, la masse de la lame de cuivre reste inchangée.

\textbf{Conclusion.} Le cuivre métallique \emph{ne réduit pas} les ions \ce{Zn^2+} ni les ions \ce{Fe^2+}. Les réactions inverses de celles des expériences 1 et 2 ne se produisent pas :
\[ \ce{Cu(s) + Zn^2+(aq) ->} \text{rien} \qquad \ce{Cu(s) + Fe^2+(aq) ->} \text{rien} \]
\end{experience}

\courssub{Généralisation : la réaction n'a lieu que dans un sens}

\begin{propriete}[Sens unique de la réaction métal / ion métallique]
Un métal peut \cle{réduire} les ions d'un autre métal, mais la réaction n'a lieu que \emph{dans un seul sens} : le zinc et le fer réduisent les ions \ce{Cu^2+}, mais le cuivre ne réduit ni les ions \ce{Zn^2+} ni les ions \ce{Fe^2+}. On dit que le zinc et le fer ont un \cle{pouvoir réducteur} plus grand que celui du cuivre. On note cette hiérarchie :
\[ \text{pouvoir réducteur : } \ce{Zn} > \ce{Fe} > \ce{Cu} \]
De manière générale, une réaction spontanée a lieu entre un métal $M$ et les ions d'un métal $M'$ \textbf{si et seulement si} $M$ est \emph{plus réducteur} que $M'$.
\end{propriete}

\begin{methode}[Prévoir si une réaction a lieu entre un métal et un ion métallique]
\begin{enumerate}
\item \textbf{Identifier} les deux métaux en présence : le métal solide $M$ et le métal $M'$ dont les ions $M'^{n+}$ sont en solution.
\item \textbf{Comparer} leurs pouvoirs réducteurs à l'aide des expériences de référence (ou, plus tard, de la classification électrochimique).
\item \textbf{Conclure} : si $M$ est plus réducteur que $M'$, la réaction $M(s) + M'^{n+}(aq) \rightarrow M^{n+}(aq) + M'(s)$ a lieu ; sinon, il ne se passe rien.
\item \textbf{Vérifier} la cohérence : le métal solide doit \emph{disparaître} (ou s'user) et l'autre métal doit \emph{se déposer}.
\end{enumerate}
\end{methode}

\begin{exemplebox}[Prévisions rapides]
\begin{itemize}
\item Zinc + ions \ce{Fe^2+} : le zinc est plus réducteur que le fer, donc $\ce{Zn(s) + Fe^2+(aq) -> Zn^2+(aq) + Fe(s)}$ : dépôt gris de fer sur le zinc.
\item Fer + ions \ce{Zn^2+} : le fer est moins réducteur que le zinc : \emph{aucune réaction}.
\item Cuivre + ions \ce{Ag+} : le cuivre est plus réducteur que l'argent, donc $\ce{Cu(s) + 2Ag+(aq) -> Cu^2+(aq) + 2Ag(s)}$ : dépôt gris-argenté d'argent (c'est la célèbre expérience de « l'arbre de Diane »).
\end{itemize}
\end{exemplebox}

\courssub{Lien avec l'action des acides sur les métaux}

Reprenons les résultats de la section précédente. Les métaux attaqués par les acides chlorhydrique et sulfurique dilués sont le zinc, le fer, l'aluminium : ce sont précisément ceux qui \emph{réduisent les ions \ce{H+}} en dihydrogène \ce{H2}. Le cuivre, l'argent et l'or, qui ne réduisent pas \ce{H+}, ne sont pas attaqués. Tout se passe donc comme si l'ion \ce{H+} se comportait, dans cette comparaison, comme un ion métallique intermédiaire : les métaux plus réducteurs que l'hydrogène (au sens de ces expériences) réduisent \ce{H+} ; les métaux moins réducteurs (cuivre, argent, or) ne le réduisent pas.

\begin{aretenir}[Une même logique pour toutes ces réactions]
Action d'un acide sur un métal et action d'un ion métallique sur un autre métal obéissent à la \emph{même règle} : c'est toujours le \cle{plus réducteur} qui cède ses électrons. Le zinc réduit \ce{H+} et \ce{Cu^2+} ; le cuivre ne réduit ni \ce{H+}, ni \ce{Zn^2+}, ni \ce{Fe^2+}. Une seule hiérarchie des pouvoirs réducteurs permet donc de prévoir \emph{toutes} ces réactions à la fois.
\end{aretenir}

\courssub{Exemple chiffré : quelle masse de cuivre une lame de zinc peut-elle déposer ?}

\begin{exoresolu}[Zinc et sulfate de cuivre : calcul de la masse de cuivre déposée]
On plonge une lame de zinc de masse $m_{\ce{Zn}} = \SI{3,27}{g}$ dans un volume important d'une solution de sulfate de cuivre en large excès. On donne $M(\ce{Zn}) = \SI{65,4}{g.mol^{-1}}$ et $M(\ce{Cu}) = \SI{63,5}{g.mol^{-1}}$. Quelle est la masse maximale de cuivre que cette lame peut déposer ?
\tcblower
\textbf{Étape 1 : quantité de matière de zinc disponible.}
\[ n(\ce{Zn}) = \frac{m_{\ce{Zn}}}{M(\ce{Zn})} = \frac{3,27}{65,4} = \SI{0,050}{mol} \]

\textbf{Étape 2 : exploitation de l'équation-bilan.}
\[ \ce{Zn(s) + Cu^2+(aq) -> Zn^2+(aq) + Cu(s)} \]
Les coefficients stœchiométriques sont tous égaux à 1 : \emph{une mole de zinc dépose exactement une mole de cuivre}. Comme \ce{Cu^2+} est en excès, le zinc est le réactif limitant et il est entièrement consommé :
\[ n(\ce{Cu})_{\text{déposé}} = n(\ce{Zn})_{\text{consommé}} = \SI{0,050}{mol} \]

\textbf{Étape 3 : masse de cuivre déposée.}
\[ m(\ce{Cu}) = n(\ce{Cu}) \times M(\ce{Cu}) = 0,050 \times 63,5 = \SI{3,18}{g} \]

\textbf{Conclusion.} La lame de zinc de \SI{3,27}{g} peut déposer au maximum \SI{3,18}{g} de cuivre. Remarque intéressante : la masse du système « lame + dépôt » diminue légèrement ($3,18 < 3,27$), car une mole de cuivre est un peu moins lourde qu'une mole de zinc. La différence de masse de la lame, $3,27 - 3,18 = \SI{0,09}{g}$, peut d'ailleurs servir à \emph{vérifier expérimentalement} que la réaction a bien eu lieu.
\end{exoresolu}

\courssub{Application : identification d'un métal inconnu}

La famille de situations de cette leçon est la \emph{prévision des réactions chimiques}, et l'une de ses applications concrètes est l'\cle{identification d'un métal inconnu}. Supposons qu'un atelier de Douala reçoive un lot de tiges métalliques grises, non étiquetées : s'agit-il de zinc ou de fer ? Les deux métaux se ressemblent visuellement. Une simple solution de sulfate de cuivre ne permet pas de les distinguer (tous deux se couvrent de cuivre). En revanche, une solution contenant des ions \ce{Fe^2+} tranche immédiatement :

\begin{itemize}
\item si la tige se recouvre d'un dépôt gris de fer, c'est du zinc : $\ce{Zn(s) + Fe^2+(aq) -> Zn^2+(aq) + Fe(s)}$ ;
\item si rien ne se produit, c'est du fer : le fer ne peut pas réduire ses propres ions.
\end{itemize}

On peut compléter l'identification par l'action d'un acide dilué (dégagement de \ce{H2} dans les deux cas, mais à des vitesses différentes) puis par des tests d'ions en solution (précipité blanc de \ce{Zn(OH)2} ou vert de \ce{Fe(OH)2} avec la soude). Cette démarche — formuler une hypothèse, choisir un réactif qui donne des résultats \emph{différents} selon l'hypothèse, conclure — est au cœur du savoir-faire expérimental attendu au Probatoire.

\begin{savaistu}[De la métallurgie à la récupération des métaux précieux]
Le principe « un métal réducteur dépose un métal moins réducteur de ses ions » est exploité industriellement. En \cle{métallurgie}, on récupère certains métaux de leurs solutions par \emph{cémentation} : on fait passer des solutions contenant des ions \ce{Cu^2+} (issues du traitement des minerais) sur de la ferraille ; le fer se dissout et le cuivre se dépose — c'est exactement l'expérience du clou, à grande échelle ! De même, les ateliers de traitement de surface et les laboratoires récupèrent l'\cle{argent} des solutions usées (bains de photographie, de dorure) en y plongeant du cuivre : $\ce{Cu(s) + 2Ag+(aq) -> Cu^2+(aq) + 2Ag(s)}$. Ce même principe explique aussi pourquoi il ne faut jamais verser une solution d'ions cuivre dans une canalisation en zinc ou en fer galvanisé : le tuyau serait littéralement « mangé » par la réaction d'oxydoréduction.
\end{savaistu}

\begin{exercice}[À toi de jouer]
On dispose de trois béchers contenant respectivement des solutions de \ce{CuSO4}, de \ce{ZnSO4} et de \ce{FeSO4}, et de trois lames : une de zinc, une de fer, une de cuivre.
\begin{enumerate}
\item Prévoir, en justifiant, toutes les réactions possibles entre une lame et une solution.
\item On plonge une lame de fer de masse \SI{2,80}{g} dans un excès de solution de sulfate de cuivre. Calculer la masse de cuivre déposée et la masse finale de la lame (solide) en fin de réaction. On donne $M(\ce{Fe}) = \SI{56,0}{g.mol^{-1}}$ et $M(\ce{Cu}) = \SI{63,5}{g.mol^{-1}}$.
\end{enumerate}
\end{exercice}

\begin{corrige}[Exercice : lame de fer dans le sulfate de cuivre]
\textbf{1.} D'après la hiérarchie $\ce{Zn} > \ce{Fe} > \ce{Cu}$ des pouvoirs réducteurs, les seules réactions possibles sont :
\begin{itemize}
\item $\ce{Zn(s) + Cu^2+(aq) -> Zn^2+(aq) + Cu(s)}$ et $\ce{Zn(s) + Fe^2+(aq) -> Zn^2+(aq) + Fe(s)}$ (le zinc réduit les ions des deux autres métaux) ;
\item $\ce{Fe(s) + Cu^2+(aq) -> Fe^2+(aq) + Cu(s)}$ (le fer réduit seulement les ions cuivre) ;
\item le cuivre ne réagit avec aucune des trois solutions.
\end{itemize}

\textbf{2.} Quantité de fer : $n(\ce{Fe}) = \dfrac{2,80}{56,0} = \SI{0,050}{mol}$. D'après $\ce{Fe(s) + Cu^2+(aq) -> Fe^2+(aq) + Cu(s)}$, il se dépose $n(\ce{Cu}) = \SI{0,050}{mol}$, soit $m(\ce{Cu}) = 0,050 \times 63,5 = \SI{3,18}{g}$. Le fer est entièrement consommé : la lame de fer disparaît, et le solide restant est le cuivre déposé, de masse finale $\SI{3,18}{g}$ (supérieure à la masse initiale de la lame, car une mole de cuivre est plus lourde qu'une mole de fer).
\end{corrige}

En résumé, les expériences de cette section nous ont révélé une \emph{hiérarchie} entre les métaux : certains cèdent facilement leurs électrons (zinc, fer), d'autres beaucoup moins (cuivre, argent, or). Mais pourquoi parle-t-on d'\emph{oxydation} et de \emph{réduction} ? Que signifient exactement les termes \emph{oxydant} et \emph{réducteur}, et comment écrire proprement les transferts d'électrons que nous venons de décrire ? C'est l'objet de la section suivante, où nous donnerons à ces phénomènes leur cadre théorique rigoureux : les définitions de l'oxydoréduction et les demi-équations électroniques.

\coursec{Oxydation, réduction, oxydant, réducteur}

Dans les sections précédentes, nous avons observé deux familles de phénomènes : un métal plongé dans un acide dilué se dissout en dégageant du dihydrogène, et un métal plongé dans une solution d'ions métalliques se recouvre d'un autre métal. Nous avons écrit les équations-bilans, par exemple :

\[ \ce{Zn(s) + Cu^2+(aq) -> Zn^2+(aq) + Cu(s)} \]

Mais une question fondamentale reste posée : \textbf{que s'est-il réellement passé au niveau microscopique ?} Pourquoi l'atome de zinc s'est-il transformé en ion \ce{Zn^2+} ? Pourquoi l'ion \ce{Cu^2+} s'est-il transformé en atome de cuivre ? La réponse tient en un mot : \cle{électrons}. Toute cette leçon repose sur un échange invisible mais bien réel : le \cle{transfert d'électrons} d'une espèce chimique vers une autre. C'est ce transfert que nous allons maintenant décortiquer, nommer et généraliser.

\courssub{Analyse électronique de la réaction entre le zinc et les ions cuivre(II)}

Reprenons l'expérience de la section précédente : une lame de zinc plongée dans une solution bleue de sulfate de cuivre(II). La lame se recouvre d'un dépôt rougeâtre de cuivre, la solution se décolore progressivement, et le zinc se consume. L'équation-bilan est :

\[ \ce{Zn(s) + Cu^2+(aq) -> Zn^2+(aq) + Cu(s)} \]

Regardons chaque espèce séparément, en comptant les charges.

\courssub{Ce qui arrive au zinc : une perte d'électrons}

L'atome de zinc \ce{Zn} est électriquement neutre. À la fin de la réaction, il est devenu un ion \ce{Zn^2+}, porteur de \textbf{deux charges positives}. Comment un atome neutre peut-il devenir chargé positivement ? En \textbf{perdant deux électrons} (particules chargées négativement). On écrit :

\[ \ce{Zn(s) -> Zn^2+(aq) + 2e-} \]

Vérifions l'équilibre : à gauche, charge totale $0$ ; à droite, $(+2) + 2\times(-1) = 0$. Les charges sont conservées, les éléments aussi. Cette écriture est une \cle{demi-équation électronique}.

\courssub{Ce qui arrive à l'ion cuivre(II) : un gain d'électrons}

De son côté, l'ion \ce{Cu^2+}, porteur de deux charges positives, devient un atome de cuivre \ce{Cu} neutre. Il a donc \textbf{gagné deux électrons} :

\[ \ce{Cu^2+(aq) + 2e- -> Cu(s)} \]

\courssub{Le bilan : un transfert direct d'électrons}

Les deux électrons perdus par le zinc sont \textbf{exactement} ceux gagnés par l'ion cuivre(II). Il n'apparaît aucun électron libre dans l'équation-bilan : les électrons ne se promènent pas seuls dans la solution, ils passent directement du donneur (le zinc) à l'accepteur (l'ion \ce{Cu^2+}) lors des chocs entre le métal et les ions en solution. En additionnant les deux demi-équations, les électrons s'éliminent :

\begin{align*}
\ce{Zn(s) &-> Zn^2+(aq) + 2e-} \\
\ce{Cu^2+(aq) + 2e- &-> Cu(s)} \\ \hline
\ce{Zn(s) + Cu^2+(aq) &-> Zn^2+(aq) + Cu(s)}
\end{align*}

\begin{popfigure}
\begin{center}
\begin{tikzpicture}[scale=1.0, every node/.style={font=\small}]
% Atome de zinc
\draw[fill=popblue!35, draw=popblue, thick] (0,0) circle (0.85);
\node[font=\bfseries\large, popdark] at (0,0) {\ce{Zn}};
\node[below=2pt, popdark] at (0,-0.95) {atome de zinc (neutre)};
% Ion Cu2+
\draw[fill=poporange!35, draw=poporange, thick] (6.5,0) circle (0.85);
\node[font=\bfseries\large, popdark] at (6.5,0) {\ce{Cu^2+}};
\node[below=2pt, popdark] at (6.5,-0.95) {ion cuivre(II)};
% Flèches d'électrons
\draw[-{Stealth[length=3mm]}, very thick, poppink] (0.9,0.45) .. controls (2.5,1.6) and (4,1.6) .. (5.6,0.45)
  node[midway, above, poppink, font=\bfseries] {$1\,e^-$};
\draw[-{Stealth[length=3mm]}, very thick, poppink] (0.9,-0.45) .. controls (2.5,-1.7) and (4,-1.7) .. (5.6,-0.45)
  node[midway, below, poppink, font=\bfseries] {$1\,e^-$};
% Annotations
\node[align=center, popblue, font=\bfseries] at (0,1.7) {perd $2\,e^-$ \\ (donneur)};
\node[align=center, poporange, font=\bfseries] at (6.5,1.7) {gagne $2\,e^-$ \\ (accepteur)};
\end{tikzpicture}
\end{center}
\legende{Transfert de deux électrons d'un atome de zinc vers un ion cuivre(II) : le zinc devient \ce{Zn^2+}, l'ion \ce{Cu^2+} devient \ce{Cu}.}
\end{popfigure}

\begin{aretenir}[Idée fondamentale]
Dans la réaction $\ce{Zn + Cu^2+ -> Zn^2+ + Cu}$, il y a \cle{transfert d'électrons} du zinc vers les ions cuivre(II). Le nombre d'électrons cédés par le donneur est \textbf{toujours égal} au nombre d'électrons captés par l'accepteur : les électrons ne se créent ni ne se perdent.
\end{aretenir}

\courssub{Oxydation et réduction : les deux faces d'un même transfert}

Nous pouvons maintenant donner un nom précis à chacun des deux phénomènes observés.

\begin{definition}[Oxydation]
L'\cle{oxydation} est la \textbf{perte d'un ou plusieurs électrons} par une espèce chimique (atome, ion ou molécule). On dit que l'espèce est \emph{oxydée}.

Exemple : le zinc est oxydé selon $\ce{Zn -> Zn^2+ + 2e-}$.
\end{definition}

\begin{definition}[Réduction]
La \cle{réduction} est le \textbf{gain d'un ou plusieurs électrons} par une espèce chimique. On dit que l'espèce est \emph{réduite}.

Exemple : l'ion cuivre(II) est réduit selon $\ce{Cu^2+ + 2e- -> Cu}$.
\end{definition}

\begin{attention}[Ne pas se fier au sens courant des mots]
Dans le langage courant, « réduire » signifie diminuer. En chimie, la réduction est un \textbf{gain} d'électrons ! Pour t'en souvenir, raisonne sur la charge : gagner des électrons (négatifs) fait \emph{diminuer} la charge algébrique de l'espèce : \ce{Cu^2+} devient \ce{Cu} (de $+2$ à $0$). La réduction est bien une « réduction de la charge ».
\end{attention}

\courssub{Oxydant et réducteur : qui fait quoi ?}

Une espèce ne se contente pas de subir : elle \emph{agit} sur l'autre. Le zinc, en cédant ses électrons, \textbf{provoque la réduction} de l'ion \ce{Cu^2+}. L'ion \ce{Cu^2+}, en captant des électrons, \textbf{provoque l'oxydation} du zinc. D'où deux nouveaux rôles :

\begin{definition}[Réducteur et oxydant]
\begin{itemize}
\item Un \cle{réducteur} est une espèce chimique capable de \textbf{céder des électrons}. En cédant des électrons, il est \textbf{oxydé}.
\item Un \cle{oxydant} est une espèce chimique capable de \textbf{capter des électrons}. En captant des électrons, il est \textbf{réduit}.
\end{itemize}
\end{definition}

Dans notre exemple :
\begin{itemize}
\item le zinc \ce{Zn} cède $2$ électrons : c'est le \textbf{réducteur}, et il est \textbf{oxydé} en \ce{Zn^2+} ;
\item l'ion \ce{Cu^2+} capte $2$ électrons : c'est l'\textbf{oxydant}, et il est \textbf{réduit} en \ce{Cu}.
\end{itemize}

\begin{methode}[Moyens mnémotechniques pour ne jamais se tromper]
Retiens l'une de ces deux astuces :
\begin{enumerate}
\item \textbf{Le RÉducteur RÉduit l'autre et se fait OXYder.} (Et symétriquement : l'OXydant OXYde l'autre et se fait RÉduire.)
\item \textbf{OIL RIG} (en anglais) : \emph{Oxidation Is Loss} (l'oxydation est une perte), \emph{Reduction Is Gain} (la réduction est un gain) — d'électrons.
\end{enumerate}
Vérifie toujours ta réponse en comptant les électrons dans les demi-équations : c'est le seul juge de paix infaillible.
\end{methode}

\begin{attention}[L'erreur classique du Probatoire]
Beaucoup d'élèves écrivent : « le zinc est oxydé, donc c'est un oxydant ». \textbf{C'est faux !} Le zinc est \emph{oxydé}, mais c'est un \emph{réducteur}. Le nom du rôle (oxydant/réducteur) désigne ce que l'espèce \textbf{fait subir à l'autre}, pas ce qu'elle subit elle-même. Oxydé $\neq$ oxydant : ce sont même deux notions toujours associées à des espèces \emph{opposées}. Une espèce oxydée est forcément un réducteur ; une espèce réduite est forcément un oxydant.
\end{attention}

\courssub{La réaction d'oxydoréduction : un couple indissociable}

\begin{definition}[Réaction d'oxydoréduction]
Une réaction d'\cle{oxydoréduction} est une réaction chimique au cours de laquelle il y a \textbf{transfert d'électrons} entre un réducteur (qui les cède) et un oxydant (qui les capte).
\end{definition}

\begin{propriete}[Simultanéité de l'oxydation et de la réduction]
L'oxydation et la réduction sont \textbf{simultanées et indissociables} : il ne peut pas y avoir de donneur d'électrons sans accepteur, ni d'accepteur sans donneur. Les électrons n'existent pas à l'état libre en solution aqueuse : chaque électron cédé par un réducteur est immédiatement capté par un oxydant. C'est pourquoi on parle d'\emph{oxydo}\emph{réduction} : un seul mot pour deux phénomènes inséparables.
\end{propriete}

\begin{popfigure}
\begin{center}
\begin{tikzpicture}[font=\small]
\draw[rounded corners=4pt, fill=popblue!12, draw=popblue, thick] (0,3.2) rectangle (3.6,4.6);
\node[align=center] at (1.8,3.9) {\textbf{OXYDATION}\\ perte d'électrons\\ $\ce{Zn -> Zn^2+ + 2e-}$};
\draw[rounded corners=4pt, fill=poporange!12, draw=poporange, thick] (5,3.2) rectangle (8.6,4.6);
\node[align=center] at (6.8,3.9) {\textbf{RÉDUCTION}\\ gain d'électrons\\ $\ce{Cu^2+ + 2e- -> Cu}$};
\draw[rounded corners=4pt, fill=poppink!12, draw=poppink, thick] (0,0.6) rectangle (3.6,2.0);
\node[align=center] at (1.8,1.3) {\textbf{RÉDUCTEUR}\\ donne des électrons\\ il est \emph{oxydé} (ex. \ce{Zn})};
\draw[rounded corners=4pt, fill=popgreen!12, draw=popgreen, thick] (5,0.6) rectangle (8.6,2.0);
\node[align=center] at (6.8,1.3) {\textbf{OXYDANT}\\ capte des électrons\\ il est \emph{réduit} (ex. \ce{Cu^2+})};
\draw[-{Stealth[length=2.5mm]}, very thick, poppurple] (1.8,3.1) -- (1.8,2.1);
\draw[-{Stealth[length=2.5mm]}, very thick, poppurple] (6.8,3.1) -- (6.8,2.1);
\draw[{Stealth[length=2.5mm]}-{Stealth[length=2.5mm]}, very thick, popdark] (3.7,3.9) -- (4.9,3.9) node[midway, above, font=\bfseries\small] {$e^-$};
\node[align=center, font=\itshape, popmuted] at (4.3,0) {Le réducteur subit l'oxydation ; l'oxydant subit la réduction.};
\end{tikzpicture}
\end{center}
\legende{Tableau de synthèse : les quatre notions croisées de l'oxydoréduction, illustrées sur le couple zinc/cuivre.}
\end{popfigure}

\courssub{Réinterprétation des expériences avec les acides dilués}

Revenons maintenant sur les expériences de la section 1 : l'action de l'acide chlorhydrique ou sulfurique dilué sur le zinc, le fer ou l'aluminium. Nous avions écrit, pour le fer :

\[ \ce{Fe(s) + 2H+(aq) -> Fe^2+(aq) + H2(g)} \]

Analysons cette réaction avec notre nouveau vocabulaire :
\begin{itemize}
\item le fer devient \ce{Fe^2+} : il perd $2$ électrons ($\ce{Fe -> Fe^2+ + 2e-}$), il est \textbf{oxydé} : c'est le \textbf{réducteur} ;
\item les ions \ce{H+} deviennent \ce{H2} : chaque ion gagne $1$ électron ($\ce{2H+ + 2e- -> H2}$), ils sont \textbf{réduits} : \ce{H+} est l'\textbf{oxydant}.
\end{itemize}

\begin{exemplebox}[L'attaque des métaux par les acides est une oxydoréduction]
Lorsqu'un clou en fer est plongé dans l'acide chlorhydrique dilué, le dégagement de dihydrogène et la dissolution du fer traduisent un transfert d'électrons du fer vers les ions \ce{H+}. C'est exactement le même type de phénomène que la réaction entre le zinc et les ions \ce{Cu^2+} : un métal (réducteur) cède des électrons à un accepteur (oxydant : \ce{H+} ou \ce{Cu^2+}). À l'inverse, le cuivre, l'argent et l'or ne réagissent pas avec \ce{H+} : ce sont des réducteurs trop faibles pour céder leurs électrons aux ions hydrogène — nous y reviendrons dans la section 5.
\end{exemplebox}

\courssub{Les couples oxydant/réducteur}

Observe les demi-équations rencontrées jusqu'ici : $\ce{Zn -> Zn^2+ + 2e-}$ et $\ce{Cu^2+ + 2e- -> Cu}$. Chacune met en regard deux formes d'un même élément : une forme \emph{oxydée} (qui peut capter des électrons) et une forme \emph{réduite} (qui peut en céder). On les regroupe en un \cle{couple oxydant/réducteur}, noté \textbf{Ox/Red}, et on écrit la demi-équation dans le sens conventionnel de la réduction :

\[ \ce{Ox + $n$ e- <=> Red} \]

où $n$ est le nombre d'électrons échangés. La double flèche rappelle que le transfert peut, selon les conditions, se faire dans un sens ou dans l'autre.

\begin{center}
\begin{tabularx}{\linewidth}{|l|c|X|}
\hline
\rowcolor{popblueL} \textbf{Couple Ox/Red} & \textbf{Demi-équation} & \textbf{Exemple de réaction} \\
\hline
\ce{Cu^2+}/\ce{Cu} & $\ce{Cu^2+ + 2e- <=> Cu}$ & \ce{Cu^2+} oxyde le zinc (dépôt de cuivre rougeâtre) \\
\hline
\ce{Zn^2+}/\ce{Zn} & $\ce{Zn^2+ + 2e- <=> Zn}$ & \ce{Zn} réduit \ce{Cu^2+} (le zinc se dissout) \\
\hline
\ce{H+}/\ce{H2} & $\ce{2H+ + 2e- <=> H2}$ & \ce{H+} des acides dilués oxyde le fer, le zinc, l'aluminium \\
\hline
\ce{Fe^2+}/\ce{Fe} & $\ce{Fe^2+ + 2e- <=> Fe}$ & \ce{Fe} réduit \ce{Cu^2+} ; \ce{Fe^2+} peut être réduit par le zinc \\
\hline
\end{tabularx}
\end{center}

\begin{aretenir}[Convention d'écriture]
Dans un couple Ox/Red, on écrit \textbf{toujours l'oxydant en premier} : \ce{Cu^2+}/\ce{Cu}, \ce{H+}/\ce{H2}. La demi-équation associée s'écrit dans le sens de la réduction : $\ce{Ox + $n$ e- <=> Red}$. Toute réaction d'oxydoréduction met en jeu \textbf{deux couples} : le réducteur de l'un cède ses électrons à l'oxydant de l'autre.
\end{aretenir}

\begin{exoresolu}[Identifier oxydant et réducteur dans une réaction]
Le trempage d'une lame de fer dans une solution de sulfate de cuivre(II) est une technique artisanale connue : la lame se recouvre de cuivre. L'équation de la réaction est :
\[ \ce{Fe(s) + Cu^2+(aq) -> Fe^2+(aq) + Cu(s)} \]
\begin{enumerate}
\item Écrire les deux demi-équations électroniques.
\item Identifier l'espèce oxydée, l'espèce réduite, l'oxydant et le réducteur.
\item Citer les deux couples mis en jeu.
\end{enumerate}
\tcblower
\textbf{Solution détaillée.}
\begin{enumerate}
\item Le fer passe de \ce{Fe} (charge $0$) à \ce{Fe^2+} (charge $+2$) : il perd $2$ électrons :
\[ \ce{Fe(s) -> Fe^2+(aq) + 2e-} \]
L'ion cuivre(II) passe de \ce{Cu^2+} à \ce{Cu} : il gagne $2$ électrons :
\[ \ce{Cu^2+(aq) + 2e- -> Cu(s)} \]
Le nombre d'électrons cédés égale le nombre d'électrons captés ($2$) : le transfert est équilibré, et la somme des deux demi-équations redonne bien l'équation-bilan.
\item \begin{itemize}
\item Espèce \textbf{oxydée} : le fer \ce{Fe} (il perd des électrons) ;
\item espèce \textbf{réduite} : l'ion \ce{Cu^2+} (il gagne des électrons) ;
\item \textbf{réducteur} : le fer \ce{Fe} (c'est lui qui cède les électrons — il est oxydé) ;
\item \textbf{oxydant} : l'ion \ce{Cu^2+} (c'est lui qui capte les électrons — il est réduit).
\end{itemize}
\item Les deux couples mis en jeu sont \ce{Cu^2+}/\ce{Cu} et \ce{Fe^2+}/\ce{Fe}. Le réducteur du couple \ce{Fe^2+}/\ce{Fe} a cédé ses électrons à l'oxydant du couple \ce{Cu^2+}/\ce{Cu}.
\end{enumerate}
\textbf{Vérification anti-piège :} le fer est oxydé \emph{et} réducteur ; dire « le fer est oxydé donc c'est un oxydant » serait l'erreur classique à ne jamais commettre.
\end{exoresolu}

\begin{savaistu}[Pourquoi dit-on « oxydation » alors qu'il n'y a pas d'oxygène ?]
Historiquement, au XVIII\textsuperscript{e} siècle, Lavoisier a baptisé « oxydation » les réactions de combinaison avec le \emph{oxygène} (comme la rouille du fer ou la combustion du charbon), et « réduction » la récupération d'un métal à partir de son oxyde — le minerai « réduit » donnant un métal dont la masse... diminuait ! Au début du XX\textsuperscript{e} siècle, les chimistes ont compris que le point commun de toutes ces réactions était le transfert d'électrons : combiner un métal à l'oxygène, c'est lui faire perdre des électrons au profit de l'oxygène. La définition moderne (perte et gain d'électrons) est donc beaucoup plus générale : elle s'applique même quand l'oxygène est totalement absent, comme dans $\ce{Zn + Cu^2+ -> Zn^2+ + Cu}$. Le nom est resté, le sens a grandi.
\end{savaistu}

\begin{exercice}[À toi de jouer]
L'aluminium réagit avec l'acide chlorhydrique dilué selon :
\[ \ce{2Al(s) + 6H+(aq) -> 2Al^3+(aq) + 3H2(g)} \]
\begin{enumerate}
\item Écris les deux demi-équations électroniques et vérifie que le nombre d'électrons échangés est le même.
\item Identifie l'oxydant, le réducteur, l'espèce oxydée et l'espèce réduite.
\item Cite les deux couples Ox/Red mis en jeu.
\end{enumerate}
\end{exercice}

\begin{corrige}[Exercice]
\begin{enumerate}
\item Oxydation de l'aluminium : $\ce{Al -> Al^3+ + 3e-}$, multipliée par $2$ : $\ce{2Al -> 2Al^3+ + 6e-}$. Réduction des ions hydrogène : $\ce{2H+ + 2e- -> H2}$, multipliée par $3$ : $\ce{6H+ + 6e- -> 3H2}$. Les deux demi-équations échangent bien $6$ électrons : le transfert est équilibré.
\item L'aluminium \ce{Al} cède des électrons : c'est le \textbf{réducteur}, il est \textbf{oxydé}. Les ions \ce{H+} captent des électrons : \ce{H+} est l'\textbf{oxydant}, il est \textbf{réduit} en \ce{H2}.
\item Couples mis en jeu : \ce{Al^3+}/\ce{Al} et \ce{H+}/\ce{H2}.
\end{enumerate}
\end{corrige}

En résumé, toute réaction d'oxydoréduction est un transfert d'électrons d'un réducteur vers un oxydant, mettant en jeu deux couples Ox/Red. Reste une question décisive : peut-on \emph{prévoir} si un transfert aura lieu, et dans quel sens ? Pourquoi le zinc réduit-il \ce{Cu^2+} alors que le cuivre ne réduit pas \ce{Zn^2+} ? C'est l'objet de la fin de cette leçon, où nous classerons les couples grâce à la notion de force des oxydants et des réducteurs.

\coursec{Demi-équations électroniques et équation-bilan}

Dans la section précédente, nous avons appris à reconnaître qui est l'\cle{oxydant} et qui est le \cle{réducteur} dans une réaction comme celle du zinc avec les ions cuivre(II) : le zinc \emph{perd} des électrons (oxydation), l'ion \ce{Cu^2+} les \emph{gagne} (réduction). Mais un transfert d'électrons, c'est un peu comme un paiement au marché de Mokolo : il y a un donneur et un receveur, et la somme donnée doit être \emph{exactement égale} à la somme reçue. Comment écrire proprement cette comptabilité électronique ? C'est toute l'ambition de cette section : séparer la réaction en deux \cle{demi-équations électroniques}, puis les recombiner pour obtenir l'\cle{équation-bilan} rigoureusement équilibrée.

\courssub{Qu'est-ce qu'une demi-équation électronique ?}

\textbf{L'intuition d'abord.}\par
Quand un atome de zinc plongé dans une solution de sulfate de cuivre(II) devient un ion \ce{Zn^2+}, il ne disparaît pas dans la nature : il \emph{lâche} deux électrons qui restent « disponibles » pour un receveur. On peut écrire cet événement seul, indépendamment du reste :

\[ \ce{Zn -> Zn^2+ + 2e-} \]

De son côté, l'ion \ce{Cu^2+} capte deux électrons pour se transformer en cuivre métallique :

\[ \ce{Cu^2+ + 2e- -> Cu} \]

Chacune de ces écritures ne décrit que la \emph{moitié} de l'histoire : c'est pourquoi on les appelle \textbf{demi-équations}. Aucune ne peut se produire seule en solution : les électrons libérés par l'une sont immédiatement consommés par l'autre. Il n'existe pas d'électrons « en liberté » dans un bécher !

\begin{definition}[Demi-équation électronique]
Une \cle{demi-équation électronique} est une écriture formelle qui traduit l'oxydation ou la réduction d'un couple oxydant/réducteur en faisant apparaître explicitement les électrons échangés. Elle s'écrit sous la forme générale :
\[ \text{Ox} + n\,\text{e}^- \rightleftharpoons \text{Red} \]
où $n$ est le nombre d'électrons échangés. La double flèche rappelle que l'écriture est \emph{formelle} : le sens réel du transfert dépend de l'espèce avec laquelle le couple est mis en présence.
\end{definition}

\textbf{Les demi-équations des couples étudiés dans cette leçon.}\par
Reprenons tous les couples rencontrés dans les sections précédentes et écrivons leur demi-équation :

\begin{center}
\begin{tabularx}{\linewidth}{|l|X|l|}
\hline
\rowcolor{popblueL} \textbf{Couple Ox/Red} & \textbf{Demi-équation électronique} & \textbf{Électrons} \\
\hline
\ce{Zn^2+}/\ce{Zn} & $\ce{Zn^2+} + 2\,\text{e}^- \rightleftharpoons \ce{Zn}$ & $n = 2$ \\
\hline
\ce{Cu^2+}/\ce{Cu} & $\ce{Cu^2+} + 2\,\text{e}^- \rightleftharpoons \ce{Cu}$ & $n = 2$ \\
\hline
\ce{Fe^2+}/\ce{Fe} & $\ce{Fe^2+} + 2\,\text{e}^- \rightleftharpoons \ce{Fe}$ & $n = 2$ \\
\hline
\ce{H+}/\ce{H2} & $2\,\ce{H+} + 2\,\text{e}^- \rightleftharpoons \ce{H2}$ & $n = 2$ \\
\hline
\ce{Al^3+}/\ce{Al} & $\ce{Al^3+} + 3\,\text{e}^- \rightleftharpoons \ce{Al}$ & $n = 3$ \\
\hline
\end{tabularx}
\end{center}

\begin{attention}[Le couple \ce{H+}/\ce{H2} : un piège classique]
Pour le couple de l'ion hydrogène, il faut écrire \textbf{deux} ions \ce{H+} car le dihydrogène \ce{H2} contient \textbf{deux} atomes d'hydrogène : $2\,\ce{H+} + 2\,\text{e}^- \rightleftharpoons \ce{H2}$. Écrire $\ce{H+} + \text{e}^- \rightleftharpoons \ce{H2}$ est une double faute : ni les atomes ni les charges ne sont conservés. On équilibre \emph{d'abord} les atomes, \emph{ensuite} les charges avec les électrons.
\end{attention}

\courssub{Les deux lois de conservation dans une demi-équation}

Une demi-équation correcte obéit à deux lois inviolables, héritées de la physique la plus fondamentale :

\begin{propriete}[Conservations dans une demi-équation]
\begin{enumerate}
\item \textbf{Conservation de la matière} : chaque élément chimique apparaît en même nombre dans les deux membres.
\item \textbf{Conservation de la charge électrique} : la somme algébrique des charges (en comptant chaque électron comme $-1$) est identique dans les deux membres.
\end{enumerate}
\end{propriete}

\begin{exemplebox}[Vérification sur le couple \ce{Al^3+}/\ce{Al}]
Demi-équation : $\ce{Al^3+} + 3\,\text{e}^- \rightleftharpoons \ce{Al}$.
\begin{itemize}
\item \emph{Matière} : 1 atome Al à gauche, 1 atome Al à droite. \checkmark
\item \emph{Charge} : à gauche $(+3) + 3\times(-1) = 0$ ; à droite $0$. \checkmark
\end{itemize}
La demi-équation est correcte. Si l'on avait écrit $\ce{Al^3+} + 2\,\text{e}^- \rightleftharpoons \ce{Al}$, la charge à gauche vaudrait $+1$ contre $0$ à droite : la demi-équation serait \textbf{fausse}, même si les atomes sont équilibrés.
\end{exemplebox}

\begin{attention}[Deux équilibrages distincts à ne pas confondre]
Beaucoup d'élèves mélangent deux étapes différentes :
\begin{itemize}
\item l'\textbf{équilibrage des atomes} se fait avec des \emph{coefficients stœchiométriques} (ex. : $2\,\ce{H+}$ pour obtenir \ce{H2}) ;
\item l'\textbf{équilibrage des charges} se fait \emph{uniquement} avec des électrons $\text{e}^-$, jamais en modifiant les coefficients des espèces chimiques.
\end{itemize}
On ne « répare » jamais une charge déséquilibrée en ajoutant des atomes, ni un atome manquant en ajoutant des électrons !
\end{attention}

\courssub{Méthode : des demi-équations à l'équation-bilan}

Voici maintenant la grande question : comment passer des deux demi-équations à l'équation-bilan de la réaction observée au laboratoire ? Rappelons le principe comptable énoncé en introduction : \textbf{tous les électrons cédés par le réducteur doivent être captés par l'oxydant}. Le nombre d'électrons perdus doit donc être égal au nombre d'électrons gagnés.

\begin{methode}[Établir une équation-bilan d'oxydoréduction en 4 étapes]
\begin{enumerate}
\item \textbf{Identifier les deux couples} mis en jeu et repérer qui s'oxyde (le réducteur) et qui se réduit (l'oxydant).
\item \textbf{Écrire les deux demi-équations} dans le \emph{bon sens} : la réduction dans le sens Ox $+ n\,\text{e}^-$ $\to$ Red, l'oxydation dans le sens Red $\to$ Ox $+ n\,\text{e}^-$.
\item \textbf{Égaliser les électrons} : multiplier chaque demi-équation entière par un coefficient tel que le nombre d'électrons cédés égale le nombre d'électrons captés (on cherche le plus petit commun multiple).
\item \textbf{Additionner membre à membre} et simplifier : les électrons, présents en nombre égal des deux côtés, disparaissent. \textbf{Vérifier} enfin la conservation des atomes et de la charge globale.
\end{enumerate}
\end{methode}

\begin{popfigure}
\begin{center}
\begin{tikzpicture}[
  node distance=0.55cm,
  etape/.style={rectangle, rounded corners=4pt, draw=popblue, thick, fill=popblueL, text width=0.9\linewidth, align=left, inner sep=6pt, font=\small},
  num/.style={circle, draw=poporange, fill=poporange, text=white, font=\bfseries\small, inner sep=2pt, minimum size=0.55cm},
  fleche/.style={-{Stealth[length=3mm]}, thick, poporange}
]
\node[etape] (e1) {\textbf{Identifier les couples} : qui est l'oxydant ? qui est le réducteur ? (ex. : \ce{Cu^2+}/\ce{Cu} et \ce{Zn^2+}/\ce{Zn})};
\node[num, left=0.25cm of e1] {1};
\node[etape, below=of e1] (e2) {\textbf{Écrire les demi-équations dans le bon sens} : réduction de l'oxydant, oxydation du réducteur};
\node[num, left=0.25cm of e2] {2};
\node[etape, below=of e2] (e3) {\textbf{Multiplier pour égaliser les électrons} : chercher le PPCM des nombres d'électrons échangés};
\node[num, left=0.25cm of e3] {3};
\node[etape, below=of e3] (e4) {\textbf{Additionner, simplifier, vérifier} : les $\text{e}^-$ disparaissent ; contrôler atomes et charges};
\node[num, left=0.25cm of e4] {4};
\draw[fleche] (e1) -- (e2);
\draw[fleche] (e2) -- (e3);
\draw[fleche] (e3) -- (e4);
\end{tikzpicture}
\end{center}
\legende{Organigramme de la méthode en 4 étapes pour construire une équation-bilan d'oxydoréduction à partir des demi-équations électroniques.}
\end{popfigure}

\courssub{Exemple guidé 1 : le zinc et les ions cuivre(II)}

Reprenons l'expérience de la section 2 : une lame de zinc plongée dans une solution bleue de sulfate de cuivre(II) se couvre d'un dépôt rougeâtre de cuivre, et la solution se décolore progressivement.

\textbf{Étape 1 — Couples en jeu.} Le zinc métal s'oxyde (couple \ce{Zn^2+}/\ce{Zn}) ; l'ion \ce{Cu^2+} se réduit (couple \ce{Cu^2+}/\ce{Cu}).

\textbf{Étape 2 — Demi-équations dans le bon sens.}
\begin{align*}
\text{Oxydation : } & \ce{Zn -> Zn^2+ + 2e-} \\
\text{Réduction : } & \ce{Cu^2+ + 2e- -> Cu}
\end{align*}

\textbf{Étape 3 — Égalisation des électrons.} Ici, c'est le cas le plus simple : $2$ électrons sont cédés, $2$ électrons sont captés. Aucune multiplication n'est nécessaire (coefficients $\times 1$ des deux côtés).

\textbf{Étape 4 — Addition membre à membre.}
\[ \ce{Zn + Cu^2+ + 2e- -> Zn^2+ + 2e- + Cu} \]
Les $2\,\text{e}^-$ apparaissent des deux côtés : on les supprime. Il reste :
\[ \ce{Zn + Cu^2+ -> Zn^2+ + Cu} \]

\textbf{Vérification.} Atomes : 1 Zn et 1 Cu de chaque côté. Charge : $+2$ à gauche, $+2$ à droite. L'équation est correcte. Cette réaction est \emph{totale} et \emph{spontanée} à température ambiante ; elle est à la base du fonctionnement de la pile Daniell, que tu étudieras plus tard.

\courssub{Exemple guidé 2 : l'aluminium et les ions hydrogène}

Passons à un cas plus riche, où les deux demi-équations n'échangent pas le même nombre d'électrons. Une poudre d'aluminium versée dans l'acide chlorhydrique dilué produit un vif dégagement de dihydrogène (réaction exothermique, à réaliser avec précaution).

\textbf{Étape 1 — Couples en jeu.} L'aluminium s'oxyde (couple \ce{Al^3+}/\ce{Al}) ; l'ion \ce{H+} se réduit (couple \ce{H+}/\ce{H2}).

\textbf{Étape 2 — Demi-équations dans le bon sens.}
\begin{align*}
\text{Oxydation : } & \ce{Al -> Al^3+ + 3e-} \\
\text{Réduction : } & \ce{2H+ + 2e- -> H2}
\end{align*}

\textbf{Étape 3 — Égalisation des électrons.} L'aluminium libère $3$ électrons ; la réduction de \ce{H+} n'en consomme que $2$. Il faut trouver le \cle{plus petit commun multiple} de 3 et 2 : c'est \textbf{6}. On multiplie donc :
\begin{itemize}
\item la demi-équation de l'aluminium par $\times 2$ : $\ce{2Al -> 2Al^3+ + 6e-}$ ;
\item la demi-équation de l'ion hydrogène par $\times 3$ : $\ce{6H+ + 6e- -> 3H2}$.
\end{itemize}

\begin{exemplebox}[Pourquoi 6 électrons ?]
Chaque atome d'aluminium donne $3$ électrons, mais chaque molécule de \ce{H2} formée n'en réclame que $2$. Pour que rien ne se perde, il faut un nombre d'électrons divisible à la fois par 3 et par 2 : le plus petit est $6$. Ainsi $2$ atomes Al (qui donnent $2\times 3 = 6$ électrons) réagissent exactement avec $6$ ions \ce{H+} (qui forment $3$ molécules \ce{H2}). C'est la même logique que lorsqu'on achète des beignets vendus par 3 et des sachets de piment vendus par 2 : pour n'avoir aucun reste, il faut en prendre 6 de chaque !
\end{exemplebox}

\textbf{Étape 4 — Addition et simplification.}
\[ \ce{2Al + 6H+ + 6e- -> 2Al^3+ + 3H2 + 6e-} \]
Les $6\,\text{e}^-$ s'éliminent :
\[ \ce{2Al + 6H+ -> 2Al^3+ + 3H2} \]

\textbf{Vérification finale.} Atomes : 2 Al et 6 H de chaque côté. Charge : $6\times(+1) = +6$ à gauche ; $2\times(+3) = +6$ à droite. L'équation est équilibrée. Cette réaction, vive et exothermique, explique pourquoi on ne conserve jamais d'acide dans un récipient en aluminium.

\begin{popfigure}
\begin{center}
\begin{tikzpicture}[
  demi/.style={rectangle, rounded corners=3pt, draw=popdark, thick, fill=popcream, align=center, inner sep=6pt, font=\small},
  mult/.style={rectangle, draw=poporange, fill=poporangeL, rounded corners=3pt, font=\small\bfseries, inner sep=4pt},
  res/.style={rectangle, rounded corners=4pt, draw=popgreen, very thick, fill=popgreenL, align=center, inner sep=7pt, font=\small},
  fleche/.style={-{Stealth[length=3mm]}, thick, popdark}
]
\node[demi] (ox) at (0,0) {Oxydation\\ $\ce{Al -> Al^3+ + 3e-}$};
\node[demi] (red) at (0,-2.2) {Réduction\\ $\ce{2H+ + 2e- -> H2}$};
\node[mult, right=1.6cm of ox] (m1) {$\times\,2$};
\node[mult, right=1.6cm of red] (m2) {$\times\,3$};
\node[demi, right=1.6cm of m1] (ox2) {$\ce{2Al -> 2Al^3+ + 6e-}$};
\node[demi, right=1.6cm of m2] (red2) {$\ce{6H+ + 6e- -> 3H2}$};
\node[res, below=2.6cm of $(ox2)!0.5!(red2)$] (bil) {Équation-bilan : $\ce{2Al + 6H+ -> 2Al^3+ + 3H2}$\\ les $6\,\text{e}^-$ ont disparu};
\draw[fleche] (ox) -- (m1);
\draw[fleche] (m1) -- (ox2);
\draw[fleche] (red) -- (m2);
\draw[fleche] (m2) -- (red2);
\draw[fleche] (ox2.south) -- (bil.north -| ox2.south);
\draw[fleche] (red2.north) -- (bil.north -| red2.north);
\node[font=\footnotesize\itshape, popmuted] at ($(m1)!0.5!(m2) + (0,-0.05)$) {PPCM de 3 et 2 : 6 électrons};
\end{tikzpicture}
\end{center}
\legende{Combinaison des demi-équations des couples \ce{Al^3+}/\ce{Al} et \ce{H+}/\ce{H2} : les multiplicateurs $\times 2$ et $\times 3$ égalisent les électrons échangés, qui disparaissent de l'équation-bilan.}
\end{popfigure}

\begin{attention}[Jamais d'électrons dans l'équation-bilan finale !]
L'erreur la plus fréquente au Probatoire est de laisser des électrons $\text{e}^-$ dans l'équation-bilan, par exemple : $\ce{Zn + Cu^2+ -> Zn^2+ + Cu + 2e-}$. C'est \textbf{faux} : dans une solution, les électrons ne circulent pas librement ; ils passent directement du réducteur à l'oxydant. Si des électrons subsistent dans ton bilan, c'est le signe que l'étape 3 (égalisation) a été mal faite. Recommence en cherchant le PPCM des électrons échangés.
\end{attention}

\begin{exoresolu}[Action du fer sur une solution de sulfate de cuivre(II)]
Un clou en fer de masse $m = \SI{5,6}{g}$ est plongé dans une solution contenant des ions \ce{Cu^2+} en excès. On observe un dépôt de cuivre sur le clou. (1) Écrire les demi-équations électroniques. (2) Établir l'équation-bilan. (3) Calculer la quantité de cuivre déposé si tout le fer réagit. Donnée : $M(\text{Fe}) = \SI{56}{g.mol^{-1}}$.
\tcblower
\textbf{(1) Demi-équations.} Le fer est le réducteur, il s'oxyde ; l'ion \ce{Cu^2+} est l'oxydant, il se réduit :
\begin{align*}
\text{Oxydation : } & \ce{Fe -> Fe^2+ + 2e-} \\
\text{Réduction : } & \ce{Cu^2+ + 2e- -> Cu}
\end{align*}
\textbf{(2) Équation-bilan.} Les deux demi-équations échangent le même nombre d'électrons ($n = 2$) : on additionne directement, les électrons s'éliminent :
\[ \ce{Fe + Cu^2+ -> Fe^2+ + Cu} \]
Vérification : 1 Fe et 1 Cu de chaque côté ; charge $+2$ dans les deux membres.
\textbf{(3) Quantité de cuivre.} La quantité de fer consommé est :
\[ n(\text{Fe}) = \frac{m}{M} = \frac{5,6}{56} = \SI{0,10}{mol} \]
D'après l'équation-bilan, $1$ mol de Fe produit $1$ mol de Cu, donc $n(\text{Cu}) = \SI{0,10}{mol}$, soit une masse $m(\text{Cu}) = 0,10 \times 63,5 = \SI{6,35}{g}$ de cuivre déposé sur le clou.
\end{exoresolu}

\begin{aretenir}[L'essentiel de la section]
\begin{itemize}
\item Une \cle{demi-équation électronique} s'écrit Ox $+ n\,\text{e}^- \rightleftharpoons$ Red ; elle conserve la matière \emph{et} la charge.
\item Pour obtenir l'\cle{équation-bilan} : identifier les couples, écrire les demi-équations dans le bon sens, multiplier pour égaliser les électrons (PPCM), additionner et simplifier.
\item Les électrons ne doivent \textbf{jamais} apparaître dans l'équation-bilan finale.
\item Toujours terminer par la double vérification : conservation des atomes et conservation de la charge globale.
\end{itemize}
\end{aretenir}

\begin{exercice}[À toi de jouer]
Le zinc réagit avec l'acide chlorhydrique dilué en dégageant du dihydrogène. (1) Identifier les deux couples mis en jeu. (2) Écrire les deux demi-équations électroniques. (3) Établir l'équation-bilan et vérifier les conservations. (4) Même question pour l'action du fer sur une solution contenant des ions \ce{Zn^2+} : la réaction a-t-elle lieu ? Justifier à l'aide de ce que tu as observé dans les sections précédentes.
\end{exercice}

\begin{corrige}[Exercice]
(1) Couples : \ce{Zn^2+}/\ce{Zn} (le zinc s'oxyde) et \ce{H+}/\ce{H2} (l'ion \ce{H+} se réduit).
(2) Oxydation : $\ce{Zn -> Zn^2+ + 2e-}$ ; réduction : $\ce{2H+ + 2e- -> H2}$.
(3) Les électrons sont déjà en nombre égal ($n = 2$) : $\ce{Zn + 2H+ -> Zn^2+ + H2}$. Atomes : 1 Zn, 2 H de chaque côté ; charge $+2$ dans les deux membres. L'équation est correcte.
(4) Avec le fer et les ions \ce{Zn^2+}, aucune réaction n'est observée expérimentalement : le fer ne peut pas réduire les ions \ce{Zn^2+}. Écrire $\ce{Fe + Zn^2+ -> Fe^2+ + Zn}$ serait formellement équilibré, mais cette réaction ne se produit pas spontanément : c'est la réaction \emph{inverse} (action du zinc sur les ions \ce{Fe^2+}) qui a lieu. On verra dans la section suivante comment prévoir le sens des réactions grâce à la force relative des oxydants et des réducteurs.
\end{corrige}

\begin{savaistu}[Des demi-équations à la pile Daniell]
En 1836, le chimiste britannique John Frederic Daniell a eu l'idée géniale de \emph{séparer physiquement} les deux demi-équations $\ce{Zn -> Zn^2+ + 2e-}$ et $\ce{Cu^2+ + 2e- -> Cu}$ dans deux compartiments reliés par un fil : les électrons, obligés de passer par le circuit extérieur, y produisent un courant électrique utilisable ! C'est la pile Daniell, ancêtre de toutes nos batteries. La prochaine fois que tu recharges ton téléphone à Yaoundé ou à Douala, souviens-toi : tout part de deux demi-équations électroniques bien équilibrées.
\end{savaistu}

\coursec{Prévision des réactions et applications}

Dans la section précédente, tu as appris à écrire les demi-équations électroniques et à les combiner pour obtenir l'équation-bilan d'une oxydoréduction, comme l'attaque du zinc par l'acide chlorhydrique ou le dépôt de cuivre sur un clou de fer. Une question se pose alors naturellement : \emph{peut-on savoir à l'avance si une réaction aura lieu ?} Un bijoutier de Douala doit-il craindre de plonger un anneau d'or dans un acide ? Un plombier peut-il fixer un tuyau de cuivre sur une gouttière en zinc ? Cette dernière section te donne la clé : le \cle{classement des métaux par pouvoir réducteur}, qui permet de \emph{prévoir} les réactions au lieu de les subir.

\courssub{Classement qualitatif des métaux par pouvoir réducteur}

\textbf{Intuition.} Reprenons les expériences des sections 1 et 2. Nous avons observé que l'aluminium, le zinc et le fer sont attaqués par les acides chlorhydrique et sulfurique dilués avec dégagement de dihydrogène (l'aluminium parfois après un temps de latence dû à sa couche d'oxyde), alors que le cuivre, l'argent et l'or restent intacts. Nous avons aussi vu que le zinc dépose le cuivre à partir d'une solution de \ce{Cu^2+}, que le fer dépose aussi le cuivre, et que le zinc dépose le fer à partir de \ce{Fe^2+}, mais que les réactions inverses sont impossibles. Tout cela suggère que les métaux n'ont pas tous la même « générosité » vis-à-vis de leurs électrons : certains les cèdent très facilement, d'autres les gardent jalousement.

\begin{definition}[Pouvoir réducteur d'un métal]
Le \cle{pouvoir réducteur} d'un métal traduit sa tendance thermodynamique à céder des électrons, c'est-à-dire à s'oxyder. Plus un métal cède facilement ses électrons (plus son potentiel standard d'oxydation est élevé, ou son potentiel de réduction est bas), plus il est \emph{réducteur}.
\end{definition}

Les expériences réalisées en classe, corroborées par les potentiels standards, permettent d'établir le classement qualitatif suivant, par pouvoir réducteur \emph{décroissant} (du plus fort réducteur au plus faible) :

\begin{popfigure}
\begin{center}
\begin{tikzpicture}[scale=1, every node/.style={font=\large}]
\draw[-{Stealth[length=3mm]}, line width=1.2pt, popblue] (0,0) -- (12.5,0);
\node[below, popblue, font=\small\itshape] at (6.25,-0.15) {pouvoir réducteur décroissant $\longrightarrow$};
\foreach \x/\m/\c in {0.8/{Al}/popink, 2.6/{Zn}/popink, 4.4/{Fe}/popink, 6.4/{H}/popmuted, 8.2/{Cu}/popgreen, 9.6/{Ag}/popgreen, 11.0/{Au}/popgreen}{
  \node[circle, draw=\c, fill=\c!15, minimum size=0.85cm, font=\bfseries] at (\x,0.55) {\m};
}
\node[above, align=center, font=\small, popink] at (2.6,1.15) {réduisent \ce{H+}\\ (attaqués par les acides dilués)};
\node[above, align=center, font=\small, popgreen] at (9.6,1.15) {ne réduisent pas \ce{H+}\\ (métaux « nobles »)};
\draw[dashed, popmuted] (5.4,-0.4) -- (5.4,1.5);
\node[font=\tiny, popmuted, rotate=90] at (5.4,0.55) {Hydrogène};
\end{tikzpicture}
\end{center}
\legende{Classement qualitatif des métaux étudiés par pouvoir réducteur décroissant. L'hydrogène sert de frontière : à sa gauche, les métaux sont attaqués par les acides dilués non oxydants ; à sa droite, ils ne le sont pas. L'aluminium est le plus fort réducteur de la série, mais sa réactivité est souvent masquée par la passivation.}
\end{popfigure}

\begin{aretenir}[Classement à connaître par cœur]
\[\text{Al} > \text{Zn} > \text{Fe} > (\text{H}) > \text{Cu} > \text{Ag} > \text{Au}\]
\begin{itemize}
\item \ce{Al}, \ce{Zn}, \ce{Fe} \textbf{réduisent} les ions \ce{H+} : ils sont attaqués par \ce{HCl} et \ce{H2SO4} dilués, avec dégagement de \ce{H2} (pour l'Al, réaction souvent lente au démarrage à cause de la couche d'oxyde).
\item \ce{Cu}, \ce{Ag}, \ce{Au} \textbf{ne réduisent pas} \ce{H+} : ils résistent aux acides dilués non oxydants. C'est pourquoi l'or et l'argent servent en bijouterie : ils ne s'altèrent pas au contact de la sueur (acide) ou de l'air.
\end{itemize}
\end{aretenir}

\begin{attention}[Nuance cruciale sur l'aluminium]
L'aluminium est le \emph{meilleur} réducteur de cette liste (potentiel standard \ce{Al^3+/Al} = \SI{-1,66}{V}), pourtant une casserole en aluminium résiste à l'air, à l'eau et même aux acides dilués froids ! La raison : il se recouvre spontanément d'une fine couche d'oxyde \ce{Al2O3}, imperméable et adhérente, qui le \emph{passive}. Cette couche protectrice masque sa réactivité thermodynamique réelle. Ne conclus donc jamais, à partir d'un objet en aluminium brillant, que l'aluminium est un métal noble. Si l'on gratte la couche d'oxyde (ou si l'on utilise de l'aluminium amalgamé), la réaction avec les acides est violente.
\end{attention}

\courssub{Règle de prévision des réactions métal / ion métallique}

\textbf{Intuition.} Imagine deux métaux en compétition pour des électrons. Le plus réducteur des deux, le plus « généreux », cède ses électrons ; l'autre, sous forme d'ion, les reçoit et se transforme en métal. La réaction ne peut se produire spontanément que dans ce sens-là, jamais dans l'autre.

\begin{propriete}[Règle de prévision (Métal + Ion métallique)]
Un métal \textbf{réduit les ions d'un métal moins réducteur que lui}. Autrement dit, la réaction d'oxydoréduction a lieu si le métal introduit (solide) se trouve \emph{au-dessus} (à gauche) du métal dont l'ion est présent en solution, dans le classement par pouvoir réducteur.
\end{propriete}

\begin{methode}[Prévoir une réaction métal / ion en 4 étapes]
\begin{enumerate}
\item \textbf{Repérer} le métal solide introduit ($M_{\text{solide}}$) et l'ion métallique présent en solution ($M'^{n+}$).
\item \textbf{Comparer} leurs positions dans le classement par pouvoir réducteur : $M_{\text{solide}}$ est-il plus réducteur que $M'$ ?
\item \textbf{Conclure} : si $M_{\text{solide}}$ est plus réducteur que $M'$ (il est à gauche dans le classement), la réaction a lieu ; sinon, il ne se passe rien (pas de réaction spontanée).
\item \textbf{Écrire} les deux demi-équations électroniques (oxydation du métal, réduction de l'ion), puis l'équation-bilan en simplifiant les électrons.
\end{enumerate}
\end{methode}

\begin{exemplebox}[Trois prévisions classiques]
\begin{itemize}
\item \textbf{Zinc dans une solution de \ce{Cu^2+}} : Zn est plus réducteur que Cu ($Zn > Cu$) $\Rightarrow$ réaction :
\[ \ce{Zn(s) + Cu^2+(aq) -> Zn^2+(aq) + Cu(s)} \]
(Demi-équations : \ce{Zn -> Zn^2+ + 2e-} ; \ce{Cu^2+ + 2e- -> Cu}). On observe un dépôt rougeâtre de cuivre et la décoloration de la solution bleue.
\item \textbf{Fer dans une solution de \ce{Cu^2+}} : Fe est plus réducteur que Cu ($Fe > Cu$) $\Rightarrow$ réaction :
\[ \ce{Fe(s) + Cu^2+(aq) -> Fe^2+(aq) + Cu(s)} \]
\item \textbf{Zinc dans une solution de \ce{Fe^2+}} : Zn est plus réducteur que Fe ($Zn > Fe$) $\Rightarrow$ réaction :
\[ \ce{Zn(s) + Fe^2+(aq) -> Zn^2+(aq) + Fe(s)} \]
\end{itemize}
\end{exemplebox}

\begin{attention}[L'erreur classique du Probatoire]
Beaucoup d'élèves écrivent : \ce{Cu + Zn^2+ -> Cu^2+ + Zn}. \textbf{Cette réaction n'a pas lieu !} Le cuivre est \emph{moins} réducteur que le zinc ($Cu < Zn$) : il ne peut pas céder ses électrons aux ions \ce{Zn^2+}. Une lame de cuivre plongée dans une solution de sulfate de zinc reste rigoureusement intacte. Retiens : \emph{c'est toujours le métal le plus réducteur (le plus à gauche) qui donne ses électrons}. Vérifie systématiquement le classement avant d'écrire une équation.
\end{attention}

\courssub{Tableau récapitulatif des réactions possibles}

Le tableau ci-dessous synthétise toutes les situations que tu peux rencontrer au Probatoire : en ligne, les métaux solides ; en colonne, les ions en solution (dont \ce{H+} des acides dilués).

\begin{popfigure}
\begin{center}
\begin{tikzpicture}[scale=0.95]
% Grille : 5 lignes (entête + 4 métaux), 6 colonnes (Métal + 5 ions)
\foreach \i in {0,...,4} \draw[popline] (0,-\i*1.1) -- (13.1,-\i*1.1);
\foreach \j in {0,...,5} \draw[popline] ({1.6+\j*2.3},0.6) -- ({1.6+\j*2.3},-4.4);
\draw[popline] (0,0.6) -- (0,-4.4);
% En-têtes colonnes
\node[font=\small\bfseries, popdark] at (0.8,-0.275) {Métal solide};
\node[font=\small\bfseries, popblue] at (2.75,0.325) {\ce{Cu^2+}};
\node[font=\small\bfseries, popblue] at (5.05,0.325) {\ce{Zn^2+}};
\node[font=\small\bfseries, popblue] at (7.35,0.325) {\ce{Fe^2+}};
\node[font=\small\bfseries, popblue] at (9.65,0.325) {\ce{Al^3+}};
\node[font=\small\bfseries, popblue] at (11.95,0.325) {\ce{H+}};
% Lignes métaux (ordre décroissant pouvoir réducteur)
\node[font=\small\bfseries, popdark] at (0.8,-0.825) {\ce{Al}};
\node[font=\small\bfseries, popdark] at (0.8,-1.925) {\ce{Zn}};
\node[font=\small\bfseries, popdark] at (0.8,-3.025) {\ce{Fe}};
\node[font=\small\bfseries, popdark] at (0.8,-4.125) {\ce{Cu}};
% Cases : oui (vert) / non (rouge) / lent (orange)
% Al (ligne 1, y=-0.825) : réduit tout ce qui est à droite (Zn, Fe, H). Pas Zn2+, Al3+.
\node[font=\small, popgreen] at (2.75,-0.825) {\textbf{oui}};
\node[font=\small, popink]  at (5.05,-0.825) {non};
\node[font=\small, popgreen] at (7.35,-0.825) {\textbf{oui}};
\node[font=\small, popink]  at (9.65,-0.825) {non};
\node[font=\small, poporange] at (11.95,-0.825) {\textbf{oui}*};
% Zn (ligne 2, y=-1.925) : réduit Fe, H. Pas Cu, Zn, Al.
\node[font=\small, popgreen] at (2.75,-1.925) {\textbf{oui}};
\node[font=\small, popink]  at (5.05,-1.925) {non};
\node[font=\small, popgreen] at (7.35,-1.925) {\textbf{oui}};
\node[font=\small, popink]  at (9.65,-1.925) {non};
\node[font=\small, popgreen] at (11.95,-1.925) {\textbf{oui} (\ce{H2})};
% Fe (ligne 3, y=-3.025) : réduit Cu, H. Pas Zn, Fe, Al.
\node[font=\small, popgreen] at (2.75,-3.025) {\textbf{oui}};
\node[font=\small, popink]  at (5.05,-3.025) {non};
\node[font=\small, popink]  at (7.35,-3.025) {non};
\node[font=\small, popink]  at (9.65,-3.025) {non};
\node[font=\small, popgreen] at (11.95,-3.025) {\textbf{oui} (\ce{H2})};
% Cu (ligne 4, y=-4.125) : ne réduit rien.
\node[font=\small, popink]  at (2.75,-4.125) {non};
\node[font=\small, popink]  at (5.05,-4.125) {non};
\node[font=\small, popink]  at (7.35,-4.125) {non};
\node[font=\small, popink]  at (9.65,-4.125) {non};
\node[font=\small, popink]  at (11.95,-4.125) {non};
% Légende
\node[font=\small\itshape, popmuted, anchor=west] at (0,-5.0) {\textbf{oui} : réaction spontanée observée ; \textbf{oui}* : réaction thermodynamique mais lente au démarrage (passivation) ; \textbf{non} : aucune réaction.};
\node[font=\small\itshape, popmuted, anchor=west] at (0,-5.4) {Un métal ne réduit jamais ses propres ions (diagonale « non »).};
\end{tikzpicture}
\end{center}
\legende{Prévision des réactions entre un métal (ligne) et un ion en solution (colonne). On lit par exemple : le zinc réduit \ce{Cu^2+}, \ce{Fe^2+} et \ce{H+}, mais pas \ce{Zn^2+} ni \ce{Al^3+}. L'aluminium réduit \ce{Cu^2+}, \ce{Fe^2+} et \ce{H+} (souvent après grattage de la couche d'oxyde).}
\end{popfigure}

Remarque la logique implacable du tableau : un métal ne réagit qu'avec les ions situés \emph{à sa droite} dans le classement (métaux moins réducteurs, ou \ce{H+} s'il est à gauche de H). La diagonale « non » correspond au cas trivial où un métal est en présence de ses propres ions.

\courssub{Application 1 : identification d'un métal inconnu}

\textbf{Situation.} Un élève du lycée de Ngaoundéré trouve dans le laboratoire une lame métallique grise, sans étiquette. Est-ce de l'aluminium, du zinc, du fer ou du cuivre ? Grâce au classement des pouvoirs réducteurs, trois tests simples suffisent.

\begin{methode}[Identifier un métal par tests successifs]
\begin{enumerate}
\item \textbf{Test à l'acide chlorhydrique dilué (ou \ce{H2SO4} dilué)} :
\begin{itemize}
\item Si \textbf{aucun dégagement gazeux} n'a lieu (même après quelques minutes) : le métal est \emph{noble} (\ce{Cu}, \ce{Ag} ou \ce{Au}). La couleur achève l'identification (rougeâtre \ce{Cu}, blanc brillant \ce{Ag}, jaune \ce{Au}).
\item Si un \textbf{dégagement de \ce{H2}} se produit (bulles, détonation à la flamme = « aboiement ») : le métal est \ce{Al}, \ce{Zn} ou \ce{Fe}. Passer au test 2.
\end{itemize}
\item \textbf{Test avec une solution de \ce{Fe^2+} (ex. \ce{FeSO4}, verdâtre)} : plonger un échantillon frais (gratté si aspect terne).
\begin{itemize}
\item \textbf{Réaction rapide} (dépôt grisâtre de fer, solution se décolore) : c'est du \textbf{zinc} ($\ce{Zn} > \ce{Fe}$ / \ce{Zn + Fe^2+ -> Zn^2+ + Fe}).
\item \textbf{Aucune réaction visible} (même après attente) : c'est du \textbf{fer} (\ce{Fe} ne réduit pas \ce{Fe^2+}) \emph{ou} de l'\textbf{aluminium passivé}. Passer au test 3.
\end{itemize}
\item \textbf{Test avec une solution de \ce{Cu^2+} (ex. \ce{CuSO4}, bleue)} : plonger un échantillon frais (gratté).
\begin{itemize}
\item \textbf{Réaction rapide} (dépôt rouge de cuivre) : si le test 2 était négatif, c'est du \textbf{fer} (\ce{Fe + Cu^2+ -> Fe^2+ + Cu}).
\item \textbf{Réaction lente / après grattage} (dépôt rouge, solution se décolore lentement) : c'est de l'\textbf{aluminium} (\ce{2Al + 3Cu^2+ -> 2Al^3+ + 3Cu}, mais passivé).
\item \textbf{Réaction rapide} aux deux tests (2 et 3) : c'est du \textbf{zinc}.
\end{itemize}
\item \textbf{Conclure} en écrivant les équations-bilans des réactions observées.
\end{enumerate}
\end{methode}

\begin{exoresolu}[Le clou de fer et la solution de sulfate de cuivre]
Un clou en fer de masse $m = \SI{5,6}{g}$ est plongé dans un volume important (excès) d'une solution de sulfate de cuivre(II) \ce{CuSO4}. Après réaction complète du clou, on récupère le dépôt de cuivre, on le lave, on le sèche et on le pèse. Quelle masse de cuivre obtient-on ? On donne $M(\ce{Fe}) = \SI{56}{g.mol^{-1}}$ et $M(\ce{Cu}) = \SI{63,5}{g.mol^{-1}}$.
\tcblower
\textbf{1. Prévision.} Le fer est plus réducteur que le cuivre ($Fe > Cu$) : la réaction a lieu.
\[ \ce{Fe(s) + Cu^2+(aq) -> Fe^2+(aq) + Cu(s)} \]
\textbf{2. Quantité de matière de fer.} $n(\ce{Fe}) = \dfrac{m}{M(\ce{Fe})} = \dfrac{5,6}{56} = \SI{0,10}{mol}$.
\textbf{3. Tableau d'avancement} (les ions \ce{Cu^2+} sont en excès, le fer est le réactif limitant) :
\begin{center}
\begin{tabular}{lcccc}
\toprule
 & \ce{Fe(s)} & \ce{Cu^2+(aq)} & \ce{Fe^2+(aq)} & \ce{Cu(s)} \\
\midrule
État initial (mol) & 0,10 & excès & 0 & 0 \\
État final (mol) & 0 & excès & 0,10 & 0,10 \\
\bottomrule
\end{tabular}
\end{center}
D'après la stœchiométrie (coefficients 1:1), $n(\ce{Cu})_{\text{formé}} = n(\ce{Fe})_{\text{consommé}} = \SI{0,10}{mol}$.
\textbf{4. Masse de cuivre déposé :}
\[ m(\ce{Cu}) = n \times M(\ce{Cu}) = 0,10 \times 63,5 = \SI{6,35}{g} \]
\textbf{Conclusion.} Le clou se recouvre de \SI{6,35}{g} de cuivre rougeâtre. Remarque que la masse du dépôt ($\SI{6,35}{g}$) est supérieure à la masse de fer disparu ($\SI{5,6}{g}$) car $M(\ce{Cu}) > M(\ce{Fe})$ : c'est un bon moyen de vérifier la cohérence de tes résultats.
\end{exoresolu}

\courssub{Application 2 : la protection des métaux — la galvanisation}

\textbf{Intuition.} Si un métal plus réducteur donne volontiers ses électrons, on peut en profiter : placé au contact d'un métal moins réducteur, il s'oxydera \emph{à sa place}, comme un garde du corps qui prend les coups destinés à son protégé.

\begin{propriete}[Protection cathodique par un métal plus réducteur (Galvanisation)]
Un métal $M_1$ (ex. fer) peut être protégé de la corrosion par un métal $M_2$ \emph{plus réducteur} (ex. zinc) placé à son contact : c'est $M_2$ qui s'oxyde (il est « sacrifié »), tant qu'il en reste. C'est le principe de la \cle{galvanisation} : recouvrir le fer (ou l'acier) d'une couche de zinc.
\end{propriete}

\begin{popfigure}
\begin{center}
\begin{tikzpicture}[scale=1]
% Tôle : couche de fer
\fill[popmuted!40] (0,0) rectangle (10,1.2);
\draw[popdark, thick] (0,0) rectangle (10,1.2);
\node[font=\bfseries, popdark] at (5,0.6) {Fer (\ce{Fe}) — la tôle à protéger};
% Couche de zinc
\fill[popblue!45] (0,1.2) rectangle (10,1.75);
\draw[popblue, thick] (0,1.2) rectangle (10,1.75);
\node[font=\bfseries, popblue] at (5,1.475) {Zinc (\ce{Zn}) — couche protectrice};
% Rayure
\fill[popmuted!40] (6.6,1.2) rectangle (7.1,1.75);
\draw[popink, thick, dashed] (6.6,1.2) -- (6.6,1.75) (7.1,1.2) -- (7.1,1.75);
\node[font=\small, popink, align=center] at (6.85,2.35) {rayure :\\ le fer est exposé};
\draw[-{Stealth}, popink] (6.85,2.05) -- (6.85,1.8);
% Oxydation du zinc
\node[font=\small, popblue, align=center] at (2.6,2.6) {Le zinc s'oxyde (anode) :\\ \ce{Zn -> Zn^2+ + 2e-}};
\draw[-{Stealth}, popblue] (2.6,2.15) -- (2.6,1.8);
% Protection du fer
\node[font=\small, popgreen, align=center] at (9.0,-0.75) {Le fer est protégé (cathode) :\\ il reçoit les $e^-$, ne s'oxyde pas};
\draw[-{Stealth}, popgreen] (8.6,-0.35) -- (8.2,-0.05);
% Flux d'électrons
\draw[-{Stealth}, poporange, thick] (3.5,1.475) .. controls (5,2.0) and (6,2.0) .. (6.85,1.6);
\node[font=\small, poporange] at (5.1,2.15) {$e^-$};
% Ions Zn2+ partant
\draw[-{Stealth}, popblue, dashed] (2.6,1.475) -- (2.6,2.0);
\node[font=\small, popblue] at (2.6,2.15) {\ce{Zn^2+}};
\end{tikzpicture}
\end{center}
\legende{Principe de la tôle galvanisée : la couche de zinc barre le contact avec l'air et l'eau (protection barrière) ; même si elle est rayée, le zinc, plus réducteur, s'oxyde à la place du fer (protection cathodique / sacrificielle) et continue de le protéger.}
\end{popfigure}

\begin{savaistu}[La galvanisation au Cameroun et dans le monde]
Les tôles ondulées qui couvrent les toits de Yaoundé, Douala, Bafoussam ou Garoua sont des tôles \emph{galvanisées} : de l'acier recouvert d'une fine couche de zinc (par trempage à chaud ou électrolyse). Sans cette protection, le fer rouillerait en quelques saisons des pluies. Le zinc joue un double rôle : il forme d'abord une barrière physique, puis, si la tôle est rayée ou percée, il continue de protéger le fer en s'oxydant à sa place (anode sacrificielle). Le même principe protège les coques en acier des navires (blocs de zinc appelés « anodes sacrificielles » fixés sur la coque), les pipelines enterrés (oléoducs, gazoducs) et les pylônes électriques haute tension. À l'inverse, les métaux précieux n'ont besoin d'aucune protection : l'or ne réagissant ni avec les acides dilués ni avec la plupart des ions métalliques, un bijoutier du marché central peut tester l'authenticité d'un bijou avec de l'acide nitrique dilué (eau forte) — si le métal est attaqué, ce n'est pas de l'or pur (l'or ne réagit qu'avec l'eau régale, mélange \ce{HCl} + \ce{HNO3}) !
\end{savaistu}

\begin{attention}[Galvanisation et contact entre métaux différents (Corrosion galvanique)]
Ne jamais assembler directement, en milieu humide (eau de mer, eau de pluie, humidité), un métal noble et un métal réducteur sans isolation : c'est le métal réducteur (l'anode) qui sera rongé par corrosion galvanique. Par exemple, fixer une gouttière en \ce{Cu} sur une charpente en \ce{Fe} (ou utiliser des vis en acier sur du cuivre) accélère dramatiquement la corrosion du fer. À l'inverse, des vis en \ce{Zn} (galvanisées) sur une tôle en \ce{Fe} protègent la tôle. Avant d'assembler deux métaux en milieu humide, consulte toujours le classement par pouvoir réducteur : le plus réducteur sera l'anode sacrifiée.
\end{attention}

\courssub{Complément Probatoire : la réaction naturelle (hors strict programme mais très valorisé)}

Ce complément, très apprécié des correcteurs au Probatoire, généralise la règle de prévision à \emph{tous} les couples oxydant/réducteur, pas seulement aux métaux.

\begin{propriete}[Réaction naturelle]
Lorsque plusieurs oxydants et plusieurs réducteurs sont présents en solution, la réaction qui se produit spontanément en premier — la \cle{réaction naturelle} — est celle entre le \textbf{réducteur le plus fort} et l'\textbf{oxydant le plus fort} présents. Les réactions secondaires éventuelles suivent le même ordre de priorité (deuxième réducteur avec premier oxydant, etc.) tant qu'il reste des réactifs.
\end{propriete}

\begin{exemplebox}[Justifier qu'une réaction est naturelle]
On plonge une lame de zinc dans une solution contenant \emph{simultanément} des ions \ce{Cu^2+} et \ce{Fe^2+}.
\begin{itemize}
\item \textbf{Réducteurs présents :} \ce{Zn(s)} (le seul métal) — c'est le réducteur le plus fort disponible.
\item \textbf{Oxydants présents :} \ce{Cu^2+} et \ce{Fe^2+}. Or \ce{Cu} est \emph{moins} réducteur que \ce{Fe} ($Cu < Fe$), donc \ce{Cu^2+} est un \emph{meilleur} oxydant que \ce{Fe^2+} (il a plus tendance à capter des électrons).
\item La réaction naturelle est donc \textbf{d'abord} : \ce{Zn(s) + Cu^2+(aq) -> Zn^2+(aq) + Cu(s)}.
\item Ce n'est que lorsque \ce{Cu^2+} est \emph{épuisé} que le zinc (s'il en reste) attaque \ce{Fe^2+} : \ce{Zn(s) + Fe^2+(aq) -> Zn^2+(aq) + Fe(s)}.
\end{itemize}
Justifier qu'une réaction est « naturelle » revient donc à montrer qu'elle met en jeu le meilleur donneur d'électrons (réducteur le plus fort) et le meilleur capteur d'électrons (oxydant le plus fort) du milieu.
\end{exemplebox}

\courssub{Lien avec la vie quotidienne et bilan de la leçon}

\textbf{Autour de toi.} La rouille qui attaque un portail en fer non peint, la patine verte (\ce{Cu2(OH)2CO3}) des toitures en cuivre des bâtiments administratifs, le test à l'acide du bijoutier du marché Mokolo, les batteries de téléphone et les piles alcalines (qui ne sont que des oxydoréductions contrôlées produisant un courant électrique) : toutes ces situations relèvent du transfert d'électrons que tu viens d'étudier. Savoir classer les métaux par pouvoir réducteur, c'est savoir prévoir la corrosion, choisir les bons matériaux pour la construction, identifier un métal inconnu au laboratoire et comprendre le fonctionnement d'une pile.

\begin{aretenir}[Bilan de la leçon — À maîtriser pour le Probatoire]
\begin{itemize}
\item Les métaux se classent par pouvoir réducteur décroissant (thermodynamique) : $\text{Al} > \text{Zn} > \text{Fe} > (\text{H}) > \text{Cu} > \text{Ag} > \text{Au}$.
\item Les métaux placés \textbf{avant H} (\ce{Al}, \ce{Zn}, \ce{Fe}) réduisent \ce{H+} des acides dilués (\ce{HCl}, \ce{H2SO4}) avec dégagement de \ce{H2} (attention : \ce{Al} passivé = réaction lente/déclenchée par grattage) ; ceux placés \textbf{après H} (\ce{Cu}, \ce{Ag}, \ce{Au}) ne réagissent pas avec ces acides.
\item Un métal réduit les ions d'un métal \emph{moins réducteur} que lui (plus à droite dans le classement) : c'est la règle de prévision universelle pour les couples métalliques.
\item Toutes ces réactions sont des \cle{oxydoréductions} : le réducteur (le métal solide) \emph{cède} des électrons (demi-équation d'oxydation : $M \to M^{n+} + n\,e^-$), l'oxydant (l'ion \ce{H+} ou l'ion métallique $M'^{n+}$) les \emph{gagne} (demi-équation de réduction : $M'^{n+} + n\,e^- \to M'$). L'équation-bilan s'obtient en combinant les deux demi-équations de manière à éliminer les électrons.
\item Applications concrètes : identification d'un métal inconnu (tests \ce{HCl}, \ce{Fe^2+}, \ce{Cu^2+}), protection cathodique par galvanisation (zinc sur fer), prévision de la corrosion galvanique (contact métaux dissemblables), compréhension du sens spontané des réactions (réaction naturelle).
\end{itemize}
\end{aretenir}

\begin{exercice}[À toi de jouer — Type Probatoire]
On dispose de quatre lames métalliques : aluminium (grattée), zinc, fer, cuivre, et de quatre solutions : sulfate de cuivre \ce{CuSO4} (bleue), sulfate de zinc \ce{ZnSO4} (incolore), sulfate de fer(II) \ce{FeSO4} (verdâtre), acide chlorhydrique dilué \ce{HCl}.
\begin{enumerate}
\item Prévoir, en justifiant par le classement, toutes les réactions possibles entre ces métaux et ces solutions. Écrire les équations-bilans correspondantes.
\item Une lame inconnue, de couleur grise, est attaquée par \ce{HCl} dilué avec dégagement d'un gaz qui détonne à la flamme. Elle ne réagit \textbf{pas} avec la solution de \ce{FeSO4} (même après grattage et attente). En revanche, elle se recouvre rapidement d'un dépôt rougeâtre dans la solution de \ce{CuSO4}. Quel est ce métal ? Justifier.
\item Calculer la masse de zinc consommée lorsque \SI{3,27}{g} de zinc réagissent avec un excès d'acide chlorhydrique, et le volume de dihydrogène dégagé dans les conditions où $V_m = \SI{24}{L.mol^{-1}}$. On donne $M(\ce{Zn}) = \SI{65,4}{g.mol^{-1}}$.
\end{enumerate}
\end{exercice}

\begin{corrige}[Exercice]
\textbf{1.} D'après le classement $\text{Al} > \text{Zn} > \text{Fe} > (\text{H}) > \text{Cu}$ :
\begin{itemize}
\item \textbf{Aluminium (gratté)} réduit \ce{Cu^2+}, \ce{Fe^2+}, \ce{H+} : 
$\ce{2Al + 3Cu^2+ -> 2Al^3+ + 3Cu}$ ; $\ce{2Al + 3Fe^2+ -> 2Al^3+ + 3Fe}$ ; $\ce{2Al + 6H+ -> 2Al^3+ + 3H2}$.
Ne réduit pas \ce{Zn^2+}, \ce{Al^3+}.
\item \textbf{Zinc} réduit \ce{Cu^2+}, \ce{Fe^2+}, \ce{H+} : 
$\ce{Zn + Cu^2+ -> Zn^2+ + Cu}$ ; $\ce{Zn + Fe^2+ -> Zn^2+ + Fe}$ ; $\ce{Zn + 2H+ -> Zn^2+ + H2}$.
Ne réduit pas \ce{Zn^2+}, \ce{Al^3+}.
\item \textbf{Fer} réduit \ce{Cu^2+}, \ce{H+} : 
$\ce{Fe + Cu^2+ -> Fe^2+ + Cu}$ ; $\ce{Fe + 2H+ -> Fe^2+ + H2}$.
Ne réduit pas \ce{Zn^2+}, \ce{Fe^2+}, \ce{Al^3+}.
\item \textbf{Cuivre} ne réduit aucun des ions proposés (\ce{Cu^2+}, \ce{Zn^2+}, \ce{Fe^2+}, \ce{Al^3+}, \ce{H+}). Aucune réaction.
\end{itemize}

\textbf{2.} La lame est attaquée par \ce{HCl} (dégagement \ce{H2}) : ce n'est pas du cuivre. Elle ne réagit pas avec \ce{Fe^2+} : elle n'est pas plus réductrice que le fer, donc ce n'est ni de l'aluminium (qui réduit \ce{Fe^2+}, même lentement après grattage), ni du zinc (qui réduit \ce{Fe^2+} rapidement). Elle réagit avec \ce{Cu^2+} (dépôt Cu) : elle est plus réductrice que le cuivre. Le seul métal vérifiant $Fe > M > Cu$ (ou $M = Fe$) est le \textbf{fer}.
\textit{Note : Si c'était de l'aluminium gratté, il aurait réagi avec \ce{Fe^2+} (lentement). Si c'était du zinc, il aurait réagi rapidement avec \ce{Fe^2+}.}

\textbf{3.} Équation-bilan : $\ce{Zn(s) + 2H+(aq) -> Zn^2+(aq) + H2(g)}$.
Quantité de zinc : $n(\ce{Zn}) = \dfrac{m}{M(\ce{Zn})} = \dfrac{3,27}{65,4} = \SI{0,050}{mol}$.
D'après la stœchiométrie (1 mol Zn $\to$ 1 mol \ce{H2}), $n(\ce{H2}) = n(\ce{Zn}) = \SI{0,050}{mol}$.
Volume de dihydrogène : $V(\ce{H2}) = n \times V_m = 0,050 \times 24 = \SI{1,2}{L}$.
Le zinc est entièrement consommé (masse consommée = \SI{3,27}{g}) et il se dégage \SI{1,2}{L} de dihydrogène.
\end{corrige}

\newpage

\coursec{Travaux pratiques}

\courssub{Objectifs du TP}

Ce travail pratique te place dans la situation d'un technicien de laboratoire qui doit identifier des métaux inconnus à partir de leur comportement chimique. À l'issue de la séance, tu seras capable de :

\begin{itemize}
\item réaliser l'action des acides chlorhydrique et sulfurique \emph{dilués} sur les métaux usuels (zinc, fer, cuivre, aluminium) ;
\item identifier le gaz dégagé par le test à la flamme (détonation caractéristique du dihydrogène) ;
\item mettre en évidence les ions métalliques formés à l'aide d'une solution d'hydroxyde de sodium ;
\item réaliser l'action d'un métal (zinc, fer) sur des solutions contenant des ions métalliques (\ce{Cu^2+}, \ce{Zn^2+}, \ce{Fe^2+}) ;
\item écrire les demi-équations électroniques et les équations-bilans des réactions observées ;
\item établir un \cle{classement qualitatif} des métaux étudiés selon leur pouvoir réducteur.
\end{itemize}

\courssub{Matériel et produits}

\begin{center}
\rowcolors{2}{popblueL}{white}
\begin{tabularx}{\linewidth}{|l|X|}
\hline
\rowcolor{popblue!40} \textbf{Matériel} & \textbf{Produits} \\
\hline
8 tubes à essai propres et secs & Acide chlorhydrique dilué (\SI{1}{mol.L^{-1}}) \\
Portoir à tubes & Acide sulfurique dilué (\SI{0,5}{mol.L^{-1}}) \\
Pipette graduée de \SI{5}{mL} & Solution de sulfate de cuivre \ce{CuSO4} (bleue) \\
Bouchons à doigt (non hermétiques) & Solution de sulfate de zinc \ce{ZnSO4} (incolore) \\
Allumettes et briquet & Solution de sulfate de fer(II) \ce{FeSO4} (verdâtre) \\
Spatule, pince en bois & Hydroxyde de sodium dilué \ce{NaOH} (\SI{1}{mol.L^{-1}}) \\
Grenaille de zinc & Eau distillée \\
Clous en fer décapés & \\
Lame de cuivre (fil électrique dénudé) & \\
Feuille d'aluminium (papier d'emballage alimentaire) & \\
\hline
\end{tabularx}
\end{center}

\begin{savaistu}[Alternatives avec le matériel local]
Si le sulfate de cuivre manque, on peut utiliser la \emph{bleu de vitriol} vendue en quincaillerie comme fongicide. Le zinc peut être récupéré dans le boîtier des piles usagées (à nettoyer soigneusement), et le cuivre dans les fils électriques dénudés. Les clous de charpentier, bien décapés à la toile émeri, remplacent avantageusement la tournure de fer. L'acide sulfurique dilué peut être remplacé par l'acide de batterie (dilué avec prudence : \emph{verser l'acide dans l'eau}).
\end{savaistu}

\courssub{Sécurité}

\begin{attention}[Consignes de sécurité impératives]
\begin{itemize}
\item \textbf{Pictogramme « corrosif »} (liquide attaquant une main et une surface) : les acides chlorhydrique et sulfurique, même dilués, ainsi que la soude, attaquent la peau et les yeux. Port de la \textbf{blouse} et des \textbf{lunettes} obligatoires. Toute manipulation d'acide se fait \textbf{sous la supervision de l'enseignant}.
\item En cas de contact avec la peau : rincer \textbf{abondamment à l'eau claire} pendant plusieurs minutes, puis prévenir l'enseignant.
\item \textbf{Ne jamais boucher hermétiquement} un tube où se dégage un gaz : la pression pourrait faire exploser le tube. On bouche « au doigt » de façon lâche, le temps du test.
\item Le dihydrogène est \textbf{inflammable} (pictogramme « flamme ») : éloigner toute flamme des tubes pendant le dégagement, n'approcher l'allumette qu'au moment du test.
\item Le sulfate de cuivre est \textbf{nocif et polluant} : ne pas le jeter à l'évier ; le recueillir dans le bidon de récupération prévu.
\item Ne jamais goûter ni sentir directement un produit : pour percevoir une odeur, on \emph{évente} avec la main au-dessus du tube.
\end{itemize}
\end{attention}

\courssub{Schéma du dispositif}

\begin{popfigure}
\begin{center}
\begin{tikzpicture}[scale=0.95, every node/.style={font=\small}]
% Tube 1 : acide + zinc
\draw[popdark, thick] (0,0) -- (0,3.2);
\draw[popdark, thick] (1.2,0) -- (1.2,3.2);
\draw[popdark, thick] (0,0) arc (180:360:0.6);
\fill[popblue!25] (0,0.05) arc (180:360:0.6) -- (1.2,1.4) -- (0,1.4) -- cycle;
\draw[popblue, thick] (0,1.4) -- (1.2,1.4);
% grenaille de zinc
\foreach \x/\y in {0.25/0.25,0.55/0.18,0.85/0.28,0.4/0.4,0.75/0.45}{
  \fill[popmuted] (\x,\y) circle (0.07);}
% bulles
\foreach \x/\y in {0.3/0.9,0.6/1.1,0.8/0.75,0.45/1.25,0.9/1.2}{
  \draw[popblue] (\x,\y) circle (0.045);}
\node[popdark] at (0.6,-0.6) {Tube 1 : \ce{HCl} + \ce{Zn}};
% Tube 2 : test à la flamme
\begin{scope}[xshift=3.6cm]
\draw[popdark, thick] (0,0) -- (0,3.2);
\draw[popdark, thick] (1.2,0) -- (1.2,3.2);
\draw[popdark, thick] (0,0) arc (180:360:0.6);
\fill[popblue!25] (0,0.05) arc (180:360:0.6) -- (1.2,1.4) -- (0,1.4) -- cycle;
\draw[popblue, thick] (0,1.4) -- (1.2,1.4);
% allumette enflammée
\draw[thick, brown] (2.1,3.9) -- (1.35,3.35);
\fill[poporange] (1.3,3.3) .. controls (1.15,3.6) and (1.45,3.65) .. (1.3,3.3);
\fill[poppink] (1.3,3.32) .. controls (1.22,3.5) and (1.38,3.52) .. (1.3,3.32);
\node[poporange] at (2.6,4.1) {détonation};
\node[popdark] at (0.6,-0.6) {Tube 2 : test du gaz \ce{H2}};
\end{scope}
% Tube 3 : Zn dans CuSO4
\begin{scope}[xshift=7.2cm]
\draw[popdark, thick] (0,0) -- (0,3.2);
\draw[popdark, thick] (1.2,0) -- (1.2,3.2);
\draw[popdark, thick] (0,0) arc (180:360:0.6);
\fill[popblue!45] (0,0.05) arc (180:360:0.6) -- (1.2,1.8) -- (0,1.8) -- cycle;
\draw[popblue, thick] (0,1.8) -- (1.2,1.8);
% lame de zinc
\draw[thick, fill=popmuted!60] (0.45,0.1) rectangle (0.75,2.6);
\fill[poporange!80] (0.45,0.1) rectangle (0.75,0.9);
\node[popdark, align=center] at (0.6,-0.6) {Tube 3 : lame de \ce{Zn}\\ dans \ce{CuSO4}};
\node[poporange, align=center] at (2.3,0.6) {dépôt\\ rouge};
\draw[->, poporange] (1.7,0.6) -- (0.85,0.55);
\end{scope}
\end{tikzpicture}
\end{center}
\legende{Dispositifs du TP : action de l'acide chlorhydrique dilué sur le zinc avec dégagement gazeux (tube 1), test d'identification du dihydrogène à la flamme (tube 2), action du zinc sur une solution de sulfate de cuivre avec dépôt rouge de cuivre (tube 3).}
\end{popfigure}

\courssub{Protocole}

\textbf{Partie A : action des acides dilués sur les métaux}\par

\begin{enumerate}
\item Numéroter huit tubes à essai de A1 à A8 et les placer sur le portoir.
\item Introduire dans les tubes A1 à A4 une pointe de spatule de grenaille de zinc, un clou décapé, un morceau de lame de cuivre et un petit carré d'aluminium respectivement.
\item Verser avec la pipette graduée environ \SI{2}{mL} d'acide chlorhydrique dilué dans chacun des tubes A1 à A4. Observer immédiatement.
\item Recommencer avec les tubes A5 à A8 en utilisant l'acide sulfurique dilué à la place de l'acide chlorhydrique.
\item Pour les tubes où un dégagement gazeux est visible : boucher \emph{lâchement} le tube au doigt pendant une minute, puis approcher rapidement une allumette enflammée de l'ouverture. Noter le résultat (une \cle{détonation} caractéristique signale le dihydrogène \ce{H2}).
\item Ajouter ensuite quelques gouttes de solution d'hydroxyde de sodium dans chaque tube ayant réagi : la formation d'un \cle{précipité} met en évidence les ions métalliques formés (blanc pour \ce{Zn^2+}, verdâtre pour \ce{Fe^2+}, blanc gélatineux pour \ce{Al^3+} soluble en excès de soude).
\end{enumerate}

\textbf{Partie B : action d'un métal sur un ion métallique}\par

\begin{enumerate}
\item Verser environ \SI{3}{mL} de solution de sulfate de cuivre (bleue) dans les tubes B1 et B2, et \SI{3}{mL} de solution de sulfate de zinc (incolore) dans le tube B3.
\item Plonger une lame de zinc dans le tube B1, un clou en fer dans le tube B2, et une lame de cuivre dans le tube B3.
\item Laisser reposer \textbf{10 minutes} sans agiter, puis observer : aspect du métal (dépôt éventuel), couleur de la solution.
\item Consigner toutes les observations dans le tableau de résultats.
\end{enumerate}

\courssub{Observations et résultats attendus}

\begin{center}
\rowcolors{2}{popgreenL}{white}
\begin{tabularx}{\linewidth}{|l|X|X|X|}
\hline
\rowcolor{popgreen!40} \textbf{Tube} & \textbf{Système} & \textbf{Observations} & \textbf{Conclusion} \\
\hline
A1 / A5 & \ce{Zn} + acide dilué & Dégagement gazeux vif ; détonation à la flamme ; précipité blanc avec \ce{NaOH} & Réaction : \ce{H2} et \ce{Zn^2+} formés \\
A2 / A6 & \ce{Fe} + acide dilué & Dégagement lent ; détonation faible ; précipité verdâtre avec \ce{NaOH} & Réaction : \ce{H2} et \ce{Fe^2+} formés \\
A3 / A7 & \ce{Cu} + acide dilué & Aucune observation & Pas de réaction \\
A4 / A8 & \ce{Al} + acide dilué & Dégagement après un temps de latence (couche d'alumine) ; détonation ; précipité blanc gélatineux soluble en excès de \ce{NaOH} & Réaction : \ce{H2} et \ce{Al^3+} formés \\
B1 & \ce{Zn} dans \ce{CuSO4} & Dépôt rouge brique sur le zinc ; la solution bleue se décolore & Réaction : \ce{Cu} et \ce{Zn^2+} formés \\
B2 & \ce{Fe} dans \ce{CuSO4} & Dépôt rouge sur le clou ; décoloration partielle & Réaction : \ce{Cu} et \ce{Fe^2+} formés \\
B3 & \ce{Cu} dans \ce{ZnSO4} & Aucune observation & Pas de réaction \\
\hline
\end{tabularx}
\end{center}

\courssub{Exploitation}

\begin{definition}[Vocabulaire de l'oxydoréduction]
\begin{itemize}
\item \textbf{Oxydation} : perte d'un ou plusieurs électrons par une espèce chimique. Le nombre d'oxydation augmente.
\item \textbf{Réduction} : gain d'un ou plusieurs électrons par une espèce chimique. Le nombre d'oxydation diminue.
\item \textbf{Réducteur} : espèce qui perd des électrons (elle s'oxyde).
\item \textbf{Oxydant} : espèce qui gagne des électrons (elle se réduit).
\item \textbf{Couple oxydant/réducteur} : noté \ce{Ox/Red}, il associe l'oxydant et le réducteur d'une même transformation (ex. \ce{Zn^2+/Zn}, \ce{H+/H2}, \ce{Cu^2+/Cu}).
\item \textbf{Réaction d'oxydoréduction} : réaction au cours de laquelle se produit un transfert d'électrons d'un réducteur vers un oxydant.
\end{itemize}
\end{definition}

\begin{methode}[Écrire l'équation-bilan d'une réaction d'oxydoréduction]
\begin{enumerate}
\item Identifier les deux couples en présence (\ce{Ox1/Red1} et \ce{Ox2/Red2}).
\item Écrire les deux demi-équations électroniques (réduction de l'oxydant, oxydation du réducteur) en équilibrant les atomes et les charges (ajout d'électrons \ce{e-}).
\item Multiplier les demi-équations par des coefficients entiers pour obtenir le même nombre d'électrons des deux côtés.
\item Additionner membre à membre les deux demi-équations et simplifier les électrons (et les éventuelles espèces communes).
\item Vérifier l'équilibrage final : atomes et charges.
\end{enumerate}
\end{methode}

\begin{exoresolu}[Questions guidées d'exploitation]
\textbf{Question 1.} Quel est le gaz dégagé dans les tubes A1, A2, A4 ? Comment l'identifie-t-on ?

\textbf{Question 2.} Écrire les demi-équations électroniques puis l'équation-bilan de l'action de l'acide chlorhydrique sur le zinc.

\textbf{Question 3.} Le cuivre est-il attaqué par les acides dilués ? Proposer une interprétation.

\textbf{Question 4.} Écrire les demi-équations et l'équation-bilan de la réaction observée dans le tube B1. Quelle espèce est oxydée ? Quelle espèce est réduite ?

\textbf{Question 5.} En comparant les tubes B1 et B3, établir un classement qualitatif des métaux zinc, fer et cuivre selon leur pouvoir réducteur.
\tcblower
\textbf{Réponse 1.} Le gaz est le \cle{dihydrogène} \ce{H2}. Approché d'une flamme, il produit une \emph{détonation} (un « aboiement ») caractéristique : il brûle avec l'oxygène de l'air selon \ce{2H2 + O2 -> 2H2O}.

\textbf{Réponse 2.} Le zinc perd des électrons (oxydation) : \ce{Zn -> Zn^2+ + 2e-}. Les ions hydrogène gagnent des électrons (réduction) : \ce{2H+ + 2e- -> H2}. En additionnant membre à membre, les électrons s'éliminent :
\[ \ce{Zn + 2H+ -> Zn^2+ + H2} \]

\textbf{Réponse 3.} Non : aucun dégagement n'est observé avec le cuivre. Le cuivre n'est \emph{pas assez réducteur} pour céder ses électrons aux ions \ce{H+} des acides dilués ; il appartient, comme l'argent et l'or, aux métaux dits « nobles » vis-à-vis de ces acides.

\textbf{Réponse 4.} Oxydation du zinc : \ce{Zn -> Zn^2+ + 2e-} ; réduction des ions cuivre : \ce{Cu^2+ + 2e- -> Cu}. Équation-bilan :
\[ \ce{Zn + Cu^2+ -> Zn^2+ + Cu} \]
Le zinc, qui perd des électrons, est \emph{oxydé} (c'est le réducteur, couple \ce{Zn^2+/Zn}) ; l'ion \ce{Cu^2+}, qui gagne des électrons, est \emph{réduit} (c'est l'oxydant, couple \ce{Cu^2+/Cu}). Le dépôt rouge est du cuivre métallique ; la décoloration traduit la disparition des ions \ce{Cu^2+} bleus.

\textbf{Réponse 5.} Le zinc et le fer réduisent les ions \ce{Cu^2+} (tubes B1 et B2), donc le zinc et le fer sont \emph{plus réducteurs} que le cuivre. Inversement, le cuivre ne réduit pas les ions \ce{Zn^2+} (tube B3), ce qui confirme ce classement. D'où l'ordre de pouvoir réducteur décroissant :
\[ \ce{Zn} > \ce{Fe} > \ce{Cu} \]
\end{exoresolu}

\begin{aretenir}[Bilan du TP : réactivité des métaux vis-à-vis des acides et déplacement]
\begin{minipage}{\linewidth}
\begin{itemize}
\item \textbf{Action des acides dilués ($\ce{H+}$) :} Les métaux \ce{Zn}, \ce{Fe}, \ce{Al} (après latence) réagissent : $\ce{M} + n\,\ce{H+} \rightarrow \ce{M}^{n+} + \frac{n}{2}\,\ce{H2}$. Le cuivre (\ce{Cu}), l'argent (\ce{Ag}), l'or (\ce{Au}) ne réagissent pas.
\item \textbf{Réactions de déplacement :} Un métal $M_1$ réduit l'ion $M_2^{n+}$ si $M_1$ est un meilleur réducteur que $M_2$ (i.e. $M_1$ se place au-dessus de $M_2$ dans l'échelle des pouvoirs réducteurs). La réaction est : $M_1 + M_2^{n+} \rightarrow M_1^{n+} + M_2$.
\item \textbf{Classement établi :} $\ce{Zn} > \ce{Fe} > \ce{Cu}$. Ce classement permet de \emph{prévoir} si une réaction aura lieu entre un métal et un ion métallique, ou entre un métal et un acide dilué.
\end{itemize}
\end{minipage}
\end{aretenir}

\courssub{Conclusion}

Ce TP a permis de vérifier expérimentalement deux grandes familles de réactions d'\cle{oxydoréduction} en solution aqueuse. D'une part, les acides chlorhydrique et sulfurique \emph{dilués} attaquent les métaux comme le zinc, le fer et l'aluminium avec dégagement de dihydrogène et formation d'ions métalliques, mais restent sans action sur le cuivre (comme sur l'argent et l'or). D'autre part, un métal peut réduire les ions d'un autre métal : le zinc et le fer se recouvrent de cuivre rouge dans une solution de sulfate de cuivre, tandis que le cuivre est sans effet sur une solution de sulfate de zinc. Chaque réaction s'interprète par un \emph{transfert d'électrons} du métal (réducteur, qui s'oxyde) vers l'ion (oxydant, qui se réduit). L'ensemble des observations conduit au classement qualitatif $\ce{Zn} > \ce{Fe} > \ce{Cu}$ selon le pouvoir réducteur, classement que la suite du module permettra de généraliser et d'exploiter pour \emph{prévoir} les réactions possibles — par exemple pour identifier un métal inconnu, comme le ferait un laboratoire d'analyse.

\newpage

\coursec{Méthodes et exercices résolus}

\coursec{Généralités sur l'oxydoréduction en solution aqueuse}

\courssub{1. Action d'un acide dilué sur un métal}

\begin{experience}[Action de l'acide chlorhydrique dilué sur la grenaille de zinc]
\textbf{Matériel :} tube à essai, bec Bunsen, grenaille de zinc, acide chlorhydrique dilué (\SI{1}{mol.L^{-1}}), tube à essai retourné, allumette.
\textbf{Protocole :}
\begin{enumerate}
\item Introduire quelques grammes de grenaille de zinc dans le tube à essai.
\item Verser \SI{10}{mL} d'acide chlorhydrique dilué.
\item Observer l'évolution au contact du métal.
\item Recueillir le gaz dégagé par déplacement de l'air (tube retourné) et approcher une flamme.
\end{enumerate}
\textbf{Observations :} dégagement gazeux immédiat et abondant à la surface du zinc (effervescence), le zinc se dissout progressivement, la solution reste incolore. Le gaz recueilli produit une détonation caractéristique (aboiement) à la flamme : c'est le dihydrogène \ce{H2}.
\textbf{Interprétation :} le zinc est un métal \cle{plus réducteur que l'hydrogène} : il cède des électrons aux ions \ce{H+} de l'acide. Il y a réaction d'oxydoréduction.
\begin{popfigure}
\begin{tikzpicture}[scale=0.9, every node/.style={font=\small}]
% Becher
\draw[thick, popblue] (0,0) -- (0,4) -- (3,4) -- (3,0);
\draw[thick, popblue] (-0.2,0) -- (3.2,0);
% Zinc granules
\foreach \x/\y in {0.5,0.5 / 1.5,0.5 / 2.5,0.5 / 0.8,1.2 / 1.8,1.2 / 2.8,1.2} {
    \fill[popgray!50] (\x,\y) circle (0.15);
}
% Acid liquid
\fill[popblueL!50] (0.1,0.1) rectangle (2.9,3.9);
% Bubbles
\foreach \i in {1,...,15} {
    \pgfmathsetmacro{\bx}{0.3+rnd*2.4}
    \pgfmathsetmacro{\by}{0.3+rnd*3.4}
    \pgfmathsetmacro{\br}{0.05+rnd*0.1}
    \draw[popblue, thick] (\bx,\by) circle (\br);
}
% Gas collection tube
\draw[thick, popdark] (3.5,2) -- (3.5,5) -- (6,5) -- (6,2);
\draw[thick, popdark] (3.3,2) -- (6.2,2);
\fill[popgray!20] (3.5,2) rectangle (6,5);
% Flame
\draw[poporange, thick] (6.5,1.5) -- (6.5,0.5);
\draw[poporange, decorate, decoration={snake, amplitude=1pt, segment length=3pt}] (6.5,1.5) -- (6.5,2.5);
\node[popdark, anchor=south] at (1.5,4.3) {Solution incolore};
\node[popdark, anchor=north] at (1.5,-0.5) {Grenaille de zinc};
\node[popdark, anchor=west] at (6.2,3.5) {Gaz : \ce{H2}};
\node[popdark, anchor=west] at (6.2,1) {Flamme $\rightarrow$ \emph{aboiement}};
\end{tikzpicture}
\legende{Dispositif d'action de l'acide chlorhydrique dilué sur le zinc et test du dihydrogène.}
\end{popfigure}
\end{experience}

\begin{aretenir}[Condition de réaction acide / métal]
Un acide dilué (\ce{HCl}, \ce{H2SO4}) fournit des ions \ce{H+} (souvent notés \ce{H3O+} en solution aqueuse).
\begin{itemize}
\item \textbf{Métaux attaqués (réducteurs plus forts que \ce{H2}) :} Zn, Fe, Al, Mg, etc. $\rightarrow$ dégagement de \ce{H2}, dissolution du métal, formation des ions métalliques correspondants (\ce{Zn^2+}, \ce{Fe^2+}, \ce{Al^3+}...).
\item \textbf{Métaux non attaqués (réducteurs plus faibles que \ce{H2}) :} Cu, Ag, Au, Pt... $\rightarrow$ \textbf{aucune réaction} visible (pas de gaz, pas de dissolution).
\end{itemize}
\textbf{Attention à l'aluminium :} sa couche superficielle d'alumine \ce{Al2O3} (passivation) peut retarder le démarrage de la réaction.
\end{aretenir}

\begin{methode}[Décrire et interpréter l'action d'un acide dilué sur un métal]
\begin{enumerate}
\item \textbf{Identifier les espèces présentes.} L'acide dilué apporte \ce{H+} ; le métal apporte des atomes $M$.
\item \textbf{Se poser la question clé : le métal est-il plus réducteur que l'hydrogène ?} Consulter le classement des métaux (Zn, Fe, Al, Mg $>$ H $>$ Cu, Ag, Au).
\item \textbf{Observer et décrire} : effervescence (\ce{H2} gazeux, test de l'aboiement), disparition du métal, éventuelle coloration de la solution (\ce{Fe^2+} vert pâle, \ce{Cu^2+} bleu, \ce{Zn^2+}/\ce{Al^3+} incolores).
\item \textbf{Écrire les demi-équations électroniques} :
      \begin{itemize}
      \item Oxydation du métal : $M \ce{->} M^{n+} + n\,e^-$
      \item Réduction des ions hydrogène : $\ce{2H+ + 2e- -> H2}$
      \end{itemize}
\item \textbf{Écrire l'équation-bilan} en égalisant les électrons (PPCM), additionner membre à membre, simplifier les électrons, vérifier atomes et charges.
\end{enumerate}
\end{methode}

\begin{exoresolu}[Le zinc dans l'acide chlorhydrique]
On introduit \SI{6,54}{g} de grenaille de zinc dans \SI{100}{mL} d'une solution d'acide chlorhydrique de concentration \SI{1,00}{mol.L^{-1}}. On donne $M(\ce{Zn}) = \SI{65,4}{g.mol^{-1}}$ et $V_m = \SI{24,0}{L.mol^{-1}}$ (conditions normales de température et de pression).
\begin{enumerate}
\item Décrire les observations expérimentales.
\item Écrire les demi-équations et l'équation-bilan.
\item Calculer le volume de gaz dégagé.
\end{enumerate}
\tcblower
\textbf{1. Observations.} On observe un dégagement gazeux abondant et immédiat à la surface du zinc (effervescence), qui se dissout progressivement. Le gaz recueilli produit une détonation brève (aboiement) à l'approche d'une flamme : c'est le dihydrogène \ce{H2}. La solution reste incolore car l'ion \ce{Zn^2+} est incolore.

\textbf{2. Équations.} Le zinc est un métal plus réducteur que l'hydrogène : il est attaqué.
\begin{itemize}
\item Oxydation du zinc : \ce{Zn -> Zn^2+ + 2e-}
\item Réduction des ions hydrogène : \ce{2H+ + 2e- -> H2}
\end{itemize}
Les deux demi-équations échangent le même nombre d'électrons (2) ; on les additionne membre à membre :
\[ \ce{Zn + 2H+ -> Zn^2+ + H2} \]
\textbf{Vérification :} Atomes : 1 Zn, 2 H de chaque côté. Charges : $+2$ à gauche, $+2$ à droite. L'équation est équilibrée.

\textbf{3. Volume de gaz.} Quantités de matière initiales :
\[ n_i(\ce{Zn}) = \frac{m}{M} = \frac{6,54}{65,4} = \SI{0,100}{mol} \qquad n_i(\ce{H+}) = C \times V = 1,00 \times 0,100 = \SI{0,100}{mol} \]

\begin{center}
\begin{tabularx}{\linewidth}{|l|X|X|X|X|}
\hline
\rowcolor{popblueL} Équation & \ce{Zn} & $+$ \ce{2H+} & \ce{->} \ce{Zn^2+} & $+$ \ce{H2} \\
\hline
État initial (mol) & 0,100 & 0,100 & 0 & 0 \\
\hline
État final (mol) & $0,100 - x_f$ & $0,100 - 2x_f$ & $x_f$ & $x_f$ \\
\hline
\end{tabularx}
\end{center}

Hypothèses sur l'avancement final $x_f$ :
\begin{itemize}
\item Si Zn limitant : $x_f = \SI{0,100}{mol}$.
\item Si \ce{H+} limitant : $0,100 - 2x_f = 0 \Rightarrow x_f = \SI{0,050}{mol}$.
\end{itemize}
Le plus petit avancement l'emporte : $x_f = \SI{0,050}{mol}$ (les ions \ce{H+} sont le réactif limitant).

D'après le tableau, $n(\ce{H2}) = x_f = \SI{0,050}{mol}$, d'où :
\[ V(\ce{H2}) = n \times V_m = 0,050 \times 24,0 = \SI{1,20}{L} \]

\textbf{Conclusion :} il se dégage \SI{1,20}{L} de dihydrogène. Il reste $0,100 - 0,050 = \SI{0,050}{mol}$ de zinc non attaqué (soit \SI{3,27}{g}).
\end{exoresolu}

\courssub{2. Réaction entre un ion métallique et un métal}

\begin{experience}[Déplacement du cuivre par le fer : « l'arbre de Noël » chimique]
\textbf{Matériel :} bécher, solution de sulfate de cuivre(II) \ce{CuSO4} (\SI{0,5}{mol.L^{-1}}, bleue), clou ou limaille de fer.
\textbf{Protocole :} Plonger le fer dans la solution bleue. Observer pendant quelques minutes.
\textbf{Observations :} apparition rapide d'un dépôt poudreux rouge-orangé (cuivre métallique) à la surface du fer. La solution bleue se décolore progressivement (disparition des \ce{Cu^2+}). La solution peut verdire (apparition des \ce{Fe^2+}).
\textbf{Interprétation :} Le fer est plus réducteur que le cuivre. Il réduit les ions \ce{Cu^2+} en cuivre métal et s'oxyde lui-même en ions \ce{Fe^2+}.
\begin{popfigure}
\begin{tikzpicture}[scale=0.9, every node/.style={font=\small}]
% Beaker
\draw[thick, popblue] (0,0) -- (0,4.5) -- (4,4.5) -- (4,0);
\draw[thick, popblue] (-0.2,0) -- (4.2,0);
% Solution initial (bleu)
\fill[popblueL!30] (0.1,0.1) rectangle (3.9,4.4);
% Iron nail
\draw[popgray, line width=3pt] (2,4.5) -- (2,1.5);
\draw[popgray, line width=3pt] (1.8,1.5) -- (2.2,1.5); % head
% Copper deposit (reddish)
\foreach \i in {1,...,30} {
    \pgfmathsetmacro{\bx}{1.7+rnd*0.6}
    \pgfmathsetmacro{\by}{1.6+rnd*2.5}
    \fill[poporange!80!popdark] (\bx,\by) circle (0.05);
}
% Color gradient for solution fading
\shade[left color=popblueL!60, right color=popgreenL!30] (0.1,0.1) rectangle (3.9,4.4);
\node[popdark, anchor=south] at (2,4.7) {Clou en fer};
\node[popdark, anchor=north] at (2,1.2) {Dépôt de Cu (rouge)};
\node[popdark, anchor=west] at (4.2,3) {Solution : bleu $\rightarrow$ incolore/vert};
\node[popdark, anchor=west] at (4.2,2.5) {(\ce{Cu^2+} $\downarrow$, \ce{Fe^2+} $\uparrow$)};
\end{tikzpicture}
\legende{Réaction de déplacement : fer + ions cuivre(II) $\rightarrow$ fer(II) + cuivre.}
\end{popfigure}
\end{experience}

\begin{methode}[Étudier l'action d'un métal sur un ion métallique]
\begin{enumerate}
\item \textbf{Identifier les deux couples} en présence : $M_1^{n+}/M_1$ (ion en solution) et $M_2^{p+}/M_2$ (métal introduit).
\item \textbf{Prévoir le sens de la réaction} : le métal le \cle{plus réducteur} (celui qui perd le plus facilement ses électrons, ex. Zn, Fe) s'oxyde ; l'ion du métal le moins réducteur (ex. \ce{Cu^2+}, \ce{Ag+}) se réduit et se dépose.
\item \textbf{Décrire les observations} : dépôt métallique (rougeâtre pour Cu, grisâtre pour Ag), décoloration ou changement de couleur de la solution (bleu \ce{Cu^2+} $\rightarrow$ incolore), attaque du métal introduit.
\item \textbf{Écrire les demi-équations} puis l'équation-bilan en égalisant les électrons (PPCM).
\item \textbf{Cas sans réaction} : si le métal introduit est moins réducteur que celui de l'ion (ex. Cu dans \ce{Zn^2+}), il ne se passe rien : l'écrire explicitement.
\end{enumerate}
\end{methode}

\begin{exoresolu}[Le clou de fer dans le sulfate de cuivre(II)]
On plonge un clou en fer de masse $m = \SI{5,60}{g}$ dans \SI{200}{mL} d'une solution de sulfate de cuivre(II) de concentration \SI{0,500}{mol.L^{-1}}. On donne $M(\ce{Fe}) = \SI{55,8}{g.mol^{-1}}$ et $M(\ce{Cu}) = \SI{63,5}{g.mol^{-1}}$.
\begin{enumerate}
\item Décrire les observations et écrire l'équation-bilan.
\item Déterminer la masse de cuivre déposée lorsque la réaction est totale.
\end{enumerate}
\tcblower
\textbf{1. Observations et équation.} Le fer est plus réducteur que le cuivre : il s'oxyde. On observe un dépôt rougeâtre de cuivre métallique sur le clou et la décoloration progressive de la solution bleue (disparition des ions \ce{Cu^2+}). La solution peut prendre une teinte vert pâle due aux ions \ce{Fe^2+} formés.
\begin{itemize}
\item Oxydation : \ce{Fe -> Fe^2+ + 2e-}
\item Réduction : \ce{Cu^2+ + 2e- -> Cu}
\end{itemize}
Équation-bilan (électrons déjà égaux) :
\[ \ce{Fe + Cu^2+ -> Fe^2+ + Cu} \]
\textbf{Vérification :} Atomes : 1 Fe, 1 Cu. Charges : $+2$ des deux côtés.

\textbf{2. Masse de cuivre déposée.} Quantités initiales :
\[ n_i(\ce{Fe}) = \frac{5,60}{55,8} = \SI{0,100}{mol} \qquad n_i(\ce{Cu^2+}) = 0,500 \times 0,200 = \SI{0,100}{mol} \]

\begin{center}
\begin{tabularx}{\linewidth}{|l|X|X|X|X|}
\hline
\rowcolor{popblueL} Équation & \ce{Fe} & $+$ \ce{Cu^2+} & \ce{->} \ce{Fe^2+} & $+$ \ce{Cu} \\
\hline
État initial (mol) & 0,100 & 0,100 & 0 & 0 \\
\hline
État final (mol) & $0,100 - x_f$ & $0,100 - x_f$ & $x_f$ & $x_f$ \\
\hline
\end{tabularx}
\end{center}

Les deux réactifs sont en proportions stœchiométriques exactes (coefficients 1 : 1 et quantités égales) : l'avancement final est $x_f = \SI{0,100}{mol}$.
\[ n(\ce{Cu}) = x_f = \SI{0,100}{mol} \quad\Longrightarrow\quad m(\ce{Cu}) = n \times M = 0,100 \times 63,5 = \SI{6,35}{g} \]

\textbf{Conclusion :} il se dépose \SI{6,35}{g} de cuivre sur le clou, tandis que $0,100 \times 55,8 = \SI{5,58}{g}$ de fer passent en solution : le clou a pratiquement disparu sous le dépôt de cuivre.
\end{exoresolu}

\courssub{3. Oxydation, réduction, oxydant, réducteur}

\begin{definition}[Oxydation et Réduction]
\begin{itemize}
\item \cle{Oxydation} : perte d'un ou plusieurs électrons par une espèce chimique. L'espèce qui s'oxyde voit son nombre d'oxydation \textbf{augmenter}.
\item \cle{Réduction} : gain d'un ou plusieurs électrons par une espèce chimique. L'espèce qui se réduit voit son nombre d'oxydation \textbf{diminuer}.
\end{itemize}
\textbf{Moyen mnémotechnique :} \textbf{Perte} $\rightarrow$ \textbf{Oxydation} ; \textbf{Gain} $\rightarrow$ \textbf{Réduction}. (Anode/Oxydation, Cathode/Réduction en pile).
\end{definition}

\begin{definition}[Oxydant, Réducteur et Couple redox]
\begin{itemize}
\item \cle{Oxydant} : espèce chimique qui \textbf{gagne} des électrons (elle se réduit). C'est l'espèce \emph{avant} réaction dans le couple.
\item \cle{Réducteur} : espèce chimique qui \textbf{perd} des électrons (elle s'oxyde). C'est l'espèce \emph{après} réaction dans le couple.
\item \cle{Couple oxydant/réducteur} : noté $\text{Ox}/\text{Red}$, il associe l'oxydant et le réducteur d'une même transformation (ex. \ce{Cu^2+}/\ce{Cu}, \ce{Zn^2+}/\ce{Zn}, \ce{H+}/\ce{H2}).
\end{itemize}
\textbf{Règle d'or :} Dans une réaction d'oxydoréduction, l'oxydant d'un couple réagit avec le réducteur de l'autre couple. Le réducteur \emph{réduit} l'oxydant en s'oxydant lui-même.
\end{definition}

\begin{propriete}[Conservation des électrons]
Au cours d'une réaction d'oxydoréduction, le nombre total d'électrons perdus par le(s) réducteur(s) est \textbf{exactement égal} au nombre total d'électrons gagnés par l'oxydant(s). Les électrons n'apparaissent \textbf{jamais} dans l'équation-bilan finale.
\end{propriete}

\begin{methode}[Identifier les rôles dans une réaction d'oxydoréduction]
\begin{enumerate}
\item \textbf{Repérer les transferts d'électrons} en écrivant les deux demi-équations électroniques.
\item \textbf{Appliquer les définitions} :
      \begin{itemize}
      \item L'espèce qui \textbf{perd} des électrons (électrons à droite) subit une \textbf{oxydation} : c'est le \textbf{réducteur}.
      \item L'espèce qui \textbf{gagne} des électrons (électrons à gauche) subit une \textbf{réduction} : c'est l'\textbf{oxydant}.
      \end{itemize}
\item \textbf{Vérifier la cohérence} : nombre d'électrons cédés = nombre d'électrons captés.
\item \textbf{Citer les couples} mis en jeu : $\text{Ox}_1/\text{Red}_1$ et $\text{Ox}_2/\text{Red}_2$.
\end{enumerate}
\end{methode}

\begin{exoresolu}[Qui est l'oxydant ? Qui est le réducteur ?]
On considère la réaction entre le zinc métallique et une solution de nitrate d'argent :
\[ \ce{Zn + 2Ag+ -> Zn^2+ + 2Ag} \]
\begin{enumerate}
\item Écrire les deux demi-équations électroniques.
\item Identifier l'oxydant, le réducteur, l'oxydation et la réduction.
\item Citer les deux couples oxydant/réducteur mis en jeu.
\end{enumerate}
\tcblower
\textbf{1. Demi-équations.}
\begin{itemize}
\item \ce{Zn -> Zn^2+ + 2e-} : le zinc perd 2 électrons (oxydation).
\item \ce{Ag+ + e- -> Ag} : chaque ion argent gagne 1 électron (réduction). Pour équilibrer l'échange, on multiplie par 2 : \ce{2Ag+ + 2e- -> 2Ag}.
\end{itemize}

\textbf{2. Rôles.}
\begin{itemize}
\item Le zinc \ce{Zn} \textbf{perd} des électrons : il subit une \textbf{oxydation}, c'est donc le \textbf{réducteur}.
\item L'ion \ce{Ag+} \textbf{gagne} des électrons : il subit une \textbf{réduction}, c'est donc l'\textbf{oxydant}.
\end{itemize}
Bilan des électrons : 2 électrons cédés par Zn = 2 électrons captés par 2 ions \ce{Ag+} : le transfert est équilibré.

\textbf{3. Couples mis en jeu.} \ce{Zn^2+}/\ce{Zn} et \ce{Ag+}/\ce{Ag}.

\textbf{Conclusion :} dans toute oxydoréduction, le réducteur (ici \ce{Zn}) s'oxyde pendant que l'oxydant (ici \ce{Ag+}) se réduit : les deux phénomènes sont inséparables et simultanés.
\end{exoresolu}

\courssub{4. Demi-équations électroniques et équation-bilan}

\begin{methode}[Construire une équation-bilan à partir des demi-équations]
\begin{enumerate}
\item \textbf{Écrire séparément} la demi-équation d'oxydation (réducteur à gauche, électrons à droite) et celle de réduction (oxydant et électrons à gauche).
\item \textbf{Équilibrer chaque demi-équation} : atomes (sauf O, H), puis charges avec les électrons $e^-$.
\item \textbf{Égaliser les électrons échangés} : multiplier chaque demi-équation par le coefficient qui rend le nombre d'électrons identique (PPCM des nombres d'électrons).
\item \textbf{Additionner membre à membre} et simplifier : les électrons doivent disparaître totalement.
\item \textbf{Vérifier} la conservation des atomes \emph{et} la conservation de la charge globale.
\end{enumerate}
\end{methode}

\begin{exoresolu}[L'aluminium réduit les ions cuivre(II)]
Une lame d'aluminium est plongée dans une solution de sulfate de cuivre(II). Un dépôt rouge de cuivre apparaît sur la lame. Écrire les demi-équations puis l'équation-bilan de la réaction.
\tcblower
\textbf{Étape 1 : demi-équations brutes (équilibrées atomes/charges).}
\begin{itemize}
\item Oxydation de l'aluminium : \ce{Al -> Al^3+ + 3e-}
\item Réduction des ions cuivre : \ce{Cu^2+ + 2e- -> Cu}
\end{itemize}

\textbf{Étape 2 : égalisation des électrons (PPCM de 3 et 2 = 6).}
\begin{itemize}
\item \ce{Al -> Al^3+ + 3e-} \quad ($\times 2$) $\Rightarrow$ \ce{2Al -> 2Al^3+ + 6e-}
\item \ce{Cu^2+ + 2e- -> Cu} \quad ($\times 3$) $\Rightarrow$ \ce{3Cu^2+ + 6e- -> 3Cu}
\end{itemize}

\textbf{Étape 3 : addition membre à membre et simplification.}
\[ \ce{2Al + 3Cu^2+ + 6e- -> 2Al^3+ + 6e- + 3Cu} \]
Les 6 électrons se simplifient des deux côtés :
\[ \ce{2Al + 3Cu^2+ -> 2Al^3+ + 3Cu} \]

\textbf{Étape 4 : vérification.}
\begin{itemize}
\item Atomes : 2 Al et 3 Cu de chaque côté. \checkmark
\item Charges : à gauche $3 \times (+2) = +6$ ; à droite $2 \times (+3) = +6$. \checkmark
\end{itemize}
L'équation est correcte.

\textbf{Conclusion :} l'aluminium, réducteur plus fort que le cuivre, s'oxyde en ions \ce{Al^3+} (incolores) pendant que les ions \ce{Cu^2+} bleus se réduisent en cuivre métallique rouge : la solution se décolore au fur et à mesure.
\end{exoresolu}

\courssub{5. Prévision des réactions et applications}

\begin{methode}[Prévoir l'action (ou la non-action) d'un acide dilué sur un métal]
\begin{enumerate}
\item \textbf{Classer mentalement les métaux} par rapport à l'hydrogène (couple \ce{H+}/\ce{H2}) :
      \begin{itemize}
      \item \textbf{Attaqués} (réducteurs plus forts que \ce{H2}) : Mg, Al, Zn, Fe, Sn, Pb... $\rightarrow$ dégagement de \ce{H2}.
      \item \textbf{Non attaqués} (réducteurs plus faibles que \ce{H2}) : Cu, Ag, Au, Pt... $\rightarrow$ \textbf{aucune réaction}.
      \end{itemize}
\item \textbf{Conclure sur l'existence de la réaction} et rédiger la réponse en trois temps : observation attendue, demi-équations, équation-bilan. Si aucune réaction : justifier explicitement (« le cuivre est moins réducteur que l'hydrogène »).
\item \textbf{Attention à l'aluminium} : il est théoriquement très réducteur, mais sa couche d'alumine \ce{Al2O3} protectrice peut ralentir le début de la réaction (temps de latence).
\end{enumerate}
\end{methode}

\begin{exoresolu}[Trois métaux face à l'acide sulfurique dilué]
On dispose de trois béchers contenant chacun \SI{50}{mL} d'acide sulfurique dilué. On introduit respectivement un clou en fer, une pièce en cuivre et une feuille d'aluminium. Pour chaque bécher, dire si une réaction a lieu ; le cas échéant, écrire l'équation-bilan.
\tcblower
\textbf{Bécher 1 : le fer.} Le fer est plus réducteur que l'hydrogène : il est attaqué. On observe un dégagement de \ce{H2} et la solution verdit (formation d'ions \ce{Fe^2+}).
\begin{itemize}
\item \ce{Fe -> Fe^2+ + 2e-} et \ce{2H+ + 2e- -> H2}
\end{itemize}
\[ \ce{Fe + 2H+ -> Fe^2+ + H2} \]

\textbf{Bécher 2 : le cuivre.} Le cuivre est moins réducteur que l'hydrogène : \textbf{aucune réaction} n'a lieu avec l'acide sulfurique dilué. La pièce reste intacte, la solution reste incolore. C'est pourquoi on utilise le cuivre pour les tuyauteries et les toitures (résistance à la corrosion acide).

\textbf{Bécher 3 : l'aluminium.} L'aluminium est plus réducteur que l'hydrogène : il est attaqué (parfois après un temps de latence dû à la couche d'alumine). L'aluminium s'oxyde en \ce{Al^3+} :
\begin{itemize}
\item \ce{Al -> Al^3+ + 3e-} ($\times 2$) et \ce{2H+ + 2e- -> H2} ($\times 3$)
\end{itemize}
On égalise les électrons (6 échangés) :
\[ \ce{2Al + 6H+ -> 2Al^3+ + 3H2} \]
\textbf{Vérification :} 2 Al et 6 H de chaque côté ; charge $+6$ de chaque côté.
\end{exoresolu}

\begin{methode}[Identifier un métal ou prévoir une réaction (démarche type Probatoire)]
\begin{enumerate}
\item \textbf{Exploiter les tests qualitatifs} :
      \begin{itemize}
      \item Dégagement de \ce{H2} avec un acide dilué $\Rightarrow$ métal plus réducteur que \ce{H2} (Zn, Fe, Al...).
      \item Absence de réaction avec acide dilué $\Rightarrow$ métal noble (Cu, Ag, Au).
      \end{itemize}
\item \textbf{Utiliser la couleur des ions en solution} : \ce{Cu^2+} (bleu), \ce{Fe^2+} (vert pâle), \ce{Fe^3+} (jaune-orangé), \ce{Zn^2+}, \ce{Al^3+} (incolores).
\item \textbf{Raisonner sur les dépôts métalliques} : un dépôt sur une lame indique que le métal de la lame a réduit les ions en solution $\Rightarrow$ il est plus réducteur que le métal déposé.
\item \textbf{Conclure par un classement} : ordonner les métaux du plus réducteur au moins réducteur à partir des expériences (ex. Zn $>$ Fe $>$ Cu $>$ Ag).
\item \textbf{Rédiger proprement} : observation $\rightarrow$ interprétation $\rightarrow$ équation-bilan $\rightarrow$ conclusion.
\end{enumerate}
\end{methode}

\begin{exoresolu}[Identification de deux métaux inconnus]
Un laborantin de lycée à Douala dispose de deux lamelles métalliques non étiquetées, $A$ et $B$, qui sont soit du zinc, soit du cuivre. Il réalise deux tests :
\begin{itemize}
\item Test 1 : la lamelle $A$ plongée dans l'acide chlorhydrique dilué produit un dégagement gazeux qui détone à la flamme ; la lamelle $B$ ne réagit pas.
\item Test 2 : la lamelle $A$ plongée dans une solution bleue de sulfate de cuivre se recouvre d'un dépôt rougeâtre.
\end{itemize}
Identifier $A$ et $B$ en justifiant chaque réponse par une équation-bilan.
\tcblower
\textbf{Test 1.} Le gaz qui détone à la flamme est le dihydrogène \ce{H2}. Seul un métal plus réducteur que l'hydrogène dégage \ce{H2} avec un acide dilué : c'est le cas du zinc, pas du cuivre. Donc $A$ est le \textbf{zinc} et $B$ est le \textbf{cuivre}.
\[ \ce{Zn + 2H+ -> Zn^2+ + H2} \]

\textbf{Test 2 (confirmation).} Le dépôt rougeâtre sur la lamelle $A$ est du cuivre métallique : les ions \ce{Cu^2+} de la solution ont été réduits par le métal $A$, qui s'est donc oxydé. Le zinc est effectivement plus réducteur que le cuivre :
\begin{itemize}
\item \ce{Zn -> Zn^2+ + 2e-} et \ce{Cu^2+ + 2e- -> Cu}
\end{itemize}
\[ \ce{Zn + Cu^2+ -> Zn^2+ + Cu} \]
La solution bleue se décolore progressivement, ce qui confirme la disparition des ions \ce{Cu^2+}.

\textbf{Conclusion :} $A$ est le zinc (métal réducteur, attaqué par les acides, capable de réduire \ce{Cu^2+}) et $B$ est le cuivre (métal peu réducteur, inerte vis-à-vis des acides dilués). Le classement obtenu est : \ce{Zn} $>$ \ce{H2} $>$ \ce{Cu} en pouvoir réducteur.
\end{exoresolu}

\begin{savaistu}[Oxydoréduction au Cameroun : de la corrosion à l'industrie]
\begin{itemize}
\item \textbf{Corrosion du fer (rouille) :} C'est une oxydoréduction lente où le fer s'oxyde en \ce{Fe^2+} puis \ce{Fe^3+} (rouille \ce{Fe2O3.nH2O}) sous l'action de l'oxygène de l'air et de l'eau. Problème majeur pour les infrastructures (ponts sur le Wouri, pipelines SONARA, armatures CIMENCAM). Protection : peinture, galvanisation (zinc, plus réducteur que le fer, s'oxyde à sa place), cathodique.
\item \textbf{Galvanisation à chaud :} Trempage de l'acier dans du zinc fondu (\SI{450}{\celsius}). Le zinc forme une barrière et une protection sacrificielle (anode de sacrifice).
\item \textbf{Piles et accumulateurs :} Principe de la pile Daniell (Zn/Cu) ou alcaline (Zn/MnO2) utilisé dans les lampes torches et radios rurales. La réaction d'oxydoréduction est spatialement séparée (anode/cathode) pour produire du courant électrique.
\item \textbf{Industrie chimique (SONARA, usines d'engrais) :} Procédés redox pour le raffinage (désulfuration), la production d'hydrogène (vaporeformage du méthane : \ce{CH4 + H2O -> CO + 3H2}, redox du carbone), synthèse d'ammoniac.
\item \textbf{Pharmacopée et alimentation :} Le bissap (\emph{Hibiscus sabdariffa}) contient des anthocyanes (indicateurs redox/pH naturels) et de l'acide ascorbique (vitamine C, puissant réducteur/antioxydant). Le vinaigre de palme (acide acétique) résulte de l'oxydation bactérienne de l'éthanol (vin de palme) : \ce{CH3CH2OH + O2 -> CH3COOH + H2O}.
\end{itemize}
\end{savaistu}

\begin{attention}[Les 5 erreurs qui coûtent des points au Probatoire]
\begin{enumerate}
\item \textbf{Faire apparaître les électrons dans l'équation-bilan.} Les électrons n'existent que dans les demi-équations : ils doivent obligatoirement se simplifier dans le bilan. Un « $+2e^-$ » dans l'équation finale est éliminatoire.
\item \textbf{Inverser oxydant et réducteur.} Le réducteur est celui qui \emph{perd} des électrons (et donc s'oxyde). Répète : « le réducteur réduit l'autre en s'oxydant lui-même ».
\item \textbf{Croire que tous les métaux réagissent avec les acides dilués.} Le cuivre, l'argent et l'or ne réagissent \textbf{pas} avec \ce{HCl} et \ce{H2SO4} dilués. Écrire \ce{Cu + 2H+ -> Cu^2+ + H2} est une erreur classique et grave.
\item \textbf{Oublier d'égaliser les électrons avant d'additionner.} Pour \ce{Al} (3 $e^-$) et \ce{Cu^2+} (2 $e^-$), il faut multiplier par 2 et par 3 (PPCM = 6). Additionner directement donne une équation fausse (charges non conservées).
\item \textbf{Négliger la vérification finale et le réactif limitant.} Toujours contrôler la conservation des atomes \emph{et} des charges. Dans les calculs, identifier le réactif limitant avec un tableau d'avancement avant de conclure (unités, chiffres significatifs, phrase de conclusion).
\end{enumerate}
\end{attention}

\newpage

\coursec{Exercices}

\courssub{Je vérifie mes connaissances}

\begin{exercice}[QCM : les bons réflexes]
Pour chaque question, choisis la (ou les) bonne(s) réponse(s).
\begin{enumerate}
\item Lors de l'action de l'acide chlorhydrique dilué sur le zinc, le gaz qui se dégage est :
\begin{enumerate}
\item le dioxygène ; \item le dihydrogène ; \item le dioxyde de carbone ; \item le dichlore.
\end{enumerate}
\item Un oxydant est une espèce chimique capable de :
\begin{enumerate}
\item céder des électrons ; \item capter des électrons ; \item céder des protons ; \item capter des protons.
\end{enumerate}
\item Dans la réaction \ce{Fe + Cu^2+ -> Fe^2+ + Cu}, l'espèce qui subit une oxydation est :
\begin{enumerate}
\item \ce{Fe} ; \item \ce{Cu^2+} ; \item \ce{Fe^2+} ; \item \ce{Cu}.
\end{enumerate}
\item Parmi les métaux suivants, lequel n'est \emph{pas} attaqué par l'acide chlorhydrique dilué ?
\begin{enumerate}
\item le fer ; \item le zinc ; \item le cuivre ; \item l'aluminium.
\end{enumerate}
\item La demi-équation \ce{Zn^2+ + 2e- -> Zn} correspond à :
\begin{enumerate}
\item une oxydation ; \item une réduction ; \item une réaction acido-basique ; \item une précipitation.
\end{enumerate}
\end{enumerate}
\end{exercice}

\begin{exercice}[Vrai ou faux ? Justifie]
Réponds par \textbf{vrai} ou \textbf{faux}, puis justifie ta réponse en une ou deux phrases.
\begin{enumerate}
\item Le cuivre est attaqué par l'acide sulfurique dilué avec dégagement de dihydrogène.
\item Lors d'une réduction, le nombre d'oxydation de l'élément diminue.
\item Une lame de zinc plongée dans une solution de sulfate de fer(II) se recouvre de fer.
\item Dans la réaction \ce{Zn + 2H+ -> Zn^2+ + H2}, l'ion \ce{H+} est le réducteur.
\item L'or et l'argent ne réagissent pas avec les acides chlorhydrique et sulfurique dilués.
\end{enumerate}
\end{exercice}

\begin{exercice}[Texte à trous]
Complète le texte suivant avec les mots ou expressions de la liste ci-dessous (certains peuvent être utilisés plusieurs fois, d'autres sont des intrus) : \emph{oxydation ; réduction ; oxydant ; réducteur ; électrons ; protons ; dihydrogène ; dioxygène ; augmente ; diminue}.

Une \ldots\ldots\ldots est une transformation au cours de laquelle une espèce chimique cède des \ldots\ldots\ldots ; l'espèce qui les cède est un \ldots\ldots\ldots. Une \ldots\ldots\ldots est une transformation au cours de laquelle une espèce chimique capte des \ldots\ldots\ldots ; l'espèce qui les capte est un \ldots\ldots\ldots. Lorsqu'on plonge une lame de fer dans une solution diluée d'acide chlorhydrique, on observe un dégagement de \ldots\ldots\ldots : le fer est \ldots\ldots\ldots (oxydé / réduit) et son nombre d'oxydation \ldots\ldots\ldots.
\end{exercice}

\begin{exercice}[Définitions et vocabulaire]
\begin{enumerate}
\item Définis les termes suivants : oxydation, réduction, oxydant, réducteur, réaction d'oxydoréduction.
\item Pour chacune des réactions ci-dessous, identifie l'oxydant et le réducteur, puis précise l'espèce oxydée et l'espèce réduite :
\[ \ce{Fe + Cu^2+ -> Fe^2+ + Cu} \qquad ; \qquad \ce{Zn + 2H+ -> Zn^2+ + H2} \]
\item Recopie et complète les demi-équations électroniques suivantes :
\[ \ce{Al -> Al^3+ + 3e-} \qquad ; \qquad \ce{Cu^2+ + 2e- -> Cu} \qquad ; \qquad \ce{Fe -> Fe^2+ + 2e-} \]
\end{enumerate}
\end{exercice}

\courssub{Je m'entraîne}

\begin{exercice}[Action des acides dilués sur les métaux]
On dispose de six métaux : zinc, fer, cuivre, aluminium, argent et or.
\begin{enumerate}
\item Pour chacun de ces métaux, indique s'il est attaqué par l'acide chlorhydrique dilué (à froid) et, le cas échéant, écris l'équation-bilan de la réaction.
\item Même question avec l'acide sulfurique dilué.
\item On plonge une grenaille de zinc dans une solution diluée d'acide chlorhydrique. Décris précisément les observations expérimentales et propose un test permettant d'identifier le gaz dégagé.
\end{enumerate}
\end{exercice}

\begin{exercice}[Demi-équations et équations-bilans]
Écris les demi-équations électroniques, puis l'équation-bilan de chacune des réactions suivantes, en précisant pour chacune l'oxydant et le réducteur :
\begin{enumerate}
\item action de l'acide chlorhydrique dilué sur le fer ;
\item action de l'acide sulfurique dilué sur l'aluminium ;
\item action d'une lame de zinc sur une solution contenant des ions \ce{Cu^2+} ;
\item action d'une lame de fer sur une solution contenant des ions \ce{Zn^2+} (la réaction a-t-elle lieu ? justifie).
\end{enumerate}
\end{exercice}

\begin{exercice}[Le bijou jaune du marché]
Un commerçant du marché central de Douala vend un bijou jaune qu'il affirme être en or pur. Pour vérifier, un élève de Première C plonge le bijou dans un bécher contenant de l'acide chlorhydrique dilué : il observe un dégagement gazeux.
\begin{enumerate}
\item Le bijou peut-il être en or pur ? Justifie ta réponse à l'aide des demi-équations électroniques.
\item En supposant que le bijou contient du zinc, écris l'équation-bilan de la réaction observée.
\item Propose un test simple permettant d'identifier le gaz dégagé.
\end{enumerate}
\end{exercice}

\begin{exercice}[Trois flacons non étiquetés]
Au laboratoire du lycée, trois flacons ont perdu leur étiquette. Ils contiennent des solutions aqueuses de sulfate de cuivre(II) (\ce{Cu^2+}), de sulfate de zinc (\ce{Zn^2+}) et de sulfate de fer(II) (\ce{Fe^2+}), toutes incolores ou presque. On dispose uniquement de lames de zinc, de fer et de cuivre.
\begin{enumerate}
\item Rappelle les couleurs caractéristiques des solutions contenant les ions \ce{Cu^2+} et \ce{Fe^2+}. Ce simple indice suffit-il à identifier les trois flacons ?
\item En utilisant les lames métalliques, propose un protocole expérimental précis permettant d'identifier le contenu de chaque flacon. Pour chaque test, indique les observations attendues et écris l'équation-bilan de la réaction éventuelle.
\end{enumerate}
\end{exercice}

\courssub{Je me prépare au Probatoire}

\begin{exercice}[La lame de zinc et le sulfate de cuivre]
On plonge une lame de zinc de masse $m = \SI{13,0}{g}$ dans un bécher contenant $V = \SI{200}{mL}$ d'une solution de sulfate de cuivre(II) de concentration $C = \SI{0,50}{mol.L^{-1}}$. On observe un dépôt rougeâtre sur la lame et la décoloration progressive de la solution.

\textbf{Données :} $M(\ce{Zn}) = \SI{65,4}{g.mol^{-1}}$ ; $M(\ce{Cu}) = \SI{63,5}{g.mol^{-1}}$.

\begin{enumerate}
\item Nomme et décris les observations expérimentales. Que signifie le dépôt rougeâtre ?
\item Écris les deux demi-équations électroniques mises en jeu, en précisant l'oxydant, le réducteur, l'espèce oxydée et l'espèce réduite.
\item Établis l'équation-bilan de la réaction.
\item Calcule les quantités de matière initiales des deux réactifs et détermine le réactif limitant à l'aide d'un tableau d'avancement.
\item Déduis-en :
\begin{enumerate}
\item la masse de cuivre déposée sur la lame ;
\item la masse de zinc restant à la fin de la réaction ;
\item la masse totale de la lame (zinc restant + cuivre déposé) en fin d'expérience.
\end{enumerate}
\end{enumerate}
\end{exercice}

\begin{exercice}[Attaque de l'aluminium par l'acide chlorhydrique]
Pour préparer du dihydrogène au laboratoire, on attaque une masse $m = \SI{2,7}{g}$ d'aluminium en grenaille par un excès d'acide chlorhydrique dilué (\ce{H3O+} + \ce{Cl-}). Le gaz dégagé est recueilli par déplacement d'eau dans une éprouvette graduée.

\textbf{Données :} $M(\ce{Al}) = \SI{27,0}{g.mol^{-1}}$ ; volume molaire dans les conditions de l'expérience : $V_m = \SI{22,4}{L.mol^{-1}}$.

\begin{enumerate}
\item L'aluminium est-il attaqué par l'acide chlorhydrique dilué ? Justifie.
\item Écris les demi-équations électroniques des couples \ce{Al^3+}/\ce{Al} et \ce{H3O+}/\ce{H2}, puis l'équation-bilan de la réaction (on fera intervenir \ce{H3O+}).
\item Précise l'oxydant et le réducteur de cette réaction.
\item Calcule la quantité de matière d'aluminium utilisée.
\item Construis le tableau d'avancement et détermine l'avancement maximal.
\item Déduis-en le volume de dihydrogène dégagé, mesuré dans les conditions de l'expérience.
\item Quel volume minimal de solution d'acide chlorhydrique de concentration $C = \SI{1,0}{mol.L^{-1}}$ aurait-il fallu utiliser pour attaquer exactement toute la grenaille ?
\end{enumerate}
\end{exercice}

\begin{exercice}[Prévision de réactions et identification d'un métal]
\textbf{Partie A : Prévision de réactions.}

On considère les métaux fer, zinc et cuivre, ainsi que des solutions contenant les ions \ce{Fe^2+}, \ce{Zn^2+} et \ce{H+}. On sait que le zinc réagit avec les ions \ce{Fe^2+} et avec les ions \ce{H+}, que le fer réagit avec les ions \ce{H+}, et que le cuivre ne réagit ni avec les ions \ce{H+} ni avec les ions \ce{Zn^2+}.

Pour chacune des expériences suivantes, dis en justifiant si une réaction d'oxydoréduction a lieu ; si oui, écris les demi-équations électroniques et l'équation-bilan :
\begin{enumerate}
\item une lame de fer plongée dans une solution contenant des ions \ce{Zn^2+} ;
\item une lame de zinc plongée dans une solution contenant des ions \ce{Fe^2+} ;
\item une lame de cuivre plongée dans une solution d'acide chlorhydrique dilué.
\end{enumerate}

\textbf{Partie B : Identification d'un métal inconnu.}

Un bijoutier de Yaoundé possède un clou métallique gris de masse $m = \SI{5,6}{g}$, dont il ignore la nature. Il le plonge dans un excès d'acide sulfurique dilué : le métal est entièrement attaqué et il se dégage un volume $V = \SI{2,24}{L}$ de dihydrogène, mesuré dans les CNTP. Le métal, noté M, donne des ions \ce{M^2+}.

\textbf{Données :} $V_m = \SI{22,4}{L.mol^{-1}}$ dans les CNTP ; $M(\ce{Fe}) = \SI{56}{g.mol^{-1}}$ ; $M(\ce{Zn}) = \SI{65,4}{g.mol^{-1}}$.

\begin{enumerate}
\setcounter{enumi}{3}
\item Écris l'équation-bilan de la réaction entre le métal M et l'acide sulfurique dilué.
\item Calcule la quantité de matière de dihydrogène dégagé, puis celle du métal M ayant réagi.
\item Déduis-en la masse molaire du métal M et identifie-le parmi le fer et le zinc.
\item Le bijoutier affirme que le clou aurait pu être en argent. Réfute cette hypothèse en une phrase argumentée.
\end{enumerate}
\end{exercice}

\newpage

\coursec{Corrigés des exercices}

\begin{corrige}[Exercice 1 --- QCM : les bons réflexes]
\begin{enumerate}
\item \textbf{Réponse b.} Le gaz qui se dégage est le \cle{dihydrogène} \ce{H2}. C'est le test à la flamme (petite détonation caractéristique, « aboiement ») qui permet de l'identifier. Le dioxygène \ce{O2} ravive une bûchette incandescente, le dioxyde de carbone \ce{CO2} trouble l'eau de chaux, le dichlore \ce{Cl2} est un gaz vert jaunâtre à odeur suffocante : aucun ne correspond à l'observation.
\item \textbf{Réponse b.} Un \cle{oxydant} est une espèce chimique capable de \emph{capter des électrons}. Céder des électrons est le rôle du réducteur ; céder ou capter des protons (\ce{H+}) relève des réactions acido-basiques (couples \ce{AH}/\ce{A-}).
\item \textbf{Réponse a.} Le fer \ce{Fe} perd deux électrons (\ce{Fe -> Fe^2+ + 2e-}) : il est \emph{oxydé}, c'est le réducteur. L'ion \ce{Cu^2+} capte ces deux électrons (\ce{Cu^2+ + 2e- -> Cu}) : il est \emph{réduit}, c'est l'oxydant.
\item \textbf{Réponse c.} Le \cle{cuivre} n'est pas attaqué par l'acide chlorhydrique dilué (ni par l'acide sulfurique dilué) : c'est un métal « noble » vis-à-vis des ions \ce{H+} (ou \ce{H3O+}), son pouvoir réducteur est plus faible que celui du dihydrogène. Le fer, le zinc et l'aluminium sont attaqués avec dégagement de \ce{H2}.
\item \textbf{Réponse b.} \ce{Zn^2+ + 2e- -> Zn} est un \emph{gain d'électrons} : c'est une \cle{réduction} (l'espèce \ce{Zn^2+} est l'oxydant, \ce{Zn} est le réducteur du couple \ce{Zn^2+}/\ce{Zn}).
\end{enumerate}
\end{corrige}

\begin{corrige}[Exercice 2 --- Vrai ou faux ?]
\begin{enumerate}
\item \textbf{Faux.} Le cuivre n'est \emph{pas} attaqué par l'acide sulfurique dilué : aucun dégagement de dihydrogène n'est observé. Seuls les métaux dont le pouvoir réducteur est plus fort que celui du dihydrogène (Zn, Fe, Al, Mg\dots) réagissent avec les acides dilués (équation ionique : \ce{M + 2H+ -> M^2+ + H2}).
\item \textbf{Vrai.} Une réduction est un gain d'électrons : l'espèce gagne des charges négatives, donc son nombre d'oxydation \emph{diminue}. Exemple : \ce{Cu^2+} (n.o. $= +2$) réduit en \ce{Cu} (n.o. $= 0$).
\item \textbf{Vrai.} Le zinc est plus réducteur que le fer : il cède ses électrons aux ions \ce{Fe^2+}, qui se réduisent en fer métallique se déposant sur la lame : \ce{Zn + Fe^2+ -> Zn^2+ + Fe}. Le zinc est oxydé, le fer(II) est réduit.
\item \textbf{Faux.} Dans \ce{Zn + 2H+ -> Zn^2+ + H2}, l'ion \ce{H+} \emph{captе} des électrons (il est réduit en \ce{H2}) : c'est l'\emph{oxydant}. Le réducteur est le zinc \ce{Zn}, qui cède ses électrons (il est oxydé en \ce{Zn^2+}).
\item \textbf{Vrai.} L'or et l'argent sont des métaux très peu réducteurs (très nobles) : ils ne sont attaqués ni par l'acide chlorhydrique dilué, ni par l'acide sulfurique dilué. C'est d'ailleurs une propriété utilisée pour tester les bijoux (test à l'acide nitrique dilué pour l'argent, eau régale pour l'or).
\end{enumerate}
\end{corrige}

\begin{corrige}[Exercice 3 --- Texte à trous]
Texte complété :

Une \textbf{oxydation} est une transformation au cours de laquelle une espèce chimique cède des \textbf{électrons} ; l'espèce qui les cède est un \textbf{réducteur}. Une \textbf{réduction} est une transformation au cours de laquelle une espèce chimique capte des \textbf{électrons} ; l'espèce qui les capte est un \textbf{oxydant}. Lorsqu'on plonge une lame de fer dans une solution diluée d'acide chlorhydrique, on observe un dégagement de \textbf{dihydrogène} : le fer est \textbf{oxydé} et son nombre d'oxydation \textbf{augmente} (de $0$ à $+2$).

\textbf{Intrus non utilisés :} \emph{protons} (les échanges de protons concernent les réactions acido-basiques, pas l'oxydoréduction) et \emph{dioxygène} (le gaz dégagé par l'action d'un acide dilué sur un métal est \ce{H2}, jamais \ce{O2}).
\end{corrige}

\begin{corrige}[Exercice 4 --- Définitions et vocabulaire]
\begin{enumerate}
\item \textbf{Définitions.}
\begin{itemize}
\item \cle{Oxydation} : transformation au cours de laquelle une espèce chimique \emph{cède} un ou plusieurs électrons.
\item \cle{Réduction} : transformation au cours de laquelle une espèce chimique \emph{captе} un ou plusieurs électrons.
\item \cle{Oxydant} : espèce chimique capable de \emph{capter} des électrons (il est réduit au cours de la réaction).
\item \cle{Réducteur} : espèce chimique capable de \emph{céder} des électrons (il est oxydé au cours de la réaction).
\item \cle{Réaction d'oxydoréduction} : réaction au cours de laquelle il y a \emph{transfert d'électrons} d'un réducteur vers un oxydant ; oxydation et réduction ont toujours lieu simultanément (pas d'oxydation sans réduction).
\end{itemize}

\item \textbf{Analyse des deux réactions.}

Pour \ce{Fe + Cu^2+ -> Fe^2+ + Cu} :
\begin{itemize}
\item \ce{Fe} perd 2 électrons : c'est le \textbf{réducteur}, il est \textbf{oxydé} (n.o. $0 \to +2$) ;
\item \ce{Cu^2+} gagne 2 électrons : c'est l'\textbf{oxydant}, il est \textbf{réduit} (n.o. $+2 \to 0$).
\end{itemize}

Pour \ce{Zn + 2H+ -> Zn^2+ + H2} :
\begin{itemize}
\item \ce{Zn} perd 2 électrons : c'est le \textbf{réducteur}, il est \textbf{oxydé} (n.o. $0 \to +2$) ;
\item \ce{H+} gagne des électrons : c'est l'\textbf{oxydant}, il est \textbf{réduit} (n.o. $+1 \to 0$).
\end{itemize}

\item \textbf{Demi-équations complétées.}
\[ \ce{Al -> Al^3+ + 3e-} \quad \text{(oxydation)} \qquad ; \qquad \ce{Cu^2+ + 2e- -> Cu} \quad \text{(réduction)} \]
\[ \ce{Fe -> Fe^2+ + 2e-} \quad \text{(oxydation)} \]
\textbf{Vérification systématique :} dans chaque demi-équation, les atomes \emph{et} les charges sont conservés. Par exemple pour l'aluminium : charge à gauche $= 0$, charge à droite $= (+3) + 3\times(-1) = 0$. Pour le cuivre : charge à gauche $= +2 + 2\times(-1) = 0$, charge à droite $= 0$.
\end{enumerate}
\end{corrige}

\begin{corrige}[Exercice 5 --- Action des acides dilués sur les métaux]
\begin{enumerate}
\item \textbf{Acide chlorhydrique dilué (\ce{H3O+} + \ce{Cl-}).}

\begin{center}
\begin{tabularx}{\linewidth}{|l|c|X|}
\hline
\rowcolor{popblueL} \textbf{Métal} & \textbf{Attaqué ?} & \textbf{Équation-bilan (ionique)} \\
\hline
Zinc \ce{Zn} & Oui & \ce{Zn + 2H3O+ -> Zn^2+ + H2 + 2H2O} \\
\hline
Fer \ce{Fe} & Oui & \ce{Fe + 2H3O+ -> Fe^2+ + H2 + 2H2O} \\
\hline
Aluminium \ce{Al} & Oui & \ce{2Al + 6H3O+ -> 2Al^3+ + 3H2 + 6H2O} \\
\hline
Cuivre \ce{Cu} & Non & aucune réaction \\
\hline
Argent \ce{Ag} & Non & aucune réaction \\
\hline
Or \ce{Au} & Non & aucune réaction \\
\hline
\end{tabularx}
\end{center}

\item \textbf{Acide sulfurique dilué (\ce{2H3O+} + \ce{SO4^2-}).} Les conclusions sont \emph{identiques} : l'espèce oxydante est dans les deux cas l'ion \ce{H3O+} (souvent noté simplement \ce{H+} dans les demi-équations rédox). Le zinc, le fer et l'aluminium sont attaqués avec dégagement de \ce{H2} ; le cuivre, l'argent et l'or ne réagissent pas. Les équations-bilans ioniques sont les mêmes qu'à la question 1 (les anions \ce{Cl-} ou \ce{SO4^2-} sont spectateurs).

\item \textbf{Observations et test du gaz.}

\emph{Observations :} on observe une \emph{effervescence} (dégagement gazeux) à la surface de la grenaille de zinc ; la grenaille \emph{diminue progressivement de volume} jusqu'à disparaître si l'acide est en excès ; le bécher \emph{chauffe légèrement} (réaction exothermique) ; la solution obtenue est \emph{incolore} (ions \ce{Zn^2+} incolores, ions \ce{Cl-} incolores).

\emph{Test d'identification du gaz :} on recueille le gaz dans un tube à essais retourné (par déplacement d'air), puis on approche une \emph{bûchette enflammée} de l'ouverture du tube : on entend une \emph{petite détonation} caractéristique (aboiement). Le gaz est le \cle{dihydrogène} \ce{H2}.
\end{enumerate}
\end{corrige}

\begin{corrige}[Exercice 6 --- Demi-équations et équations-bilans]
\begin{enumerate}
\item \textbf{Acide chlorhydrique dilué sur le fer.}
\begin{align*}
\text{Oxydation (réducteur : } \ce{Fe}\text{) : } & \ce{Fe -> Fe^2+ + 2e-} \\
\text{Réduction (oxydant : } \ce{H+}\text{) : } & \ce{2H+ + 2e- -> H2}
\end{align*}
\[ \ce{Fe + 2H+ -> Fe^2+ + H2} \quad \text{(équation ionique, } \ce{H+} \equiv \ce{H3O+}\text{)} \]

\item \textbf{Acide sulfurique dilué sur l'aluminium.}
\begin{align*}
\text{Oxydation (réducteur : } \ce{Al}\text{) : } & \ce{Al -> Al^3+ + 3e-} \quad (\times 2) \\
\text{Réduction (oxydant : } \ce{H+}\text{) : } & \ce{2H+ + 2e- -> H2} \quad (\times 3)
\end{align*}
On multiplie pour égaliser les électrons échangés ($6\,\ce{e-}$) :
\[ \ce{2Al + 6H+ -> 2Al^3+ + 3H2} \quad \text{(équation ionique)} \]
\emph{Remarque :} en solution aqueuse, on écrira souvent l'équation complète \ce{2Al + 6H3O+ -> 2Al^3+ + 3H2 + 6H2O}.

\item \textbf{Lame de zinc dans une solution d'ions \ce{Cu^2+}.}
\begin{align*}
\text{Oxydation (réducteur : } \ce{Zn}\text{) : } & \ce{Zn -> Zn^2+ + 2e-} \\
\text{Réduction (oxydant : } \ce{Cu^2+}\text{) : } & \ce{Cu^2+ + 2e- -> Cu}
\end{align*}
\[ \ce{Zn + Cu^2+ -> Zn^2+ + Cu} \]

\item \textbf{Lame de fer dans une solution d'ions \ce{Zn^2+}.}

\emph{La réaction n'a pas lieu.} Le fer est un réducteur \emph{moins fort} que le zinc (couple \ce{Fe^2+}/\ce{Fe} moins réducteur que \ce{Zn^2+}/\ce{Zn}) : il ne peut pas céder ses électrons aux ions \ce{Zn^2+}. Autrement dit, la réaction inverse (\ce{Zn + Fe^2+ -> Zn^2+ + Fe}) est celle qui se produit spontanément. Une lame de fer plongée dans une solution de sulfate de zinc reste intacte, la solution reste incolore.
\end{enumerate}
\end{corrige}

\begin{corrige}[Exercice 7 --- Le bijou jaune du marché]
\begin{enumerate}
\item \textbf{Le bijou ne peut pas être en or pur.} L'or est un métal très peu réducteur (très noble) : il n'est pas attaqué par l'acide chlorhydrique dilué. Si le bijou était en or pur, \emph{aucun} dégagement gazeux ne serait observé. Or on observe une effervescence : le bijou contient donc au moins un métal attaquable par l'acide (zinc, fer, aluminium, étain\dots). C'est un bijou « en or » falsifié (plaqué or sur un métal commun, ou alliage).

Formellement, la demi-équation de réduction de l'ion hydrogène \ce{2H+ + 2e- -> H2} ne peut pas être couplée à une oxydation de l'or (\ce{Au -> Au^3+ + 3e-}) dans ces conditions car le potentiel du couple \ce{Au^3+}/\ce{Au} est trop élevé (or moins réducteur que \ce{H2}).

\item \textbf{Si le bijou contient du zinc :}
\begin{align*}
\text{Oxydation : } & \ce{Zn -> Zn^2+ + 2e-} \\
\text{Réduction : } & \ce{2H+ + 2e- -> H2}
\end{align*}
\[ \ce{Zn + 2H+ -> Zn^2+ + H2} \]

\item \textbf{Test du gaz :} on recueille le gaz dans un tube retourné et on approche une flamme : une \emph{petite détonation} (aboiement) caractérise le dihydrogène \ce{H2}.
\end{enumerate}

\textbf{Point de vigilance :} ce test prouve que le bijou \emph{n'est pas} en or pur, mais ne prouve pas qu'il contient du zinc : tout métal attaquable par l'acide (fer, aluminium, étain\dots) donnerait la même observation. L'absence de réaction, en revanche, aurait été un indice fort en faveur de l'or (ou d'un métal noble comme le platine).
\end{corrige}

\begin{corrige}[Exercice 8 --- Trois flacons non étiquetés]
\begin{enumerate}
\item \textbf{Couleurs caractéristiques.}
\begin{itemize}
\item Les solutions contenant les ions \ce{Cu^2+} sont \emph{bleues} (bleu ciel à bleu soutenu) ;
\item les solutions contenant les ions \ce{Fe^2+} sont \emph{verdâtres} (vert très pâle) ;
\item les solutions contenant les ions \ce{Zn^2+} sont \emph{incolores}.
\end{itemize}
L'énoncé précise que les trois solutions sont « incolores ou presque » : l'indice de la couleur est donc \emph{insuffisant} (solutions trop diluées pour être discriminantes). Il faut des tests chimiques de réactivité.

\item \textbf{Protocole d'identification.}

On prélève un échantillon de chaque solution dans trois tubes à essais et on y plonge des lames métalliques. Règle utilisée : un métal réduit les ions d'un métal \emph{moins réducteur} que lui. Ordre des pouvoirs réducteurs (donné ou connu) : $\ce{Zn} > \ce{Fe} > \ce{Cu}$.

\textbf{Test 1 : lame de fer dans chaque solution.}
\begin{itemize}
\item Si un \emph{dépôt rougeâtre} de cuivre apparaît sur la lame et que la solution bleue se décolore : la solution contient \ce{Cu^2+}.
\[ \ce{Fe + Cu^2+ -> Fe^2+ + Cu} \]
\item Si la lame reste \emph{intacte} (pas de dépôt, pas de changement de couleur) : la solution contient \ce{Zn^2+} ou \ce{Fe^2+} (le fer ne réagit ni avec \ce{Zn^2+} --- zinc plus réducteur --- ni avec ses propres ions \ce{Fe^2+}).
\end{itemize}

\textbf{Test 2 : lame de zinc dans les deux solutions restantes (celles qui n'ont pas réagi avec le fer).}
\begin{itemize}
\item Si un \emph{dépôt gris-noirâtre} de fer apparaît sur la lame de zinc : la solution contient \ce{Fe^2+}.
\[ \ce{Zn + Fe^2+ -> Zn^2+ + Fe} \]
\item Si \emph{rien} ne se produit : la solution contient \ce{Zn^2+} (le zinc ne peut pas réduire ses propres ions).
\end{itemize}

\textbf{Bilan :} le flacon donnant un dépôt rougeâtre avec le fer contient \ce{Cu^2+} ; celui donnant un dépôt gris avec le zinc contient \ce{Fe^2+} ; celui ne réagissant avec aucune lame contient \ce{Zn^2+}.

\emph{Remarque :} la lame de cuivre n'est pas indispensable ici, mais elle peut servir de témoin : elle ne réagit avec \emph{aucune} des trois solutions, car le cuivre est le moins réducteur des trois métaux.
\end{enumerate}
\end{corrige}

\begin{corrige}[Exercice 9 --- La lame de zinc et le sulfate de cuivre]
\begin{enumerate}
\item \textbf{Observations.} La solution de sulfate de cuivre(II), initialement \emph{bleue}, se \emph{décolore progressivement} ; un \emph{dépôt rougeâtre} apparaît à la surface de la lame de zinc. Ce dépôt est du \emph{cuivre métallique} \ce{Cu} formé par réduction des ions \ce{Cu^2+} ; la décoloration traduit la disparition des ions \ce{Cu^2+} de la solution. La lame de zinc se couvre de cuivre et s'amincit.

\item \textbf{Demi-équations.}
\begin{align*}
\text{Oxydation : } & \ce{Zn -> Zn^2+ + 2e-} \quad \text{(\ce{Zn} est le réducteur, il est oxydé)} \\
\text{Réduction : } & \ce{Cu^2+ + 2e- -> Cu} \quad \text{(\ce{Cu^2+} est l'oxydant, il est réduit)}
\end{align*}

\item \textbf{Équation-bilan.}
\[ \ce{Zn + Cu^2+ -> Zn^2+ + Cu} \]

\item \textbf{Quantités initiales et réactif limitant.}
\[ n_i(\ce{Zn}) = \frac{m}{M(\ce{Zn})} = \frac{13{,}0}{65{,}4} \approx \SI{0,199}{mol} \]
\[ n_i(\ce{Cu^2+}) = C \times V = 0{,}50 \times 0{,}200 = \SI{0,100}{mol} \]

\begin{center}
\begin{tabularx}{\linewidth}{|l|c|X|X|X|X|}
\hline
\rowcolor{popblueL} \textbf{État} & \textbf{Avanc.} & \ce{Zn} & \ce{Cu^2+} & \ce{Zn^2+} & \ce{Cu} \\
\hline
Initial & $0$ & $0{,}199$ & $0{,}100$ & $0$ & $0$ \\
\hline
En cours & $x$ & $0{,}199 - x$ & $0{,}100 - x$ & $x$ & $x$ \\
\hline
Final & $x_{max}$ & $0{,}199 - x_{max}$ & $0{,}100 - x_{max}$ & $x_{max}$ & $x_{max}$ \\
\hline
\end{tabularx}
\end{center}

Hypothèses : si \ce{Zn} est limitant, $x_{max} = \SI{0,199}{mol}$ ; si \ce{Cu^2+} est limitant, $x_{max} = \SI{0,100}{mol}$. La plus petite valeur l'emporte :
\[ x_{max} = \SI{0,100}{mol} \qquad \Rightarrow \qquad \textbf{\ce{Cu^2+} est le réactif limitant.} \]

\item \textbf{Bilans de masse.}
\begin{enumerate}
\item Masse de cuivre déposée : $n(\ce{Cu}) = x_{max} = \SI{0,100}{mol}$, donc
\[ m(\ce{Cu}) = 0{,}100 \times 63{,}5 = \SI{6,35}{g} \]
\item Zinc consommé : $0{,}100 \times 65{,}4 = \SI{6,54}{g}$. Zinc restant :
\[ m_{rest}(\ce{Zn}) = 13{,}0 - 6{,}54 = \SI{6,46}{g} \]
\item Masse totale de la lame en fin d'expérience :
\[ m_{tot} = 6{,}46 + 6{,}35 = \SI{12,81}{g} \]
\end{enumerate}
\emph{Remarque :} la lame s'est légèrement \emph{allégée} car la masse molaire du zinc ($65{,}4$) est supérieure à celle du cuivre ($63{,}5$) : atome pour atome, le zinc qui part « pèse » plus lourd que le cuivre qui se dépose.
\end{enumerate}

\textbf{Barème indicatif (sur 10) :} observations 1 pt ; demi-équations 2 pts ; équation-bilan 1 pt ; quantités initiales 2 pts ; tableau d'avancement et réactif limitant 2 pts ; masses finales 2 pts.

\textbf{Points de vigilance :} convertir le volume en litres ($V = \SI{0,200}{L}$) ; ne pas oublier que le zinc est en excès (il \emph{reste} du zinc à la fin) ; vérifier la cohérence : la masse totale de la lame doit être proche de la masse initiale.
\end{corrige}

\begin{corrige}[Exercice 10 --- Attaque de l'aluminium par l'acide chlorhydrique]
\begin{enumerate}
\item \textbf{Oui}, l'aluminium est attaqué par l'acide chlorhydrique dilué : c'est un métal suffisamment réducteur pour réduire les ions \ce{H3O+} en dihydrogène. On observe une effervescence et la disparition progressive de la grenaille.

\item \textbf{Demi-équations et équation-bilan.}
On utilise le couple \ce{H3O+}/\ce{H2} (souvent noté \ce{H+}/\ce{H2}) pour la réduction.
\begin{align*}
\text{Couple } \ce{Al^3+}/\ce{Al} \text{ (oxydation) : } & \ce{Al -> Al^3+ + 3e-} \quad (\times 2) \\
\text{Couple } \ce{H3O+}/\ce{H2} \text{ (réduction) : } & \ce{2H3O+ + 2e- -> H2 + 2H2O} \quad (\times 3)
\end{align*}
En égalisant les électrons échangés ($6\,\ce{e-}$) :
\[ \ce{2Al + 6H3O+ -> 2Al^3+ + 3H2 + 6H2O} \]
\emph{Équation ionique simplifiée (avec \ce{H+}) :} \ce{2Al + 6H+ -> 2Al^3+ + 3H2}.

\item \textbf{Oxydant et réducteur.} Le réducteur est l'aluminium \ce{Al} (il cède des électrons, il est oxydé) ; l'oxydant est l'ion oxonium \ce{H3O+} (il capte des électrons, il est réduit en \ce{H2}).

\item \textbf{Quantité de matière d'aluminium.}
\[ n(\ce{Al}) = \frac{m}{M(\ce{Al})} = \frac{2{,}7}{27{,}0} = \SI{0,100}{mol} \]

\item \textbf{Tableau d'avancement} (l'acide est en excès, l'aluminium est limitant) :

\begin{center}
\begin{tabularx}{\linewidth}{|l|c|X|X|X|}
\hline
\rowcolor{popblueL} \textbf{État} & \textbf{Avanc.} & \ce{2Al} & \ce{6H3O+} & \ce{3H2} \\
\hline
Initial & $0$ & $0{,}100$ & excès & $0$ \\
\hline
En cours & $x$ & $0{,}100 - 2x$ & excès & $3x$ \\
\hline
Final & $x_{max}$ & $0{,}100 - 2x_{max}$ & excès & $3x_{max}$ \\
\hline
\end{tabularx}
\end{center}

Aluminium limitant : $0{,}100 - 2x_{max} = 0 \Rightarrow x_{max} = \SI{0,0500}{mol}$.

\item \textbf{Volume de dihydrogène dégagé.}
\[ n(\ce{H2}) = 3x_{max} = 3 \times 0{,}0500 = \SI{0,150}{mol} \]
\[ V(\ce{H2}) = n(\ce{H2}) \times V_m = 0{,}150 \times 22{,}4 = \SI{3,36}{L} \quad \text{(aux CNTP)} \]

\item \textbf{Volume minimal d'acide.} Pour attaquer exactement tout l'aluminium, il faut :
\[ n(\ce{H3O+}) = 3 \times n(\ce{Al}) = 3 \times 0{,}100 = \SI{0,300}{mol} \]
(d'après l'équation-bilan : $6$ mol d'acide pour $2$ mol d'aluminium, soit un rapport de $3$).
\[ V_{min} = \frac{n(\ce{H3O+})}{C} = \frac{0{,}300}{1{,}0} = \SI{0,300}{L} = \SI{300}{mL} \]
\end{enumerate}

\textbf{Barème indicatif (sur 10) :} question 1 : 1 pt ; demi-équations et bilan : 3 pts ; oxydant/réducteur : 1 pt ; $n(\ce{Al})$ : 1 pt ; tableau et $x_{max}$ : 2 pts ; $V(\ce{H2})$ : 1 pt ; $V_{min}$ : 1 pt.

\textbf{Points de vigilance :} bien équilibrer la demi-équation du couple \ce{H3O+}/\ce{H2} avec \ce{H2O} du côté des produits ; attention aux coefficients stoechiométriques : $n(\ce{H2}) = \frac{3}{2}\,n(\ce{Al})$ et non $n(\ce{Al})$ ; ne pas confondre $x_{max}$ ($\SI{0,0500}{mol}$) avec $n(\ce{H2})$ ($\SI{0,150}{mol}$).
\end{corrige}

\begin{corrige}[Exercice 11 --- Prévision de réactions et identification d'un métal]
\textbf{Partie A : Prévision de réactions.}

Données de l'énoncé : le zinc réagit avec \ce{Fe^2+} et \ce{H+} ; le fer réagit avec \ce{H+} ; le cuivre ne réagit ni avec \ce{H+} ni avec \ce{Zn^2+}. Le pouvoir réducteur s'ordonne donc : $\ce{Zn} > \ce{Fe} > \ce{H2} > \ce{Cu}$.

\begin{enumerate}
\item \textbf{Lame de fer dans une solution de \ce{Zn^2+} : pas de réaction.} Le fer est un réducteur plus faible que le zinc : il ne peut pas céder ses électrons aux ions \ce{Zn^2+}. La lame reste intacte, la solution reste incolore.

\item \textbf{Lame de zinc dans une solution de \ce{Fe^2+} : réaction.} Le zinc, plus réducteur que le fer, cède ses électrons aux ions \ce{Fe^2+} :
\begin{align*}
\text{Oxydation : } & \ce{Zn -> Zn^2+ + 2e-} \\
\text{Réduction : } & \ce{Fe^2+ + 2e- -> Fe}
\end{align*}
\[ \ce{Zn + Fe^2+ -> Zn^2+ + Fe} \]
Observation attendue : dépôt gris-noirâtre de fer sur la lame de zinc, la solution reste incolore.

\item \textbf{Lame de cuivre dans l'acide chlorhydrique dilué : pas de réaction.} Le cuivre est moins réducteur que \ce{H2} (donnée de l'énoncé : il ne réagit pas avec \ce{H+}) : aucun dégagement gazeux n'est observé, la lame ne se dissout pas.
\end{enumerate}

\textbf{Partie B : Identification d'un métal inconnu.}

\begin{enumerate}
\setcounter{enumi}{3}
\item \textbf{Équation-bilan.} Le métal M donne des ions \ce{M^2+} et l'acide sulfurique dilué fournit les ions \ce{H+} :
\begin{align*}
\text{Oxydation : } & \ce{M -> M^2+ + 2e-} \\
\text{Réduction : } & \ce{2H+ + 2e- -> H2}
\end{align*}
\[ \ce{M + 2H+ -> M^2+ + H2} \]

\item \textbf{Quantités de matière.}
\[ n(\ce{H2}) = \frac{V}{V_m} = \frac{2{,}24}{22{,}4} = \SI{0,100}{mol} \]
D'après l'équation-bilan, $1$ mol de métal produit $1$ mol de \ce{H2} (coefficients stœchiométriques identiques) :
\[ n(\ce{M}) = n(\ce{H2}) = \SI{0,100}{mol} \]

\item \textbf{Masse molaire et identification.}
\[ M(\ce{M}) = \frac{m}{n(\ce{M})} = \frac{5{,}6}{0{,}100} = \SI{56}{g.mol^{-1}} \]
Cette valeur correspond exactement à celle du \textbf{fer} ($M(\ce{Fe}) = \SI{56}{g.mol^{-1}}$) et non à celle du zinc ($\SI{65,4}{g.mol^{-1}}$). Le clou est donc en \emph{fer}.

\item \textbf{Réfutation de l'hypothèse « argent ».} L'argent n'est \emph{pas} attaqué par l'acide sulfurique dilué (c'est un métal noble vis-à-vis de \ce{H+}, $E^\circ(\ce{Ag+}/\ce{Ag}) > E^\circ(\ce{H+}/\ce{H2})$) : or le clou a été \emph{entièrement} attaqué avec dégagement de dihydrogène ; il ne peut donc pas être en argent.
\end{enumerate}

\textbf{Barème indicatif (sur 10) :} Partie A : 3 pts (1 pt par situation, justification comprise) ; équation-bilan : 1 pt ; $n(\ce{H2})$ : 1 pt ; relation stoechiométrique et $n(\ce{M})$ : 1 pt ; $M(\ce{M})$ et identification : 2 pts ; réfutation argent : 1 pt ; rigueur (unités, chiffres significatifs) : 1 pt.

\textbf{Points de vigilance :} la relation $n(\ce{M}) = n(\ce{H2})$ n'est valable que parce que M donne des ions \ce{M^2+} (coefficients 1 pour 1) ; bien préciser que le volume est mesuré dans les CNTP, ce qui justifie l'usage de $V_m = \SI{22,4}{L.mol^{-1}}$ ; pour identifier le métal, comparer explicitement la masse molaire calculée aux deux valeurs proposées.
\end{corrige}

\newpage

\coursec{Fiche bilan}

\courssub{Synthèse et points-clés de la leçon}

\begin{popfigure}
\begin{center}
\begin{tikzpicture}[
  mindmap/.style={font=\small\bfseries, text=white, align=center},
  central/.style={circle, fill=popink, minimum size=2.8cm, mindmap},
  branche/.style={rounded corners=4pt, align=center, font=\footnotesize\bfseries, text=white, inner sep=5pt},
  lien/.style={thick, draw=popmuted!60, line width=1.2pt, shorten >=2pt, shorten <=2pt}
]
\node[central] (C) {Oxydo-\\réduction\\en solution\\aqueuse};
\node[branche, fill=popblue, anchor=west] (A) at ($(C)+(35:3.8)$) {Acide + métal\\ $\ce{Zn + 2H+ -> Zn^2+ + H2 ^}$};
\node[branche, fill=poporange, anchor=west] (B) at ($(C)+(-15:4.2)$) {Ion métallique + métal\\ $\ce{Zn + Cu^2+ -> Zn^2+ + Cu}$};
\node[branche, fill=popgreen, anchor=east] (D) at ($(C)+(165:4.2)$) {Définitions\\ Oxydation = perte $e^-$\\ Réduction = gain $e^-$};
\node[branche, fill=poppurple, anchor=east] (E) at ($(C)+(-145:4.0)$) {Oxydant / Réducteur\\ Oxydant capte $e^-$\\ Réducteur cède $e^-$};
\node[branche, fill=poppink, anchor=south] (F) at ($(C)+(-95:3.5)$) {Demi-équations\\ $\ce{Zn -> Zn^2+ + 2e^-}$\\ $\ce{Cu^2+ + 2e^- -> Cu}$};
\node[branche, fill=popgold, anchor=south] (G) at ($(C)+(95:3.7)$) {Prévision réaction\\ Métal + ion d'un métal\\ \og moins réducteur \fg};
\draw[lien] (C) -- (A); \draw[lien] (C) -- (B); \draw[lien] (C) -- (D);
\draw[lien] (C) -- (E); \draw[lien] (C) -- (F); \draw[lien] (C) -- (G);
\end{tikzpicture}
\end{center}
\legende{Carte mentale de la leçon : les six idées-forces de l'oxydoréduction en solution aqueuse.}
\end{popfigure}

\begin{bilan}[L'essentiel de la leçon]
\textbf{Définitions clés (à savoir réciter sans ambiguïté)}
\begin{itemize}
  \item \cle{Oxydation} : transformation au cours de laquelle une espèce chimique \textbf{perd un ou plusieurs électrons}.
  \item \cle{Réduction} : transformation au cours de laquelle une espèce chimique \textbf{gagne un ou plusieurs électrons}.
  \item \cle{Oxydant} : espèce chimique capable de \textbf{capter} des électrons ; il se trouve \textbf{réduit} au cours de la réaction.
  \item \cle{Réducteur} : espèce chimique capable de \textbf{céder} des électrons ; il se trouve \textbf{oxydé} au cours de la réaction.
  \item \cle{Oxydoréduction} : réaction chimique mettant en jeu un \textbf{transfert d'électrons} du réducteur vers l'oxydant.
  \item \cle{Couple oxydant/réducteur} : ensemble de deux espèces conjuguées, noté $\ce{Ox/Red}$, reliées par la demi-équation $\ce{Ox + n e^- <=> Red}$.
\end{itemize}

\textbf{Équations-types à connaître par cœur (avec états physiques)}
\begin{itemize}
  \item Action d'un acide dilué (\ce{H+}) sur un métal réducteur :
  \[ \ce{Zn(s) + 2H+(aq) -> Zn^2+(aq) + H2(g) ^} \qquad \ce{Fe(s) + 2H+(aq) -> Fe^2+(aq) + H2(g) ^} \]
  \item Réaction de \og cimentation \fg (métal + ion métallique) :
  \[ \ce{Zn(s) + Cu^2+(aq) -> Zn^2+(aq) + Cu(s)} \qquad \ce{Fe(s) + Cu^2+(aq) -> Fe^2+(aq) + Cu(s)} \]
  \item Métaux \textbf{nobles} (\ce{Cu}, \ce{Ag}, \ce{Au}) : \textbf{pas de réaction} avec \ce{HCl} ou \ce{H2SO4} \textbf{dilués} (non oxydants).
\end{itemize}

\textbf{Tableau récapitulatif des expériences de référence}
\begin{center}
\begin{tabularx}{\linewidth}{|l|X|X|}
\hline
\rowcolor{popblueL} \textbf{Système} & \textbf{Observation} & \textbf{Conclusion (Équation-bilan)} \\
\hline
Zinc + \ce{HCl} dilué & Dégagement gazeux vif (bulles), dissolution du solide, le gaz détonne à la flamme & $\ce{Zn(s) + 2H+(aq) -> Zn^2+(aq) + H2(g) ^}$ \\
\hline
Fer + \ce{HCl} dilué & Dégagement gazeux plus lent, dissolution du solide, le gaz détonne à la flamme & $\ce{Fe(s) + 2H+(aq) -> Fe^2+(aq) + H2(g) ^}$ \\
\hline
Cuivre, Argent, Or + \ce{HCl} dilué & Aucune réaction visible, pas de dégagement gazeux & Métaux nobles : \ce{Cu}, \ce{Ag}, \ce{Au} non oxydés par \ce{H+} \\
\hline
Fer + solution de \ce{CuSO4} & Dépôt rouge-brun solide (cuivre métallique), décoloration progressive de la solution bleue & $\ce{Fe(s) + Cu^2+(aq) -> Fe^2+(aq) + Cu(s)}$ \\
\hline
Cuivre + solution de \ce{ZnSO4} & Aucune réaction visible, la solution reste incolore & \ce{Cu} moins réducteur que \ce{Zn} : pas de déplacement \\
\hline
\end{tabularx}
\end{center}

\textbf{Méthode express : écrire une équation-bilan d'oxydoréduction}
\begin{enumerate}
  \item Identifier l'oxydant (qui capte les $e^-$) et le réducteur (qui cède les $e^-$).
  \item Écrire les deux demi-équations électroniques : \textbf{oxydation} $\mathrm{Red} \rightarrow \mathrm{Ox} + n\,\ce{e^-}$ et \textbf{réduction} $\mathrm{Ox'} + n'\,\ce{e^-} \rightarrow \mathrm{Red'}$.
  \item Multiplier les demi-équations par les coefficients adéquats pour \textbf{égaliser le nombre d'électrons} échangés.
  \item Additionner les deux demi-équations et \textbf{simplifier les électrons} (ils ne figurent \textbf{jamais} dans l'équation-bilan finale).
  \item Vérifier la conservation des \textbf{éléments} (atomes) et des \textbf{charges} électriques.
\end{enumerate}

\textbf{Méthode express : prévoir si une réaction d'oxydoréduction a lieu}
\begin{itemize}
  \item \textbf{Métal + \ce{H+} (acide dilué) :} réaction seulement si le métal est \textbf{plus réducteur que le dihydrogène} (\ce{H2}).
  \begin{itemize}
    \item \textbf{Oui :} \ce{Al}, \ce{Zn}, \ce{Fe} (et métaux alcalins/alcalino-terreux).
    \item \textbf{Non :} \ce{Cu}, \ce{Ag}, \ce{Au} (métaux nobles).
  \end{itemize}
  \item \textbf{Métal $M$ + Ion métallique $M'^{n+}$ :} réaction (déplacement) seulement si $M$ est \textbf{plus réducteur} que $M'$.
  \begin{itemize}
    \item Exemple : \ce{Zn} déplace \ce{Cu^2+} ; \ce{Fe} déplace \ce{Cu^2+} ; \ce{Cu} ne déplace \textbf{pas} \ce{Zn^2+}.
  \end{itemize}
  \item \textbf{Ordre de pouvoir réducteur (force décroissante) à retenir :}
  \[ \ce{Al} > \ce{Zn} > \ce{Fe} > \ce{H2} > \ce{Cu} > \ce{Ag} > \ce{Au} \]
  \emph{(Un réducteur plus fort que \ce{H2} libère \ce{H2} avec un acide dilué ; un métal plus fort qu'un autre déplace son ion).}
\end{itemize}
\end{bilan}

\begin{attention}[Pièges classiques et points de vigilance pour le Probatoire]
\begin{itemize}
  \item \textbf{Confusion Oxydant / Réducteur :} L'oxydant \textbf{gagne} des électrons (il est réduit) ; le réducteur \textbf{perd} des électrons (il est oxydé). Mnémotechnique : \og \textbf{Ox}ydant = \textbf{Ox}ydation ? \textbf{NON !} \textbf{Ox}ydant = \textbf{Ré}duction \fg.
  \item \textbf{Électrons dans le bilan :} \textbf{Interdit} d'écrire des $\ce{e^-}$ dans l'équation-bilan finale. Ils s'annulent obligatoirement.
  \item \textbf{Aluminium et passivation :} L'aluminium est théoriquement un fort réducteur ($Al > Zn$), mais sa surface est recouverte d'une couche d'alumine $\ce{Al2O3}$ protectrice. Il ne réagit \textbf{pas immédiatement} avec un acide dilué (temps de latence = dépassivation). Ne pas écrire \og pas de réaction \fg, écrire \og réaction lente après dépassivation \fg.
  \item \textbf{Acide sulfurique concentré vs dilué :} $\ce{H2SO4}$ dilué $\approx$ source de $\ce{H+}$ (comme $\ce{HCl}$). $\ce{H2SO4}$ \textbf{concentré} est un \textbf{oxydant fort} (ion $\ce{SO4^2-}$) : il oxyde $\ce{Cu}$, $\ce{Ag}$... en donnant $\ce{SO2}$ (gaz toxique), pas $\ce{H2}$. Le programme porte sur les acides \textbf{dilués}.
  \item \textbf{Équilibrage des charges :} Vérifier systématiquement que la somme des charges est identique des deux côtés de la flèche dans chaque demi-équation et dans le bilan.
  \item \textbf{États physiques :} Au Probatoire, les états ($\ce{(s)}$, $\ce{(aq)}$, $\ce{(g)}$, $\ce{(l)}$) sont \textbf{obligatoires} dans les équations-bilan.
  \item \textbf{Caractérisation du dihydrogène :} \og Détonation à la flamme \fg (bruit sec), pas \og flamme bleue \fg (c'est la combustion), pas \og eau trouble \fg (c'est $\ce{CO2}$).
\end{itemize}
\end{attention}

\begin{aretenir}[Checklist finale avant le Probatoire]
\begin{itemize}
  \item $\square$ Je sais définir \textbf{oxydation}, \textbf{réduction}, \textbf{oxydant}, \textbf{réducteur}, \textbf{oxydoréduction} et \textbf{couple Ox/Red} avec les termes exacts (perte/gain/capte/cède/transfert).
  \item $\square$ Je sais qu'une oxydation est une \textbf{perte d'électrons} (numéro d'oxydation augmente) et une réduction un \textbf{gain d'électrons} (numéro d'oxydation diminue).
  \item $\square$ Je sais écrire et équilibrer une demi-équation électronique (atomes + charges) : ex. $\ce{Zn -> Zn^2+ + 2e^-}$, $\ce{Cu^2+ + 2e^- -> Cu}$, $\ce{2H+ + 2e^- -> H2}$.
  \item $\square$ Je sais combiner deux demi-équations pour obtenir l'équation-bilan \textbf{sans électrons}, avec les états physiques.
  \item $\square$ Je connais par cœur les équations-bilan : $\ce{Zn + 2H+}$, $\ce{Fe + 2H+}$, $\ce{Zn + Cu^2+}$, $\ce{Fe + Cu^2+}$.
  \item $\square$ Je sais que $\ce{Cu}$, $\ce{Ag}$, $\ce{Au}$ ne réagissent pas avec $\ce{HCl}$ ou $\ce{H2SO4}$ \textbf{dilués}.
  \item $\square$ Je sais décrire l'expérience \og Fer + sulfate de cuivre \fg : dépôt rouge ($\ce{Cu}$), décoloration bleue ($\ce{Cu^2+}$), et écrire son bilan.
  \item $\square$ Je sais identifier sans hésiter l'oxydant et le réducteur dans une équation-bilan donnée.
  \item $\square$ Je sais utiliser l'ordre de pouvoir réducteur $\ce{Al} > \ce{Zn} > \ce{Fe} > \ce{H2} > \ce{Cu} > \ce{Ag} > \ce{Au}$ pour prévoir le sens d'une réaction.
  \item $\square$ Je sais caractériser le dihydrogène : \textbf{détonation à la flamme} (test de la flamme).
  \item $\square$ Je vérifie \textbf{systématiquement} la conservation des éléments et des charges dans toutes mes équations.
  \item $\square$ Je n'écris \textbf{jamais} d'électrons dans l'équation-bilan finale.
  \item $\square$ Je fais attention au cas de l'aluminium (passivation) et à la nature de l'acide (dilué vs concentré).
\end{itemize}
\end{aretenir}

\findecours

\end{document}
